% Leçon 23 — Vocabulaire et compréhension de texte : exode rural — Chinois 1ère A — Cours signé OnBuch+
% Programme officiel MINESEC. Compiler avec : tectonic ZHA23-vocabulaire-exode-rural.tex
\newcommand{\DOCMATIERE}{Chinois}
\newcommand{\DOCNIVEAU}{1ère A}
\newcommand{\DOCMODULE}{Module 3 : Citoyenneté et ouverture au monde}
\newcommand{\DOCLECON}{ZHA23}
\newcommand{\DOCTITRE}{Leçon 23 — Vocabulaire et compréhension de texte : exode rural}
% OnBuch+ — préambule « cours signés » ENRICHI (compilé avec tectonic / XeLaTeX)
% Système visuel « Pop » d'OnBuch+ : fond crème (jamais blanc), bordures encre
% épaisses, ombres « dures » décalées, titres Archivo Black en capitales.
% Version MATHÉMATIQUES : graphes (pgfplots/tikz), boîtes pédagogiques
% (définition, propriété, méthode, attention, le savais-tu, exercice résolu,
% exercice, TP, bilan), page de garde et signature OnBuch+ ancrée en bas.
%
% Macros attendues dans main.tex AVANT \input{preamble} :
%   \DOCMATIERE  \DOCNIVEAU  \DOCMODULE  \DOCLECON (numéro)  \DOCTITRE
\documentclass[11pt]{article}

\usepackage{fontspec}
\usepackage[french]{babel} % français = langue principale ; le chinois via \zh{..} / textebox (xeCJK)
\usepackage{xeCJK}
\usepackage{pifont} % \ding{51} \ding{55} utilisés par les modèles
\usepackage[a4paper,margin=2cm,top=3.4cm,bottom=2.8cm,headheight=1.4cm,footskip=1.3cm]{geometry}
\usepackage{amsmath,amssymb,amsfonts}
\usepackage{mathtools} % \xmapsto, \coloneqq... souvent utilisés par les modèles
\usepackage{cancel} % simplifications barrées (bilans de forces)
% \abs{z} (module, valeur absolue) et \norm{u} (norme), utilisés par les modèles
\providecommand{\abs}{}\let\abs\relax\DeclarePairedDelimiter{\abs}{\lvert}{\rvert}
\providecommand{\norm}{}\let\norm\relax\DeclarePairedDelimiter{\norm}{\lVert}{\rVert}
\usepackage[table]{xcolor} % [table] : fournit aussi \rowcolors (alternance de lignes)
\usepackage{pagecolor}
\usepackage{tikz}
\usetikzlibrary{arrows.meta,positioning,calc,shapes.geometric,shapes.misc,shapes.arrows,decorations.pathmorphing,decorations.pathreplacing,decorations.markings,patterns,shadows,mindmap,trees,fit,backgrounds,matrix,intersections,angles,quotes}
\usepackage{pgfplots}
\usepackage[version=4]{mhchem} % équations nucléaires \ce{^{235}_{92}U}
\usepackage[european,siunitx]{circuitikz} % schémas électriques
\pgfplotsset{compat=1.18}
\usepgfplotslibrary{fillbetween,groupplots} % aires entre courbes (calcul intégral)
\usepackage{enumitem}
\usepackage{fancyhdr}
% [most] charge toutes les bibliothèques tcolorbox (listings, minted, raster,
% documentation...) et gonfle inutilement la mémoire TeX sur un document long
% (~300 boîtes) : on ne charge que ce qui est réellement utilisé.
\usepackage[skins,breakable]{tcolorbox}
\usepackage{graphicx}
\usepackage{colortbl}
\usepackage{booktabs}
\usepackage{tabularx}
\usepackage{diagbox} % case d'en-tête diagonale des tableaux à double entrée
% Colonnes extensibles centrée / gauche / droite (C, L, R), souvent utilisées par les modèles
\newcolumntype{C}{>{\centering\arraybackslash}X}
\newcolumntype{L}{>{\raggedright\arraybackslash}X}
\newcolumntype{R}{>{\raggedleft\arraybackslash}X}
\usepackage{array}
\usepackage{multicol}
\usepackage{multirow}
\usepackage{siunitx}
% parse-numbers=false : accepte aussi \SI{2,5 \times 10^{-3}}{...}
\sisetup{locale=FR,output-decimal-marker={,},group-minimum-digits=4,parse-numbers=false}
\DeclareSIUnit\ohm{\ensuremath{\symup{\Omega}}} % U+2126 absent de la police
\DeclareSIUnit\division{div} % réglages d'oscilloscope (ms/div, V/div)
\DeclareSIUnit\tr{tr} % tours (vitesse de rotation, ex. \SI{1500}{\tr\per\minute})
\DeclareSIUnit\molar{M} % concentration molaire (ex. \SI{3}{\molar}, \SI{1}{\micro\molar})
\DeclareSIUnit\an{an} % année (le modèle confond parfois avec \year, un compteur TeX) — ex. \SI{5}{\kilo\meter\per\an}
\DeclareSIUnit\FCFA{FCFA} % franc CFA (ex. \SI{500}{\FCFA\per\kilo\gram})
\DeclareSIUnit\degreeBrix{\textdegree Brix} % teneur en sucre d'un sirop/jus (réfractométrie)
\DeclareSIUnit\month{mois} % le modèle utilise parfois \month, un compteur TeX, comme unité de durée
\usepackage{qrcode}
% \degree : commande fréquemment utilisée par les modèles pour un angle ou une
% température en dehors de \SI{}{} (non fournie par les paquets chargés).
\providecommand{\degree}{^{\circ}}
% \qed (fin de démonstration), utilisé par les modèles sans amsthm
\providecommand{\qed}{\hfill$\square$}
% \barre : double barre des valeurs interdites dans un tableau de variations (tkz-tab)
\providecommand{\barre}{\ensuremath{\|}}
% environnement proof (amsthm non chargé) utilisé par les modèles
\ifdefined\proof\else\newenvironment{proof}[1][Démonstration]{\par\noindent\textit{#1.}\ }{\qed\par}\fi
\usepackage{lastpage}
\usepackage{hyperref}
\hypersetup{hidelinks}

% ============================================================
% Palette « Pop » OnBuch+ — mêmes hex que la classe P (plus_ui.dart)
% ============================================================
\definecolor{poporange}{HTML}{F59321}
\definecolor{popdark}{HTML}{E07A0C}
\definecolor{popcream}{HTML}{FFF6E8}
\definecolor{popink}{HTML}{1C1714}
\definecolor{poppurple}{HTML}{7A5AE0}
\definecolor{popgreen}{HTML}{2E9E6A}
\definecolor{poppink}{HTML}{E0457B}
\definecolor{popblue}{HTML}{2F7DE1}
\definecolor{popgold}{HTML}{C98A12}
\definecolor{popmuted}{HTML}{8A7F76}
\definecolor{popline}{HTML}{EADFD2}
\definecolor{popgray}{HTML}{8A7F76}
\colorlet{popgrey}{popgray}
% Alias tolérants (le modèle nomme parfois les couleurs différemment)
\colorlet{popwhite}{white}
\colorlet{popblack}{popink}
\colorlet{popred}{poppink}
\colorlet{popyellow}{popgold}
% Teintes claires (fonds de tableaux/encadrés) — le modèle nomme parfois
% les couleurs avec un suffixe L (ex. popblueL) sans les avoir définies.
\colorlet{poporangeL}{poporange!20}
\colorlet{popdarkL}{popdark!20}
\colorlet{poppurpleL}{poppurple!20}
\colorlet{popgreenL}{popgreen!20}
\colorlet{poppinkL}{poppink!20}
\colorlet{popblueL}{popblue!20}
\colorlet{popgoldL}{popgold!20}
\colorlet{popredL}{popred!20}
\colorlet{popgrayL}{popgray!20}
% Teintes claires pour fonds de tableaux / zones
\colorlet{popblueL}{popblue!10!white}
\colorlet{poporangeL}{poporange!14!white}
\colorlet{popgreenL}{popgreen!12!white}
\colorlet{poppurpleL}{poppurple!12!white}
\colorlet{poppinkL}{poppink!12!white}
\colorlet{popgoldL}{popgold!14!white}

\pagecolor{popcream}

% ============================================================
% Typographie
% ============================================================
\defaultfontfeatures{Ligatures=TeX}
\setmainfont{PlusJakartaSans-Medium.ttf}[Path=../../../fonts/,BoldFont=PlusJakartaSans-Bold.ttf]
% Symboles, API et emojis absents de la police principale : repli sur DejaVu Sans (caractères actifs)
\newfontface\symfont{DejaVuSans.ttf}[Path=../../../fonts/]
\makeatletter
\newcommand\symactive[1]{\catcode#1=13 \begingroup\lccode`\~=#1 \lowercase{\endgroup\def~}{{\symfont\char#1}}}
\newcommand\symblank[1]{\catcode#1=13 \begingroup\lccode`\~=#1 \lowercase{\endgroup\def~}{}}
\newcommand\symhyph[1]{\catcode#1=13 \begingroup\lccode`\~=#1 \lowercase{\endgroup\def~}{\mbox{-}}}
\newcount\symc
\newcommand\symrange[2]{\symc=#1 \@whilenum\symc<#2 \do{\expandafter\symactive\expandafter{\the\symc}\advance\symc by 1 }}
\newcommand\symblankrange[2]{\symc=#1 \@whilenum\symc<#2 \do{\expandafter\symblank\expandafter{\the\symc}\advance\symc by 1 }}
\makeatother
\symrange{"0250}{"02FF}\symactive{"2713}\symactive{"2714}\symactive{"2717}\symactive{"2718}\symactive{"2610}\symactive{"2611}\symactive{"2612}
\symrange{"2600}{"27BF}
\catcode"2705=13 \begingroup\lccode`\~="2705 \lowercase{\endgroup\def~}{{\symfont\char"2713}}\symblank{"26BD}\symblank{"26A1}
\symrange{"2460}{"257F}\symactive{"223C}\symactive{"2248}\symactive{"2260}
\symactive{"01F8}\symactive{"01F9}\symactive{"1E3E}\symactive{"1E3F}
\symhyph{"2011}\symhyph{"2E1A}\symblankrange{"1F300}{"1FAFF}\symblank{"FE0F}
% Caractères chinois : WenQuanYi Zen Hei (sous-ensemble GB2312 + pinyin), gras/italique simulés
\setCJKmainfont{WenQuanYiZenHei-subset.ttf}[Path=../../../fonts/,AutoFakeBold=3,AutoFakeSlant=0.2]
\xeCJKsetup{CJKspace=false}
% Pinyin : voyelles à caron (ǎ ǐ ǒ ǔ ǖ ǘ ǚ ǜ) absentes de la police principale -> caractères actifs
\newfontface\pyfont{WenQuanYiZenHei-subset.ttf}[Path=../../../fonts/]
\newcommand\pyactive[1]{\catcode#1=13 \begingroup\lccode`\~=#1 \lowercase{\endgroup\def~}{{\pyfont\char#1}}}
\pyactive{"01CD}\pyactive{"01CE}\pyactive{"01CF}\pyactive{"01D0}\pyactive{"01D1}\pyactive{"01D2}\pyactive{"01D3}\pyactive{"01D4}\pyactive{"01D5}\pyactive{"01D6}\pyactive{"01D7}\pyactive{"01D8}\pyactive{"01D9}\pyactive{"01DA}\pyactive{"01DB}\pyactive{"01DC}
% Maths en sans-serif (Fira Math) avec chiffres et lettres droites de Plus
% Jakarta, pour que les formules se fondent dans le texte.
\usepackage[math-style=ISO,bold-style=ISO]{unicode-math}
\setmathfont{FiraMath-Regular.otf}[Path=../../../fonts/]
\setmathfont{PlusJakartaSans-Medium.ttf}[Path=../../../fonts/,range={up/num,up/{Latin,latin},\mathup}]
\usepackage{icomma} % virgule décimale sans espace en mode math
% Caractères mathématiques Unicode absents des polices de texte (Plus Jakarta,
% Archivo Black) : rendus par leur équivalent mathématique, en texte comme en
% formule — sinon ils s'affichent en carré vide (ex. « Arithmétique dans ℤ »).
\usepackage{newunicodechar}
\newunicodechar{ℝ}{\ensuremath{\mathbb{R}}}
\newunicodechar{ℤ}{\ensuremath{\mathbb{Z}}}
\newunicodechar{ℂ}{\ensuremath{\mathbb{C}}}
\newunicodechar{ℕ}{\ensuremath{\mathbb{N}}}
\newunicodechar{ℚ}{\ensuremath{\mathbb{Q}}}
\newunicodechar{𝔻}{\ensuremath{\mathbb{D}}}
\newunicodechar{α}{\ensuremath{\alpha}}
\newunicodechar{β}{\ensuremath{\beta}}
\newunicodechar{γ}{\ensuremath{\gamma}}
\newunicodechar{δ}{\ensuremath{\delta}}
\newunicodechar{ε}{\ensuremath{\varepsilon}}
\newunicodechar{θ}{\ensuremath{\theta}}
\newunicodechar{λ}{\ensuremath{\lambda}}
\newunicodechar{μ}{\ensuremath{\mu}}
\newunicodechar{σ}{\ensuremath{\sigma}}
\newunicodechar{τ}{\ensuremath{\tau}}
\newunicodechar{φ}{\ensuremath{\varphi}}
\newunicodechar{ψ}{\ensuremath{\psi}}
\newunicodechar{ω}{\ensuremath{\omega}}
\newunicodechar{Σ}{\ensuremath{\Sigma}}
% majuscules produites par \MakeUppercase dans les titres (α-aminés -> Α)
\newunicodechar{Α}{\ensuremath{\alpha}}
\newunicodechar{Β}{\ensuremath{\beta}}
\newunicodechar{Γ}{\ensuremath{\Gamma}}
\newunicodechar{Δ}{\ensuremath{\Delta}}
\newunicodechar{Ε}{\ensuremath{\varepsilon}}
\newunicodechar{Θ}{\ensuremath{\Theta}}
\newunicodechar{Λ}{\ensuremath{\Lambda}}
\newunicodechar{Μ}{\ensuremath{\mu}}
\newunicodechar{Τ}{\ensuremath{\tau}}
\newunicodechar{Φ}{\ensuremath{\Phi}}
\newunicodechar{Ψ}{\ensuremath{\Psi}}
\newunicodechar{Ω}{\ensuremath{\Omega}}
\newunicodechar{↦}{\ensuremath{\mapsto}}
\newunicodechar{∈}{\ensuremath{\in}}
\newunicodechar{∉}{\ensuremath{\notin}}
\newunicodechar{∧}{\ensuremath{\wedge}}
\newunicodechar{⇔}{\ensuremath{\Leftrightarrow}}
\newunicodechar{⟺}{\ensuremath{\Longleftrightarrow}}
\newunicodechar{⇒}{\ensuremath{\Rightarrow}}
\newunicodechar{⟹}{\ensuremath{\Longrightarrow}}
\newunicodechar{∘}{\ensuremath{\circ}}
\newunicodechar{∪}{\ensuremath{\cup}}
\newunicodechar{∩}{\ensuremath{\cap}}
\newunicodechar{≡}{\ensuremath{\equiv}}
\newunicodechar{⊂}{\ensuremath{\subset}}
\newunicodechar{⊥}{\ensuremath{\perp}}
\newunicodechar{∃}{\ensuremath{\exists}}
\newunicodechar{∀}{\ensuremath{\forall}}
\newunicodechar{∛}{\ensuremath{\sqrt[3]{\,}}}
\newunicodechar{∜}{\ensuremath{\sqrt[4]{\,}}}
\newunicodechar{⃗}{\ensuremath{{}^{\to}}}
\newunicodechar{ⁿ}{\ifmmode^{n}\else\textsuperscript{\ensuremath{n}}\fi}
\newunicodechar{ˣ}{\ifmmode^{x}\else\textsuperscript{\ensuremath{x}}\fi}
\newunicodechar{ᵇ}{\ifmmode^{b}\else\textsuperscript{\ensuremath{b}}\fi}
\newunicodechar{ʳ}{\ifmmode^{r}\else\textsuperscript{\ensuremath{r}}\fi}
\newunicodechar{ᵘ}{\ifmmode^{u}\else\textsuperscript{\ensuremath{u}}\fi}
\newunicodechar{ʸ}{\ifmmode^{y}\else\textsuperscript{\ensuremath{y}}\fi}
\newunicodechar{ᵃ}{\ifmmode^{a}\else\textsuperscript{\ensuremath{a}}\fi}
\newunicodechar{ᵉ}{\ifmmode^{e}\else\textsuperscript{\ensuremath{e}}\fi}
\newunicodechar{ᵏ}{\ifmmode^{k}\else\textsuperscript{\ensuremath{k}}\fi}
\newunicodechar{ᵖ}{\ifmmode^{p}\else\textsuperscript{\ensuremath{p}}\fi}
\newunicodechar{ᴺ}{\ifmmode^{N}\else\textsuperscript{\ensuremath{N}}\fi}
\newunicodechar{ᶜ}{\ifmmode^{c}\else\textsuperscript{\ensuremath{c}}\fi}
\newunicodechar{ʰ}{\ifmmode^{h}\else\textsuperscript{\ensuremath{h}}\fi}
\newunicodechar{ᵠ}{\ifmmode^{q}\else\textsuperscript{\ensuremath{q}}\fi}
\newunicodechar{ₙ}{\ifmmode_{n}\else\textsubscript{\ensuremath{n}}\fi}
\newunicodechar{ₐ}{\ifmmode_{a}\else\textsubscript{\ensuremath{a}}\fi}
\newunicodechar{ᵢ}{\ifmmode_{i}\else\textsubscript{\ensuremath{i}}\fi}
\newunicodechar{ᵣ}{\ifmmode_{r}\else\textsubscript{\ensuremath{r}}\fi}
\newunicodechar{ₖ}{\ifmmode_{k}\else\textsubscript{\ensuremath{k}}\fi}
\newunicodechar{ᵦ}{\ifmmode_{\beta}\else\textsubscript{\ensuremath{\beta}}\fi}
\newunicodechar{ₓ}{\ifmmode_{x}\else\textsubscript{\ensuremath{x}}\fi}
\newfontfamily\popdisplay{ArchivoBlack-Regular.ttf}[Path=../../../fonts/]
\newfontfamily\poplogo{SpaceGrotesk-Bold.ttf}[Path=../../../fonts/]
\newfontfamily\popsemi{PlusJakartaSans-SemiBold.ttf}[Path=../../../fonts/]
\newfontfamily\popxbold{PlusJakartaSans-ExtraBold.ttf}[Path=../../../fonts/]
\newfontfamily\popxboldls{PlusJakartaSans-ExtraBold.ttf}[Path=../../../fonts/,LetterSpace=3.0]

\setlist[enumerate]{leftmargin=1.7em,itemsep=5pt,topsep=4pt}
\setlist[itemize]{leftmargin=1.7em,itemsep=4pt,topsep=4pt,label={\color{poporange}\textbullet}}
\setlength{\parindent}{0pt}
\setlength{\parskip}{5pt}
\renewcommand{\arraystretch}{1.25}

% pgfplots : style Pop par défaut
\pgfplotsset{
  every axis/.append style={axis line style={popink,line width=1pt},tick style={popink},
    grid style={popline,line width=0.5pt},label style={font=\small},tick label style={font=\footnotesize}},
  popcurve/.style={poporange,line width=1.8pt,no markers,smooth},
}
% Même style disponible dans un \draw TikZ simple (hors pgfplots)
\tikzset{popcurve/.style={poporange,line width=1.8pt}}
\tikzset{popgrid/.style={popline,very thin}}
\colorlet{popcurve}{poporange} % le nom du style est parfois utilisé comme couleur

% ============================================================
% Pill / squiggle
% ============================================================
\newcommand{\poppill}[3]{%
  \tikz[baseline=-0.55ex]\node[draw=popink,line width=1.2pt,rounded corners=2.6pt,%
    fill=#2,inner xsep=6.5pt,inner ysep=3pt]{%
    \popxboldls\fontsize{7}{7}\selectfont\color{#3}\MakeUppercase{#1}};%
}
\newcommand{\popsquiggle}[1]{%
  \tikz[baseline=-0.4ex]{
    \draw[#1,line width=2.6pt,line cap=round]
      plot [smooth] coordinates {(0,0.09) (0.34,0.01) (0.68,0.21) (1.02,0.01) (1.36,0.10)};
  }%
}
% Mot-clé surligné (marqueur orange sous le texte)
\newcommand{\cle}[1]{{\popxbold\color{popdark}#1}}

% ============================================================
% En-tête / pied
% ============================================================
\pagestyle{fancy}
\fancyhf{}
\renewcommand{\headrulewidth}{0pt}
\renewcommand{\footrulewidth}{0pt}
\lhead{\raisebox{-0.15ex}{\poplogo\fontsize{13}{13}\selectfont\color{popink}OnBuch\color{poporange}+}}
\rhead{\poppill{\DOCMATIERE\ \textbullet\ \DOCNIVEAU\ \textbullet\ Leçon \DOCLECON}{white}{popink}}
\lfoot{\popxbold\fontsize{7.6}{7.6}\selectfont\color{popmuted}\MakeUppercase{Cours signé OnBuch+}\ \textbullet\ {\popsemi\DOCTITRE}}
\rfoot{\tikz[baseline=-0.6ex]\node[rounded corners=8.5pt,fill=popink,minimum height=17pt,minimum width=17pt,inner xsep=5pt,inner sep=0]{\popxbold\fontsize{8}{8}\selectfont\color{popcream}\thepage\,/\,\pageref*{LastPage}};}

% ============================================================
% Titres
% ============================================================
\newcommand{\chaptitre}[2]{%
  \begin{tcolorbox}[enhanced,colback=poporange,colframe=popink,boxrule=2.6pt,arc=11pt,
      drop shadow=popink,left=15pt,right=15pt,top=12pt,bottom=11pt,before skip=3pt,after skip=18pt]
    \poppill{#2}{popink}{popcream}\par\vspace{9pt}
    {\raggedright\hyphenpenalty=10000\exhyphenpenalty=10000\popdisplay\fontsize{22}{24}\selectfont\color{popcream}\MakeUppercase{#1}\par}
  \end{tcolorbox}
}
% Section numérotée : pastille orange avec le numéro + titre Archivo
\newcounter{popsec}
\newcommand{\coursec}[1]{%
  \stepcounter{popsec}\setcounter{popsub}{0}%
  \par\vspace{16pt}\noindent
  \begin{minipage}{\linewidth}
  \tikz[baseline=-0.7ex]\node[draw=popink,line width=1.6pt,fill=poporange,rounded corners=4pt,minimum size=22pt,inner sep=0]{\popdisplay\fontsize{12}{12}\selectfont\color{popcream}\thepopsec};%
  \hspace{8pt}{\popdisplay\fontsize{15}{17}\selectfont\color{popink}\MakeUppercase{#1}}\par
  \vspace{4pt}\noindent{\color{poporange}\rule{36pt}{3.6pt}}
  \end{minipage}\par\nopagebreak\vspace{8pt}
}
\newcounter{popsub}
\newcommand{\courssub}[1]{%
  \stepcounter{popsub}%
  \par\vspace{9pt}\noindent{\popxbold\fontsize{11.5}{13}\selectfont\color{popink}{\color{poporange}\thepopsec.\thepopsub}\hspace{6pt}#1}\par\nopagebreak\vspace{4pt}
}

% ============================================================
% Boîtes pédagogiques — même recette : fond blanc, bordure épaisse colorée,
% ombre dure encre, badge plein. Titre optionnel [#1].
% ============================================================
\tcbset{popbase/.style={breakable,enhanced,colback=white,boxrule=2.3pt,arc=9pt,
  shadow={3pt}{-3pt}{0pt}{popink},left=14pt,right=14pt,top=11pt,bottom=11pt,before skip=11pt,after skip=15pt,
  fontupper=\popsemi\fontsize{10.6}{14.5}\selectfont\color{popink},
  fontlower=\popsemi\fontsize{10.6}{14.5}\selectfont\color{popink}}}
% Coche dessinée (le glyphe ✓ n'existe pas dans les polices du document)
\newcommand{\popcheck}[1]{\tikz[baseline=-0.5ex]\draw[#1,line width=1.6pt,line cap=round,line join=round](-0.13,0.0)--(-0.04,-0.09)--(0.13,0.1);}
\providecommand{\coche}{\popcheck{popgreen}}
\providecommand{\resultat}[1]{\par\smallskip\noindent\fcolorbox{popgreen}{popgreenL}{\parbox{\dimexpr\linewidth-2\fboxsep-2\fboxrule\relax}{\textbf{Résultat.} #1}}\par\smallskip}
\newcommand{\popboxhead}[3]{\poppill{#1}{#2}{white}\ifx\relax#3\relax\else\hspace{6pt}{\popxbold\fontsize{10.6}{12}\selectfont\color{popink}#3}\fi\par\vspace{7pt}}

\newtcolorbox{aretenir}[1][]{popbase,colframe=popgreen,before upper={\popboxhead{À retenir}{popgreen}{#1}}}
% Texte en chinois (document de lecture, dialogue, extrait) : cadre
% d'étude. Un titre optionnel (source, auteur) s'affiche dans l'en-tête.
\newtcolorbox{textebox}[1][]{popbase,colframe=popblue,before upper={\popboxhead{文本 · Texte}{popblue}{#1}},after upper={}}
% Marques ✓ / ✗ (forme correcte / fautive) souvent utilisées par les modèles
\providecommand{\cross}{\textcolor{poppink}{\ensuremath{\times}}}
\providecommand{\xmark}{\cross}\providecommand{\crossmark}{\cross}\providecommand{\cmark}{\ensuremath{\checkmark}}
% Ligne de réponse / justification à compléter (inventée par les modèles)
\providecommand{\justification}[1]{\par\noindent\textit{Justification~:} \dotfill\par}
% \textipa (package tipa non chargé) : la police ne couvre pas l'API -> simple passage
\providecommand{\textipa}[1]{#1}
% Soulignement (ulem non chargé) : \uline, \ul
\providecommand{\uline}[1]{\underline{#1}}\providecommand{\ul}[1]{\underline{#1}}
% \glossaire{\item ...} inventée par les modèles : liste de glossaire
\providecommand{\glossaire}[1]{\begin{itemize}#1\end{itemize}}
% \alert (beamer) inventée par les modèles : texte mis en évidence
\providecommand{\alert}[1]{\textbf{\textcolor{poppink}{#1}}}
% variantes ASCII d'ordinaux écrites en macro
\providecommand{\iere}{\textsuperscript{re}}\providecommand{\ieme}{\textsuperscript{e}}\providecommand{\ere}{\textsuperscript{re}}\providecommand{\eme}{\textsuperscript{e}}\providecommand{\er}{\textsuperscript{er}}\providecommand{\ie}{\textsuperscript{e}}
% Ordinaux français écrits en macro par les modèles : 1\ière, 3\ième
\newcommand{\ière}{\textsuperscript{re}}\newcommand{\ième}{\textsuperscript{e}}
% \exemple (inventée par les modèles) : étiquette « Exemple : »
\providecommand{\exemple}{\textbf{Exemple~:}\ }
% \square utilisable aussi hors mode math (cases à cocher)
\AtBeginDocument{\let\squareMath\square\renewcommand{\square}{\ensuremath{\squareMath}}}
% Mot ou phrase en chinois à l'intérieur d'un paragraphe français
\newcommand{\zh}[1]{\ifmmode\text{#1}\else{#1}\fi}
\let\eng\zh \let\esp\zh \let\all\zh % alias tolérés
% flèches d'intonation inventées par les modèles
\providecommand{\textsearrow}{}\renewcommand{\textsearrow}{\ensuremath{\searrow}}\providecommand{\textnearrow}{}\renewcommand{\textnearrow}{\ensuremath{\nearrow}}
% \fr{..} : mot français (inventé par les modèles)
\providecommand{\fr}[1]{\textit{#1}}
% zhbox : encadré de texte chinois inventé par les modèles
\newenvironment{zhbox}{\begin{quote}}{\end{quote}}
% \courssubsub : titre de niveau 3 inventé par les modèles
\providecommand{\courssubsub}[1]{\par\medskip\noindent\textbf{#1}\par\smallskip}
% \vs : « versus » inventé par les modèles
\providecommand{\vs}{\textit{vs}\ }
% \blank : case à compléter inventée par les modèles
\providecommand{\blank}[1][]{\underline{\hspace{1.6em}}}
\newtcolorbox{exemplebox}[1][]{popbase,colframe=poporange,before upper={\popboxhead{Exemple}{poporange}{#1}}}
\newtcolorbox{definition}[1][]{popbase,colframe=popblue,before upper={\popboxhead{Définition}{popblue}{#1}}}
\newtcolorbox{propriete}[1][]{popbase,colframe=poppurple,before upper={\popboxhead{Propriété}{poppurple}{#1}}}
\newtcolorbox{methode}[1][]{popbase,colframe=popgold,before upper={\popboxhead{Méthode}{popgold}{#1}}}
\newtcolorbox{attention}[1][]{popbase,colframe=poppink,before upper={\popboxhead{Attention}{poppink}{#1}}}
\newtcolorbox{experience}[1][]{popbase,colframe=popblue,colback=popblueL,before upper={\popboxhead{Expérience}{popblue}{#1}}}
% « Le savais-tu ? » : carte violette pleine (contexte, culture, Cameroun)
\newtcolorbox{savaistu}[1][]{popbase,colback=poppurple,colframe=popink,
  fontupper=\popsemi\fontsize{10.4}{14.2}\selectfont\color{white},
  before upper={\poppill{Le savais-tu\,?}{white}{poppurple}\ifx\relax#1\relax\else\hspace{6pt}{\popxbold\color{white}#1}\fi\par\vspace{7pt}}}
% Exercice résolu : énoncé en haut, \tcblower puis la solution sur fond crème
\newcounter{popexo}\newcounter{popexr}
\newtcolorbox{exoresolu}[1][]{popbase,colframe=poporange,colbacklower=popcream,
  segmentation style={popink,line width=1.2pt,dashed},
  before upper={\stepcounter{popexr}\popboxhead{Exercice résolu \thepopexr}{poporange}{#1}},
  before lower={\poppill{Solution}{popink}{popcream}\par\vspace{6pt}}}
% Exercice (non résolu en place) — la correction est dans la partie Corrigés
\newtcolorbox{exercice}[1][]{popbase,colframe=popink,
  before upper={\stepcounter{popexo}\popboxhead{Exercice \thepopexo}{popink}{#1}}}
% Corrigé
\newtcolorbox{corrige}[1][]{popbase,colframe=popgreen,colback=popgreenL,
  before upper={\popboxhead{Corrigé}{popgreen}{#1}}}
% Objectifs / prérequis / bilan
\newtcolorbox{objectifs}{popbase,colframe=popink,colback=poporangeL,
  before upper={\popboxhead{Objectifs de la leçon}{popink}{}}}
\newtcolorbox{prerequis}{popbase,colframe=popink,colback=popblueL,
  before upper={\popboxhead{Prérequis}{popblue}{}}}
\newtcolorbox{bilan}{popbase,colframe=popink,colback=white,boxrule=2.8pt,
  before upper={\popboxhead{Fiche bilan}{poporange}{L'essentiel à connaître pour l'examen}}}
% Figure encadrée avec légende
\newtcolorbox{popfigure}[1][]{enhanced,breakable=false,colback=white,colframe=popline,boxrule=1.4pt,arc=8pt,
  left=10pt,right=10pt,top=10pt,bottom=8pt,before skip=10pt,after skip=12pt,halign=center,
  lower separated=false,halign lower=center,
  fontlower=\popsemi\fontsize{9}{11.5}\selectfont\color{popmuted},
  #1}
\newcommand{\legende}[1]{\par\vspace{6pt}{\popsemi\fontsize{9}{11.5}\selectfont\color{popmuted}#1}}

% ============================================================
% Signature OnBuch+
% ============================================================
\newcommand{\poplogoinline}[1]{{\poplogo\fontsize{#1}{#1}\selectfont\color{popink}OnBuch\color{poporange}+}}
\newcommand{\signatureonbuch}{%
  \begin{tcolorbox}[enhanced,colback=popink,colframe=popink,boxrule=2.2pt,arc=10pt,
      drop shadow=poporange,left=14pt,right=14pt,top=10pt,bottom=10pt,before skip=0pt,after skip=0pt]
    \tikz[baseline=-0.7ex]\node[circle,fill=popgreen,draw=popcream,line width=1.3pt,minimum size=22pt,inner sep=0]{\popcheck{white}};%
    \hspace{9pt}\begin{minipage}[c]{0.62\linewidth}
      {\popxbold\fontsize{12}{13}\selectfont\color{popcream}Signé\ \textbullet\ {\poplogo OnBuch\color{poporange}+}}\par\vspace{2pt}
      {\raggedright\popsemi\fontsize{8}{10.5}\selectfont\color{popline}\MakeUppercase{Équipe pédagogique OnBuch+}\\\MakeUppercase{Conforme au programme officiel MINESEC}\par}
    \end{minipage}\hfill
    {\poplogo\fontsize{22}{22}\selectfont\color{popcream}OnBuch\color{poporange}+}
  \end{tcolorbox}%
}

% ============================================================
% Page de garde
% #1 titre · #2 matière · #3 niveau · #4 module · #5 n° leçon · #6 durée ·
% #7 description / objectifs (texte ou itemize)
% ============================================================
\newcommand{\pagegarde}[7]{%
  \thispagestyle{empty}
  \newgeometry{a4paper,margin=1.7cm,top=1.6cm,bottom=1.4cm}%
  \begin{tikzpicture}[remember picture,overlay]
    \fill[poporange!18] ([xshift=-3.2cm,yshift=-3.6cm]current page.north east) circle (4.6cm);
    \fill[poppurple!14] ([xshift=2.4cm,yshift=7.6cm]current page.south west) circle (3.8cm);
  \end{tikzpicture}%
  {\poplogo\fontsize{30}{30}\selectfont\color{popink}OnBuch\color{poporange}+}\hfill
  \raisebox{0.6ex}{\poppill{Cours signé \textbullet\ Premium}{popink}{popcream}}\par
  \vspace{1pt}\popsquiggle{poppurple}\par
  \vspace{8pt}
  \begin{tcolorbox}[enhanced,colback=poporange,colframe=popink,boxrule=2.8pt,arc=13pt,
      drop shadow=popink,left=18pt,right=18pt,top=12pt,bottom=13pt,after skip=10pt]
    \poppill{#2 \textbullet\ #3}{popink}{popcream}\hspace{5pt}\poppill{#4}{white}{popink}\par\vspace{8pt}
    {\popdisplay\fontsize{14}{14}\selectfont\color{popink}LEÇON #5}\par\vspace{4pt}
    {\raggedright\hyphenpenalty=10000\exhyphenpenalty=10000\popdisplay\fontsize{25}{27}\selectfont\color{popcream}\MakeUppercase{#1}\par}
    \vspace{8pt}
    {\popxbold\fontsize{9.5}{9.5}\selectfont\color{popink}\MakeUppercase{Durée conseillée : #6}}
  \end{tcolorbox}
  \begin{tcolorbox}[enhanced,colback=white,colframe=popink,boxrule=2.2pt,arc=10pt,
      drop shadow=popink,left=16pt,right=16pt,top=10pt,bottom=10pt,after skip=8pt]
    \poppill{Dans cette leçon}{poppurple}{white}\par\vspace{5pt}
    \popsemi\fontsize{9.3}{12.6}\selectfont\color{popink}#7
  \end{tcolorbox}
  \noindent\begin{minipage}[t]{0.487\linewidth}
    \begin{tcolorbox}[enhanced,colback=popgreen,colframe=popink,boxrule=2.2pt,arc=10pt,
        drop shadow=popink,left=11pt,right=11pt,top=10pt,bottom=10pt,height=3.1cm,valign=center]
      \tcbox[colback=white,colframe=popink,boxrule=1.3pt,arc=5pt,boxsep=0pt,nobeforeafter,
        left=3pt,right=3pt,top=3pt,bottom=3pt,valign=center]{\qrcode[height=1.9cm]{https://play.google.com/store/apps/details?id=com.luvvix.onbuch&hl=fr}}%
      \hspace{8pt}\begin{minipage}[c]{0.5\linewidth}
        {\popdisplay\fontsize{11}{12.5}\selectfont\color{white}\MakeUppercase{Télécharge OnBuch}}\par\vspace{4pt}
        {\raggedright\popsemi\fontsize{8.4}{11}\selectfont\color{white}Gratuit \textbullet\ Google Play\\Tuteur IA, annales, concours, résultats\par}
      \end{minipage}
    \end{tcolorbox}
  \end{minipage}\hfill
  \begin{minipage}[t]{0.487\linewidth}
    \begin{tcolorbox}[enhanced,colback=poppurple,colframe=popink,boxrule=2.2pt,arc=10pt,
        drop shadow=popink,left=11pt,right=11pt,top=10pt,bottom=10pt,height=3.1cm,valign=center]
      \tikz[baseline=-0.7ex]\node[circle,fill=white,draw=popink,line width=1.3pt,minimum size=1.6cm,inner sep=0]{\popdisplay\fontsize{24}{24}\selectfont\color{poppurple}+};%
      \hspace{8pt}\begin{minipage}[c]{0.6\linewidth}
        {\popdisplay\fontsize{11}{12.5}\selectfont\color{white}\MakeUppercase{Passe à OnBuch+}}\par\vspace{4pt}
        {\raggedright\popsemi\fontsize{8.4}{11}\selectfont\color{white}Tous les cours signés, Tuteur IA illimité, fiches PDF — onglet OnBuch+ de l'app\par}
      \end{minipage}
    \end{tcolorbox}
  \end{minipage}
  \vspace{8pt}
  \popcontenu
  \vfill
  \signatureonbuch
  \restoregeometry
  \newpage
}

% Rangée « Ce cours contient » (page de garde)
\newcommand{\poptile}[3]{% #1 couleur #2 gros texte #3 légende
  \begin{tcolorbox}[enhanced,colback=#1,colframe=popink,boxrule=2pt,arc=9pt,shadow={2.5pt}{-2.5pt}{0pt}{popink},
      left=6pt,right=6pt,top=9pt,bottom=9pt,halign=center,valign=center,height=1.75cm,width=\linewidth,nobeforeafter]
    {\popdisplay\fontsize{12.5}{13}\selectfont\color{white}\MakeUppercase{#2}}\par\vspace{3pt}
    {\centering\popsemi\fontsize{7.8}{9.5}\selectfont\color{white}#3\par}
  \end{tcolorbox}}
\newcommand{\popcontenu}{%
  \par\noindent{\popxboldls\fontsize{7.5}{7.5}\selectfont\color{popmuted}CE COURS SIGNÉ CONTIENT}\par\vspace{7pt}
  \noindent\begin{minipage}[t]{0.235\linewidth}\poptile{popblue}{Cours}{complet, illustré, pas à pas}\end{minipage}\hfill
  \begin{minipage}[t]{0.235\linewidth}\poptile{poporange}{Méthodes}{et exercices résolus}\end{minipage}\hfill
  \begin{minipage}[t]{0.235\linewidth}\poptile{poppink}{Exercices}{type examen + corrigés}\end{minipage}\hfill
  \begin{minipage}[t]{0.235\linewidth}\poptile{popgreen}{Fiche bilan}{l'essentiel à retenir}\end{minipage}\par}

% Sommaire
\newtcolorbox{sommaire}{enhanced,colback=white,colframe=popink,boxrule=2.2pt,arc=10pt,
  drop shadow=popink,left=18pt,right=18pt,top=13pt,bottom=13pt,before skip=6pt,after skip=20pt,
  before upper={\poppill{Sommaire}{poppurple}{white}\par\vspace{9pt}\popsemi\fontsize{10.6}{14.5}\selectfont\color{popink}}}

% Signature de fin de cours — poussée tout en bas de la dernière page
\newcommand{\findecours}{%
  \par\vfill\nopagebreak
  \begin{center}\popsquiggle{poporange}\end{center}\vspace{4pt}
  \signatureonbuch
}
\AtBeginDocument{\def\llbracket{\mathopen{[\mkern-3.5mu[}}\def\rrbracket{\mathclose{]\mkern-3.5mu]}}} % intervalles d'entiers (glyphe ⟦ absent de la police)
% Glyphes absents de Fira Math : redéfinis à partir de glyphes disponibles
\AtBeginDocument{%
  \def\setminus{\mathbin{\backslash}}\let\smallsetminus\setminus
  \def\vdots{\mathord{\vbox{\baselineskip3.2pt\lineskiplimit0pt\kern4pt\hbox{.}\hbox{.}\hbox{.}}}}%
  \def\ddots{\mathinner{\mkern1mu\raise6.5pt\hbox{.}\mkern2mu\raise3.6pt\hbox{.}\mkern2mu\raise0.7pt\hbox{.}\mkern1mu}}%
  \def\triangle{\mathord{\scalebox{0.9}{$\symup{\Delta}$}}}%
  \def\nabla{\mathord{\raisebox{\depth}{\scalebox{1}[-1]{$\symup{\Delta}$}}}}%
  \def\checkmark{\popcheck{popdark}}%
  \def\star{\mathbin{\ast}}\def\bigstar{\mathord{\ast}}%
  \def\diamond{\mathbin{\scalebox{0.6}{$\Diamond$}}}%
  \def\bigcup{\mathop{\mathchoice{\raisebox{-0.3em}{\scalebox{1.7}{$\cup$}}}{\scalebox{1.25}{$\cup$}}{\cup}{\cup}}\displaylimits}%
  \def\bigcap{\mathop{\mathchoice{\raisebox{-0.3em}{\scalebox{1.7}{$\cap$}}}{\scalebox{1.25}{$\cap$}}{\cap}{\cap}}\displaylimits}%
  \def\lesseqqgtr{\mathrel{\lessgtr}}\def\bigoplus{\mathop{\oplus}\displaylimits}%
}
\providecommand{\ssi}{si et seulement si} % abréviation fréquente
\AtBeginDocument{\providecommand{\wideparen}{\overparen}} % arc de cercle

%% FIGFIX-BEGIN v5 — correctifs globaux des figures (docs/onbuch-generation/tools/figfix)
\usepackage{adjustbox}\usepackage{etoolbox}\tcbuselibrary{raster}
% toute figure TikZ/pgfplots plus large que la ligne est réduite à la largeur disponible
\BeforeBeginEnvironment{tikzpicture}{\begin{adjustbox}{max width=\linewidth}}
\AfterEndEnvironment{tikzpicture}{\end{adjustbox}}
% légendes pgfplots lisibles par-dessus les courbes
\pgfplotsset{every axis/.append style={legend style={fill=white,fill opacity=0.92,text opacity=1,draw=black!15,inner sep=3pt}}}
% étiquettes d'arêtes (syntaxe edge["..."]) avec fond blanc
\tikzset{every edge quotes/.append style={fill=white,inner sep=1.2pt}}
%% FIGFIX-END
\begin{document}

\pagegarde{Leçon 23 — Vocabulaire et compréhension de texte : exode rural}{Chinois}{1ère A}{Module 3 : Citoyenneté et ouverture au monde}{ZHA23}{2 h}{Cette leçon de vocabulaire et compréhension de texte porte sur l'exode rural (农村人口外流), phénomène social majeur en Chine comme au Cameroun. À travers un texte support original, un glossaire structuré, l'étude des radicaux et des questions de compréhension, l'élève apprend à lire, comprendre et commenter un texte chinois sur ce thème.
\vspace{4pt}
\begin{multicols}{2}\small
\begin{itemize}
\item À la fin de cette leçon, je sais lire à voix haute un court texte chinois sur l'exode rural en respectant les tons
\item À la fin de cette leçon, je sais reconnaître et employer le vocabulaire clé du thème (campagne, ville, migration, travail, population)
\item À la fin de cette leçon, je sais identifier les radicaux et tracer correctement les caractères 村, 城, 流, 离, 工
\item À la fin de cette leçon, je sais répondre à des questions de compréhension en phrases complètes en chinois
\item À la fin de cette leçon, je sais utiliser les collocations du thème (离开家乡, 进城打工, 农村人口) dans des phrases
\item À la fin de cette leçon, je sais comparer brièvement l'exode rural en Chine et au Cameroun avec le vocabulaire appris
\end{itemize}\end{multicols}}

\begin{sommaire}
\begin{multicols}{2}
\begin{enumerate}
\item Texte support : 农村的年轻人
\item Glossaire du thème
\item Radicaux et ordre des traits
\item Compréhension du texte
\item Collocations et exemples
\item Production guidée
\item Activité d'intégration
\item Méthodes et exercices résolus
\item Exercices
\item Corrigés des exercices
\item Fiche bilan
\end{enumerate}
\end{multicols}
\end{sommaire}

\begin{objectifs}
\begin{itemize}
\item À la fin de cette leçon, je sais lire à voix haute un court texte chinois sur l'exode rural en respectant les tons
\item À la fin de cette leçon, je sais reconnaître et employer le vocabulaire clé du thème (campagne, ville, migration, travail, population)
\item À la fin de cette leçon, je sais identifier les radicaux et tracer correctement les caractères 村, 城, 流, 离, 工
\item À la fin de cette leçon, je sais répondre à des questions de compréhension en phrases complètes en chinois
\item À la fin de cette leçon, je sais utiliser les collocations du thème (离开家乡, 进城打工, 农村人口) dans des phrases
\item À la fin de cette leçon, je sais comparer brièvement l'exode rural en Chine et au Cameroun avec le vocabulaire appris
\item À la fin de cette leçon, je sais produire un court paragraphe guidé sur les causes de l'exode rural
\end{itemize}
\end{objectifs}

\begin{prerequis}
\begin{itemize}
\item Vocabulaire de base : 人, 家, 工作, 城市, 中国 (leçons de Seconde)
\item Structure de phrase simple : sujet + verbe + complément (SVO)
\item Lecture du pinyin et des quatre tons
\item Notion de radicaux (clés) vue en Seconde : 亻, 口, 木
\item Utilisation de 了 pour marquer un changement ou une action accomplie
\end{itemize}
\end{prerequis}

\newpage

\coursec{Texte support : 农村的年轻人}

\courssub{Situation d'accroche et problématique}

Chaque année, des milliers de jeunes quittent les villages du Nord, de l'Ouest ou de l'Extrême-Nord du Cameroun pour tenter leur chance à Douala ou à Yaoundé. Ils cherchent du travail, un salaire, une vie meilleure. Leurs villages, eux, se vident peu à peu : il ne reste souvent que les grands-parents et les enfants. Ce phénomène, les géographes l'appellent l'\cle{exode rural}.

Le même phénomène existe en Chine, mais à une échelle immense : des centaines de millions de travailleurs migrants ont quitté les campagnes pour les villes. C'est l'un des plus grands mouvements de population de l'histoire humaine.

\textbf{Problématique :} comment dit-on « exode rural » en chinois ? Quelles sont ses causes et ses conséquences ? Dans cette leçon, nous lisons un court texte chinois sur ce sujet d'actualité, et nous apprenons à en parler avec un vocabulaire précis. À la fin de la séance, vous saurez lire ce texte à voix haute, le comprendre sans dictionnaire et réutiliser ses mots dans vos propres phrases.

\courssub{Le texte : 农村的年轻人}

\begin{textebox}[农村的年轻人 — Les jeunes de la campagne]
现在，很多农村的年轻人离开家乡，到大城市去打工。\\
\emph{Xiànzài, hěn duō nóngcūn de niánqīngrén líkāi jiāxiāng, dào dà chéngshì qù dǎgōng.}\\
« Aujourd'hui, beaucoup de jeunes de la campagne quittent leur village natal pour aller travailler dans les grandes villes. »\\[4pt]
他们为什么离开农村呢？因为农村的工作机会比较少，工资也不高。\\
\emph{Tāmen wèishénme líkāi nóngcūn ne? Yīnwèi nóngcūn de gōngzuò jīhuì bǐjiào shǎo, gōngzī yě bù gāo.}\\
« Pourquoi quittent-ils la campagne ? Parce qu'à la campagne, les opportunités de travail sont assez rares et les salaires ne sont pas élevés non plus. »\\[4pt]
城市里有很多工厂和公司，年轻人可以找到更好的工作。\\
\emph{Chéngshì li yǒu hěn duō gōngchǎng hé gōngsī, niánqīngrén kěyǐ zhǎodào gèng hǎo de gōngzuò.}\\
« Dans les villes, il y a beaucoup d'usines et d'entreprises ; les jeunes peuvent trouver de meilleurs emplois. »\\[4pt]
但是，他们离开以后，农村里只剩下老人和孩子。\\
\emph{Dànshì, tāmen líkāi yǐhòu, nóngcūn li zhǐ shèngxià lǎorén hé háizi.}\\
« Mais après leur départ, il ne reste à la campagne que les personnes âgées et les enfants. »\\[4pt]
农村人口外流是一个很大的社会问题。\\
\emph{Nóngcūn rénkǒu wàiliú shì yí ge hěn dà de shèhuì wèntí.}\\
« L'exode rural est un très grand problème de société. »\\[4pt]
中国政府正在想办法帮助农村发展，让年轻人可以在家乡工作。\\
\emph{Zhōngguó zhèngfǔ zhèngzài xiǎng bànfǎ bāngzhù nóngcūn fāzhǎn, ràng niánqīngrén kěyǐ zài jiāxiāng gōngzuò.}\\
« Le gouvernement chinois est en train de chercher des solutions pour aider la campagne à se développer, afin que les jeunes puissent travailler dans leur village natal. »
\end{textebox}

\begin{methode}[Comment lire un texte chinois en quatre étapes]
\begin{enumerate}
\item \textbf{Lire le titre d'abord.} Le titre \zh{农村的年轻人} annonce le sujet : « les jeunes de la campagne ». Avant même de lire, on sait de qui on va parler.
\item \textbf{Repérer les mots déjà connus.} Soulignez les mots que vous connaissez déjà : \zh{人} (personne), \zh{工作} (travailler, travail), \zh{城市} (ville), \zh{很多} (beaucoup), \zh{为什么} (pourquoi). Ils forment le squelette du texte.
\item \textbf{Découper en mots, pas en caractères.} En chinois, un mot fait souvent deux caractères : \zh{年轻人} = \zh{年轻} (jeune) + \zh{人} (personne) = « jeune personne, jeune ». De même \zh{家乡} = \zh{家} (famille) + \zh{乡} (village, pays natal) = « village natal ».
\item \textbf{Deviner les mots nouveaux par le contexte et les radicaux.} Exemple : dans \zh{工资}, on voit \zh{工} (travail) ; la phrase parle de travail et d'argent, donc \zh{工资} (gōngzī) signifie « salaire ».
\end{enumerate}
\end{methode}

\begin{attention}[Ne lisez pas caractère par caractère !]
L'erreur classique du francophone débutant est de lire le chinois caractère par caractère, comme on épelait en CP : « nián... qīng... rén... ». Résultat : on perd le sens et la prononciation devient hachée. Lisez par \cle{groupes de sens} : \zh{年轻人} / \zh{离开} / \zh{家乡} / \zh{到大城市} / \zh{去打工}. C'est exactement comme en français : vous ne lisez pas « les-j-e-u-n-e-s », vous lisez « les jeunes / quittent / le village ». Entraînez-vous à découper le texte avec des barres obliques avant de lire à voix haute.
\end{attention}

\courssub{Traduction : littérale puis naturelle}

Comparons les deux niveaux de traduction sur des phrases clés.

\begin{center}
\begin{tabularx}{\linewidth}{|X|X|X|}
\hline
\rowcolor{popblueL} Chinois & Traduction littérale & Traduction naturelle \\
\hline
\zh{很多农村的年轻人离开家乡} \emph{hěn duō nóngcūn de niánqīngrén líkāi jiāxiāng} & « beaucoup campagne-de jeunes quitter village-natal » & « beaucoup de jeunes de la campagne quittent leur village natal » \\
\hline
\zh{农村的工作机会比较少} \emph{nóngcūn de gōngzuò jīhuì bǐjiào shǎo} & « campagne-de travail occasions relativement peu » & « les opportunités d'emploi sont assez rares à la campagne » \\
\hline
\zh{农村里只剩下老人和孩子} \emph{nóngcūn li zhǐ shèngxià lǎorén hé háizi} & « campagne-dans seulement rester vieillards et enfants » & « il ne reste à la campagne que les vieux et les enfants » \\
\hline
\zh{中国政府正在想办法} \emph{Zhōngguó zhèngfǔ zhèngzài xiǎng bànfǎ} & « Chine gouvernement en-train-de penser moyens » & « le gouvernement chinois cherche des solutions » \\
\hline
\end{tabularx}
\end{center}

Observez : le chinois dit « penser des moyens » (\zh{想办法}) là où le français dit « chercher des solutions » ; il dit « relativement peu » (\zh{比较少}) là où le français dit « assez rares ». Une bonne traduction restitue le \emph{sens}, pas les mots un à un.

\courssub{Repérage : mots connus et mots nouveaux}

\textbf{Mots déjà connus} (à surligner dans le texte) : \zh{人} rén (personne), \zh{工作} gōngzuò (travail), \zh{城市} chéngshì (ville), \zh{很多} hěn duō (beaucoup), \zh{为什么} wèishénme (pourquoi), \zh{因为} yīnwèi (parce que), \zh{但是} dànshì (mais), \zh{可以} kěyǐ (pouvoir), \zh{在} zài (à, être en train de).

\textbf{Mots nouveaux} du thème « exode rural » :

\begin{center}
\begin{tabularx}{\linewidth}{|X|X|X|X|}
\hline
\rowcolor{popblueL} Caractères & Pinyin & Nature & Traduction \\
\hline
农村 & nóngcūn & nom & campagne, monde rural \\
\hline
年轻人 & niánqīngrén & nom & jeune, jeunes gens \\
\hline
离开 & líkāi & verbe & quitter, partir de \\
\hline
家乡 & jiāxiāng & nom & village natal, pays natal \\
\hline
打工 & dǎgōng & verbe & travailler (comme ouvrier migrant) \\
\hline
机会 & jīhuì & nom & occasion, opportunité \\
\hline
工资 & gōngzī & nom & salaire \\
\hline
工厂 & gōngchǎng & nom & usine \\
\hline
公司 & gōngsī & nom & entreprise, société \\
\hline
人口外流 & rénkǒu wàiliú & expression & exode (rural), émigration \\
\hline
社会 & shèhuì & nom & société \\
\hline
政府 & zhèngfǔ & nom & gouvernement \\
\hline
办法 & bànfǎ & nom & moyen, solution \\
\hline
发展 & fāzhǎn & verbe/nom & (se) développer, développement \\
\hline
\end{tabularx}
\end{center}

\courssub{Segmentation : découper le texte en mots}

L'écriture chinoise ne met pas d'espace entre les mots : tout est collé. Le lecteur doit donc découper lui-même. Voici la première phrase segmentée :

\zh{现在 / 很多 / 农村 / 的 / 年轻人 / 离开 / 家乡，/ 到 / 大城市 / 去 / 打工。}

\emph{Xiànzài / hěn duō / nóngcūn / de / niánqīngrén / líkāi / jiāxiāng, / dào / dà chéngshì / qù / dǎgōng.}

Remarquez la logique de composition des mots :

\begin{itemize}
\item \zh{年轻人} = \zh{年轻} (jeune, adj.) + \zh{人} (personne) → « personne jeune »
\item \zh{农村} = \zh{农} (agriculture) + \zh{村} (village) → « village agricole », la campagne
\item \zh{人口外流} = \zh{人口} (population) + \zh{外} (extérieur) + \zh{流} (couler) → « la population coule vers l'extérieur » : l'exode !
\item \zh{打工} = \zh{打} (faire, frapper) + \zh{工} (travail) → « faire des boulots », travailler comme migrant
\end{itemize}

\begin{exoresolu}[Segmenter une phrase]
Énoncé : segmentez la phrase suivante en mots, puis traduisez : \zh{城市里有很多工厂和公司。}
\tcblower
Solution : \zh{城市 / 里 / 有 / 很多 / 工厂 / 和 / 公司。}\\
\emph{Chéngshì li yǒu hěn duō gōngchǎng hé gōngsī.}\\
« Dans les villes, il y a beaucoup d'usines et d'entreprises. »\\
Piège : ne pas découper \zh{城里 / 市有} ! \zh{城市} est un mot insécable (« ville ») et \zh{里} marque l'intérieur (« dans »). De même, \zh{工厂} et \zh{公司} sont des mots complets : on ne coupe jamais à l'intérieur.
\end{exoresolu}

\coursec{Clés (radicaux) et ordre des traits}

\courssub{Les radicaux des caractères du thème}

Chaque caractère chinois contient une \cle{clé} (ou radical), élément graphique qui donne souvent un indice de sens. Voici les clés des caractères importants de cette leçon :

\begin{center}
\begin{tabularx}{\linewidth}{|X|X|X|X|}
\hline
\rowcolor{popblueL} Caractère & Clé (radical) & Sens de la clé & Autres exemples \\
\hline
村 cūn (village) & 木 mù & bois, arbre & 树 shù (arbre), 机 jī (machine) \\
\hline
城 chéng (ville) & 土 tǔ & terre & 地 dì (terre), 场 chǎng (terrain, site) \\
\hline
流 liú (couler) & 氵 shuǐ & eau (trois gouttes) & 河 hé (rivière), 海 hǎi (mer) \\
\hline
机 jī (machine) & 木 mù & bois & 村 cūn (village), 校 xiào (école) \\
\hline
社 shè (société) & 礻 shì & autel, rite & 祝 zhù (souhaiter) \\
\hline
府 fǔ (gouvernement) & 广 guǎng & abri, bâtiment & 店 diàn (magasin) \\
\hline
离 lí (quitter) & 亠 tóu & couvercle (clé de tête) & 京 jīng (capitale) \\
\hline
\end{tabularx}
\end{center}

Observez la logique : les villages (\zh{村}) sont faits de bois (\zh{木}), les villes (\zh{城}) sont bâties sur la terre (\zh{土}), et ce qui « coule » (\zh{流}) porte les trois gouttes d'eau (\zh{氵}). Le radical est votre meilleur ami pour deviner le sens d'un caractère inconnu.

\courssub{Ordre des traits : quatre caractères à écrire}

Rappel des règles générales d'écriture : de \textbf{haut en bas}, de \textbf{gauche à droite}, l'\textbf{horizontal avant le vertical}, l'extérieur avant l'intérieur, et on « ferme » une boîte en dernier.

\textbf{1. 村 cūn (village) — 7 traits.} On écrit d'abord la clé \zh{木} à gauche (horizontal, vertical, puis les deux obliques), puis \zh{寸} à droite (horizontal, crochet vertical, point). Gauche avant droite.

\textbf{2. 城 chéng (ville) — 9 traits.} D'abord la clé \zh{土} à gauche (horizontal court, vertical, horizontal long), puis la partie droite \zh{成} (horizontal, oblique, puis les traits intérieurs, et enfin le long crochet \zh{戈} avec son point final).

\textbf{3. 流 liú (couler) — 10 traits.} D'abord les trois gouttes d'eau \zh{氵} à gauche (trois points, de haut en bas), puis la partie droite de haut en bas : le point et l'horizontal de tête, puis les traits du bas.

\textbf{4. 发 fā (dans 发展, développer) — 5 traits.} D'abord le trait brisé supérieur, puis l'oblique, puis les deux traits croisés du bas, et enfin le point en haut à droite. Le point final se trace toujours en dernier.

\begin{attention}[Pièges d'écriture]
\begin{itemize}
\item Ne confondez pas \zh{木} (mù, bois) et \zh{本} (běn, racine) : \zh{本} a un trait horizontal supplémentaire en bas.
\item La clé « eau » s'écrit \zh{氵} (trois gouttes) à gauche d'un caractère, mais \zh{水} quand elle est seule : comparez \zh{流} et \zh{水}.
\item \zh{发} (fā, 1\textsuperscript{er} ton, dans \zh{发展}) ne prend pas le même ton que dans \zh{头发} (tóufa, cheveux, ton neutre) : même caractère, deux lectures !
\end{itemize}
\end{attention}

\coursec{Compréhension du texte}

\courssub{Carte mentale et schéma du phénomène}

\begin{popfigure}
\begin{center}\tcbox[colback=popink,colframe=popink,coltext=white,arc=9pt,boxsep=4pt,left=6pt,right=6pt]{\bfseries\small 农村人口外流 \emph{rénkǒu wàiliú} exode rural}\end{center}
\vspace{-2pt}
\begin{tcbraster}[raster columns=2,raster equal height=rows,raster column skip=6pt,raster row skip=6pt,enhanced,blankest]
\begin{tcolorbox}[enhanced,colback=white,colframe=poporange,boxrule=0.9pt,arc=3pt,title={原因 causes},fonttitle=\bfseries\scriptsize,coltitle=white,colbacktitle=poporange,left=4pt,right=4pt,top=2pt,bottom=2pt,boxsep=1pt,toptitle=1.5pt,bottomtitle=1.5pt,halign=left]
\begin{itemize}[nosep,leftmargin=*,itemsep=1pt,font=\scriptsize]
\item 工作机会少 peu d'emplois
\item 工资不高 bas salaires
\end{itemize}
\end{tcolorbox}
\begin{tcolorbox}[enhanced,colback=white,colframe=poppink,boxrule=0.9pt,arc=3pt,title={结果 conséquences},fonttitle=\bfseries\scriptsize,coltitle=white,colbacktitle=poppink,left=4pt,right=4pt,top=2pt,bottom=2pt,boxsep=1pt,toptitle=1.5pt,bottomtitle=1.5pt,halign=left]
\begin{itemize}[nosep,leftmargin=*,itemsep=1pt,font=\scriptsize]
\item 只剩老人孩子 restent vieux et enfants
\item 社会问题 problème social
\end{itemize}
\end{tcolorbox}
\begin{tcolorbox}[enhanced,colback=white,colframe=popgreen,boxrule=0.9pt,arc=3pt,title={办法 solutions},fonttitle=\bfseries\scriptsize,coltitle=white,colbacktitle=popgreen,left=4pt,right=4pt,top=2pt,bottom=2pt,boxsep=1pt,toptitle=1.5pt,bottomtitle=1.5pt,halign=left]
\begin{itemize}[nosep,leftmargin=*,itemsep=1pt,font=\scriptsize]
\item 帮助农村发展 développer la campagne
\item 在家乡工作 travailler au village
\end{itemize}
\end{tcolorbox}
\end{tcbraster}

\legende{Carte mentale du texte : causes, conséquences et solutions de l'exode rural.}
\end{popfigure}

\begin{popfigure}
\begin{tikzpicture}[node distance=1.6cm, >=Stealth]
\node[draw=popgreen, fill=popgreenL, rounded corners, align=center, font=\small] (village) {农村\\ \emph{nóngcūn}\\ campagne};
\node[draw=poporange, fill=poporangeL, rounded corners, align=center, font=\small, right=of village] (depart) {年轻人离开\\ \emph{niánqīngrén líkāi}\\ les jeunes partent};
\node[draw=popblue, fill=popblueL, rounded corners, align=center, font=\small, right=of depart] (ville) {大城市\\ \emph{dà chéngshì}\\ grande ville};
\draw[->, very thick, popdark] (village) -- (depart);
\draw[->, very thick, popdark] (depart) -- (ville);
\end{tikzpicture}
\legende{Schéma fléché du mouvement : de la campagne vers la ville.}
\end{popfigure}

\courssub{Questions de compréhension}

\begin{exoresolu}[Répondre en chinois à une question de compréhension]
Énoncé : répondez en chinois, avec une phrase complète : \zh{年轻人为什么离开农村？} \emph{(Niánqīngrén wèishénme líkāi nóngcūn?)}
\tcblower
Solution : on reprend la structure de la question et on remplace \zh{为什么} par la cause introduite par \zh{因为} :\\
\zh{因为农村的工作机会比较少，工资也不高。}\\
\emph{Yīnwèi nóngcūn de gōngzuò jīhuì bǐjiào shǎo, gōngzī yě bù gāo.}\\
« Parce qu'à la campagne, les opportunités de travail sont assez rares et les salaires ne sont pas élevés non plus. »\\
Méthode : la réponse à une question en \zh{为什么} commence presque toujours par \zh{因为}. Il suffit de recopier la phrase du texte qui donne la cause.
\end{exoresolu}

\begin{exercice}[Questions de compréhension]
Répondez en chinois, par des phrases complètes tirées du texte.
\begin{enumerate}
\item \zh{年轻人到哪里去打工？} (Où les jeunes vont-ils travailler ?)
\item \zh{城市里有什么？} (Qu'y a-t-il dans les villes ?)
\item \zh{年轻人离开以后，农村里剩下谁？} (Après le départ des jeunes, qui reste à la campagne ?)
\item \zh{中国政府正在想办法做什么？} (Que cherche à faire le gouvernement chinois ?)
\end{enumerate}
\end{exercice}

\begin{corrige}[Questions de compréhension]
\begin{enumerate}
\item \zh{年轻人到大城市去打工。} \emph{Niánqīngrén dào dà chéngshì qù dǎgōng.} « Les jeunes vont travailler dans les grandes villes. »
\item \zh{城市里有很多工厂和公司。} \emph{Chéngshì li yǒu hěn duō gōngchǎng hé gōngsī.} « Dans les villes, il y a beaucoup d'usines et d'entreprises. »
\item \zh{农村里只剩下老人和孩子。} \emph{Nóngcūn li zhǐ shèngxià lǎorén hé háizi.} « Il ne reste à la campagne que les vieillards et les enfants. »
\item \zh{中国政府正在想办法帮助农村发展。} \emph{Zhōngguó zhèngfǔ zhèngzài xiǎng bànfǎ bāngzhù nóngcūn fāzhǎn.} « Le gouvernement chinois cherche des moyens d'aider la campagne à se développer. »
\end{enumerate}
\end{corrige}

\coursec{Collocations et exemples}

\courssub{Les collocations du thème}

Une \cle{collocation} est une combinaison de mots qui vont naturellement ensemble. En chinois comme en français, on ne combine pas les mots au hasard : on dit \zh{找工作} (« chercher un travail ») et non \zh{找工资}. Voici les collocations à mémoriser :

\begin{center}
\begin{tabularx}{\linewidth}{|X|X|X|}
\hline
\rowcolor{popblueL} Collocation & Pinyin & Traduction \\
\hline
离开家乡 & líkāi jiāxiāng & quitter son village natal \\
\hline
到大城市去 & dào dà chéngshì qù & aller dans les grandes villes \\
\hline
找工作 & zhǎo gōngzuò & chercher un emploi \\
\hline
找到工作 & zhǎodào gōngzuò & trouver un emploi \\
\hline
工作机会 & gōngzuò jīhuì & opportunité d'emploi \\
\hline
想办法 & xiǎng bànfǎ & chercher une solution \\
\hline
帮助农村发展 & bāngzhù nóngcūn fāzhǎn & aider la campagne à se développer \\
\hline
社会问题 & shèhuì wèntí & problème de société \\
\hline
\end{tabularx}
\end{center}

Remarquez la différence entre \zh{找} (zhǎo, chercher) et \zh{找到} (zhǎodào, trouver) : le verbe \zh{到} ajouté après un verbe d'action exprime le \emph{résultat atteint}. De même : \zh{看} (regarder) → \zh{看到} (voir, apercevoir) ; \zh{听} (écouter) → \zh{听到} (entendre).

\courssub{Exemples de réemploi}

\begin{exemplebox}[Le vocabulaire en action]
\zh{我的哥哥离开家乡，到杜阿拉去打工。}\\
\emph{Wǒ de gēge líkāi jiāxiāng, dào Dù'ālā qù dǎgōng.}\\
« Mon grand frère a quitté son village natal pour aller travailler à Douala. »\\[4pt]
\zh{他在一个工厂找到了工作，工资比较高。}\\
\emph{Tā zài yí ge gōngchǎng zhǎodào le gōngzuò, gōngzī bǐjiào gāo.}\\
« Il a trouvé un emploi dans une usine, le salaire est assez élevé. »\\[4pt]
\zh{但是农村里只剩下我的爷爷和奶奶。}\\
\emph{Dànshì nóngcūn li zhǐ shèngxià wǒ de yéye hé nǎinai.}\\
« Mais il ne reste au village que mon grand-père et ma grand-mère. »\\[4pt]
\zh{政府应该想办法帮助农村发展。}\\
\emph{Zhèngfǔ yīnggāi xiǎng bànfǎ bāngzhù nóngcūn fāzhǎn.}\\
« Le gouvernement devrait chercher des moyens d'aider la campagne à se développer. »
\end{exemplebox}

\begin{attention}[Erreurs fréquentes des francophones]
\begin{itemize}
\item \textbf{Place du complément de lieu} : en français, « il travaille à Douala » met le lieu après le verbe ; en chinois, le complément de lieu avec \zh{在} se place \emph{avant} le verbe : \zh{他在杜阿拉工作。} \emph{Tā zài Dù'ālā gōngzuò.} Ne dites jamais \zh{他工作杜阿拉}.
\item \textbf{\zh{离开} ne prend pas de \zh{从}} : le français dit « partir \emph{de} la campagne », mais le chinois dit directement \zh{离开农村} (sans préposition). Ne dites pas \zh{从农村离开}.
\item \textbf{\zh{比较} + adjectif} : \zh{比较少} signifie « relativement peu / assez rare », sans comparaison explicite. Pour une vraie comparaison, il faut \zh{比} : \zh{城市的工资比农村高。} « Les salaires en ville sont plus élevés qu'à la campagne. »
\end{itemize}
\end{attention}

\coursec{Production guidée}

\courssub{Lecture à voix haute : déroulement en classe}

\begin{enumerate}
\item \textbf{Lecture modèle par l'enseignant.} Le professeur lit le texte une première fois, à vitesse naturelle, les élèves écoutent le livre fermé ou les yeux sur le pinyin masqué. Puis deuxième lecture, les élèves suivent du doigt les caractères.
\item \textbf{Lecture chorale.} Toute la classe lit ensemble, phrase par phrase, en respectant les groupes de sens : \zh{很多农村的年轻人 / 离开家乡，/ 到大城市 / 去打工。}
\item \textbf{Lecture individuelle.} Chaque élève lit une phrase à tour de rôle ; le professeur corrige surtout les tons : \zh{nóngcūn} (2\textsuperscript{e} + 1\textsuperscript{er} ton), \zh{líkāi} (2\textsuperscript{e} + 1\textsuperscript{er} ton), \zh{jiāxiāng} (1\textsuperscript{er} + 1\textsuperscript{er} ton).
\end{enumerate}

\begin{attention}[Les tons des mots-clés]
Attention aux confusions de tons fréquentes chez les francophones : \zh{gōngzī} 工资 (salaire, tons 1-1) ne doit pas être prononcé avec un ton montant ; \zh{gōngchǎng} 工厂 (usine, 1-3) et \zh{gōngsī} 公司 (entreprise, 1-1) commencent tous deux par \zh{工} mais se distinguent à l'oreille par le second caractère. Répétez la série : 工作 gōngzuò, 工资 gōngzī, 工厂 gōngchǎng, 公司 gōngsī.
\end{attention}

\begin{experience}[Mini-exposé : l'exode rural au Cameroun]
Par groupes de trois, préparez un mini-exposé de cinq phrases en chinois sur l'exode rural au Cameroun, en suivant le plan du texte :
\begin{enumerate}
\item \textbf{Le phénomène} : \zh{很多年轻人离开农村，到杜阿拉和雅温得去打工。} \emph{Hěn duō niánqīngrén líkāi nóngcūn, dào Dù'ālā hé Yǎwēndé qù dǎgōng.} « Beaucoup de jeunes quittent la campagne pour aller travailler à Douala et à Yaoundé. »
\item \textbf{Les causes} : utilisez \zh{因为} + \zh{工作机会比较少} ou \zh{工资不高}.
\item \textbf{Les conséquences} : utilisez \zh{只剩下老人和孩子}.
\item \textbf{Une solution} : utilisez \zh{政府应该想办法帮助农村发展}.
\item Chaque membre du groupe lit au moins une phrase à voix haute ; la classe écoute et vérifie les tons.
\end{enumerate}
\end{experience}

\begin{savaistu}[Les 农民工 : les « ouvriers-paysans » de Chine]
En Chine, les travailleurs migrants venus de la campagne s'appellent \zh{农民工} (\emph{nóngmíngōng}), littéralement « paysans-ouvriers » : \zh{农} (agriculture) + \zh{民} (peuple) + \zh{工} (ouvrier). Ils sont des centaines de millions. Ils construisent les immeubles, travaillent dans les usines et livrent les colis dans les grandes villes comme Pékin, Shanghai ou Shenzhen. Chaque année, au moment du Nouvel An chinois, ils rentrent presque tous dans leur village : c'est la plus grande migration annuelle du monde, appelée \zh{春运} (\emph{chūnyùn}). Au Cameroun, on observe un mouvement comparable vers Douala et Yaoundé, même si les chiffres sont bien plus modestes.
\end{savaistu}

\begin{exercice}[À vous de jouer]
\begin{enumerate}
\item Retrouvez dans le texte trois mots composés contenant le caractère \zh{工} et donnez leur traduction.
\item Segmentez puis traduisez : \zh{他们为什么离开农村呢？}
\item Vrai ou faux (justifiez avec le texte) : a) Les jeunes partent parce que les salaires à la campagne sont élevés. b) Après le départ des jeunes, il reste les vieillards et les enfants au village.
\item Traduisez en chinois : « Les jeunes quittent leur village natal pour chercher du travail en ville. »
\end{enumerate}
\end{exercice}

\begin{corrige}[Exercice « À vous de jouer »]
\begin{enumerate}
\item \zh{工作} gōngzuò (travail), \zh{工资} gōngzī (salaire), \zh{打工} dǎgōng (travailler comme migrant), \zh{工厂} gōngchǎng (usine). Trois de ces quatre mots suffisent. Attention : \zh{公司} (entreprise) ne contient pas le caractère \zh{工}, il ne compte donc pas pour cette question.
\item \zh{他们 / 为什么 / 离开 / 农村 / 呢？} — \emph{Tāmen wèishénme líkāi nóngcūn ne?} — « Pourquoi quittent-ils la campagne ? »
\item a) Faux : le texte dit \zh{工资也不高} (les salaires ne sont pas élevés non plus). b) Vrai : \zh{农村里只剩下老人和孩子} (il ne reste que les vieillards et les enfants).
\item \zh{年轻人离开家乡，到城市里找工作。} \emph{Niánqīngrén líkāi jiāxiāng, dào chéngshì li zhǎo gōngzuò.} On réutilise les collocations \zh{离开家乡} et \zh{找工作}, et le complément de lieu \zh{到城市里} se place avant le verbe.
\end{enumerate}
\end{corrige}

\begin{aretenir}[L'essentiel de la leçon]
\begin{itemize}
\item L'\cle{exode rural} se dit \zh{农村人口外流} (\emph{nóngcūn rénkǒu wàiliú}) : « la population de la campagne coule vers l'extérieur ».
\item Le texte suit un plan en trois temps : \textbf{causes} (\zh{工作机会少，工资不高}), \textbf{conséquences} (\zh{只剩下老人和孩子}), \textbf{solutions} (\zh{想办法帮助农村发展}).
\item On lit un texte chinois par \textbf{groupes de sens}, jamais caractère par caractère ; on découpe en mots de deux caractères.
\item Les radicaux donnent des indices de sens : \zh{木} (bois) dans \zh{村}, \zh{土} (terre) dans \zh{城}, \zh{氵} (eau) dans \zh{流}.
\item Collocations à retenir : \zh{离开家乡}, \zh{找工作}, \zh{想办法}, \zh{帮助农村发展}.
\item Le complément de lieu avec \zh{在} se place \textbf{avant} le verbe ; \zh{离开} s'emploie \textbf{sans} préposition.
\end{itemize}
\end{aretenir}

\coursec{Leçon 23 — Vocabulaire et compréhension de texte : exode rural}

\courssub{Texte support : 农村的年轻人}

\begin{textebox}[Texte support : 农村的年轻人]
\zh{中国的农村有很多年轻人。他们喜欢自己的家乡，但是家乡的工厂不多，工资低，机会少。所以，很多年轻人离开家乡，去大城市打工。城市里机会多，工资高，生活却很忙，很累。农村人口外流很快，村子里只剩下老人和留守儿童。这是一个大的社会问题。政府想很多办法发展农村，希望年轻人能在家乡找到好工作，不必离开父母。}\\
\emph{Zhōngguó de nóngcūn yǒu hěn duō niánqīngrén. Tāmen xǐhuan zìjǐ de jiāxiāng, dànshì jiāxiāng de gōngchǎng bù duō, gōngzī dī, jīhuì shǎo. Suǒyǐ, hěn duō niánqīngrén líkāi jiāxiāng, qù dà chéngshì dǎgōng. Chéngshì lǐ jīhuì duō, gōngzī gāo, shēnghuó què hěn máng, hěn lèi. Nóngcūn rénkǒu wàiliú hěn kuài, cūnzi lǐ zhǐ shèngxià lǎorén hé liúshǒu értóng. Zhè shì yí ge dà de shèhuì wèntí. Zhèngfǔ xiǎng hěn duō bànfǎ fāzhǎn nóngcūn, xīwàng niánqīngrén néng zài jiāxiāng zhǎodào hǎo gōngzuò, bù bì líkāi fùmǔ.}\\
« La campagne chinoise compte de nombreux jeunes. Ils aiment leur pays natal, mais il y a peu d'usines, les salaires sont bas et les opportunités rares. C'est pourquoi beaucoup de jeunes quittent leur village pour aller travailler dans les grandes villes. En ville, les opportunités sont nombreuses et les salaires élevés, mais la vie est très occupée et fatigante. L'exode rural s'accélère, il ne reste au village que les personnes âgées et les enfants laissés au village. C'est un grand problème de société. Le gouvernement cherche de nombreuses solutions pour développer la campagne, dans l'espoir que les jeunes puissent trouver un bon travail dans leur pays natal, sans avoir à quitter leurs parents. »
\end{textebox}

\courssub{Glossaire du thème}

Dans la section précédente, nous avons lu le texte \zh{农村的年轻人} (\emph{Nóngcūn de niánqīngrén}, « Les jeunes de la campagne »). Ce texte raconte le départ des jeunes villageois vers les grandes villes, un phénomène que l'on observe aussi bien en Chine qu'au Cameroun, où beaucoup de jeunes quittent leurs villages pour Yaoundé ou Douala. Pour bien comprendre ce texte et pouvoir en parler, il faut d'abord maîtriser son vocabulaire. C'est l'objet de cette section : apprendre les mots, les prononcer correctement, comprendre leur emploi exact et savoir les réutiliser dans nos propres phrases.

\courssub{Le tableau de vocabulaire principal}

Observez d'abord le tableau ci-dessous. Lisez chaque mot à voix haute en suivant le pinyin, puis vérifiez la traduction. Remarquez la colonne « nature » : en chinois comme en français, connaître la nature grammaticale d'un mot (nom, verbe) est indispensable pour le placer correctement dans la phrase.

\begin{center}
\begin{tabularx}{\linewidth}{|X|X|X|X|}
\hline
\rowcolor{popblueL} Caractères & Pinyin & Nature & Traduction \\
\hline
农村 & nóngcūn & n. & campagne, monde rural \\
\hline
城市 & chéngshì & n. & ville \\
\hline
家乡 & jiāxiāng & n. & pays natal, village natal \\
\hline
离开 & líkāi & v. & quitter \\
\hline
外流 & wàiliú & v./n. & exode, flux vers l'extérieur \\
\hline
打工 & dǎgōng & v. & travailler (comme ouvrier migrant) \\
\hline
机会 & jīhuì & n. & opportunité, occasion \\
\hline
工资 & gōngzī & n. & salaire \\
\hline
工厂 & gōngchǎng & n. & usine \\
\hline
人口 & rénkǒu & n. & population \\
\hline
社会 & shèhuì & n. & société \\
\hline
问题 & wèntí & n. & problème, question \\
\hline
发展 & fāzhǎn & v./n. & (se) développer ; développement \\
\hline
办法 & bànfǎ & n. & solution, moyen \\
\hline
剩下 & shèngxià & v. & rester, subsister \\
\hline
\end{tabularx}
\end{center}

Pour l'examen, deux expressions composées sont très utiles et fréquentes dans les médias :

\begin{center}
\begin{tabularx}{\linewidth}{|X|X|X|X|}
\hline
\rowcolor{popblueL} Caractères & Pinyin & Nature & Traduction \\
\hline
移民 & yímín & n./v. & migrant ; émigrer \\
\hline
留守儿童 & liúshǒu értóng & expr. & enfants laissés au village (par leurs parents partis travailler en ville) \\
\hline
\end{tabularx}
\end{center}

\courssub{Observation : les mots dans leur contexte}

Avant d'énoncer les règles d'emploi, observons comment ces mots fonctionnent dans des phrases réelles tirées du thème du texte.

\begin{exemplebox}[Les mots du glossaire en contexte]
\zh{他离开家乡去城市打工。}\\
\emph{Tā líkāi jiāxiāng qù chéngshì dǎgōng.}\\
« Il a quitté son village pour aller travailler en ville. »\\[4pt]
\zh{城市里工作机会多，工资也高。}\\
\emph{Chéngshì lǐ gōngzuò jīhuì duō, gōngzī yě gāo.}\\
« En ville, les opportunités de travail sont nombreuses et les salaires sont élevés aussi. »\\[4pt]
\zh{农村人口外流是一个社会问题。}\\
\emph{Nóngcūn rénkǒu wàiliú shì yí ge shèhuì wèntí.}\\
« L'exode de la population rurale est un problème de société. »\\[4pt]
\zh{年轻人都走了，村子里只剩下老人和孩子。}\\
\emph{Niánqīngrén dōu zǒu le, cūnzi lǐ zhǐ shèngxià lǎorén hé háizi.}\\
« Les jeunes sont tous partis, il ne reste au village que les vieillards et les enfants. »
\end{exemplebox}

Observations grammaticales :

\begin{itemize}
\item \zh{离开} (líkāi) est un verbe \cle{transitif direct} : on quitte quelque chose, sans préposition. On dit \zh{离开家乡} (quitter son village), jamais \zh{从家乡离开} à ce niveau. Structure : Sujet + 离开 + Lieu.
\item \zh{去 + Lieu + Verbe} exprime le but du déplacement : \zh{去城市打工} « aller en ville pour travailler ». C'est la construction en \cle{série verbale} : deux verbes se suivent sans mot de liaison (pas de « pour »).
\item \zh{剩下} (shèngxià) signifie « il reste / subsister » : il s'emploie souvent avec \zh{只} (zhǐ, seulement) placé avant le verbe : \zh{只剩下老人} « il ne reste que les vieillards ».
\end{itemize}

\courssub{Le mot-clé du thème : 外流}

\begin{definition}[外流 wàiliú]
\zh{外流} (wàiliú) est composé de \zh{外} (wài, extérieur) et \zh{流} (liú, couler, flux). Il désigne un \cle{mouvement de population qui quitte une zone} pour aller ailleurs. L'expression \zh{农村人口外流} (nóngcūn rénkǒu wàiliú), littéralement « flux vers l'extérieur de la population rurale », se traduit en français par \emph{exode rural}. On rencontre aussi \zh{人才外流} (réncái wàiliú), « la fuite des cerveaux / exode des talents ».
\end{definition}

Comparaison avec le français : le français utilise un nom abstrait d'origine grecque (« exode »), le chinois construit un mot transparent dont chaque caractère garde son sens concret : « extérieur + couler ». C'est une grande force du vocabulaire chinois : en connaissant les caractères, on devine le sens des mots composés.

\courssub{Attention aux faux amis et aux emplois précis}

\begin{attention}[打工 ne veut pas dire « travailler » en général]
\zh{打工} (dǎgōng) désigne le \cle{travail salarié précaire des migrants} : ouvrier d'usine, serveur, ouvrier du bâtiment, loin de chez soi. Ne l'utilisez jamais pour un emploi stable, de cadre ou de fonctionnaire. Pour « travailler » au sens général, employez \zh{工作} (gōngzuò).\\
Correct : \zh{他在工厂打工。} \emph{Tā zài gōngchǎng dǎgōng.} « Il travaille comme ouvrier dans une usine. »\\
Incorrect pour un enseignant : on dira \zh{他是老师，在学校工作。} \emph{Tā shì lǎoshī, zài xuéxiào gōngzuò.} « Il est professeur, il travaille dans une école. »
\end{attention}

\begin{attention}[家乡, 农村, 城市 : ne pas confondre les référents]
\zh{家乡} (jiāxiāng) est le \emph{pays natal}, le lieu d'où l'on vient (ville ou village) ; \zh{农村} (nóngcūn) est la \emph{campagne} par opposition à \zh{城市} (chéngshì), la \emph{ville}. Un habitant de Douala peut dire que Douala est son \zh{家乡}, mais pas que Douala est un \zh{农村} ! De même, \zh{问题} (wèntí) signifie à la fois « problème » et « question » (à l'école) : c'est le contexte qui décide.
\end{attention}

\begin{attention}[工资 vs 工作 : distinction nom / verbe]
\zh{工资} (gōngzī) est un \cle{nom} (salaire). On dit \zh{工资高} (salaire élevé) ou \zh{工资低} (salaire bas). \zh{工作} (gōngzuò) est principalement un \cle{verbe} (travailler) ou un nom (emploi, travail). Ne confondez pas « chercher du travail » (\zh{找工作}) et « chercher un salaire » (\zh{找工资} — absurde).
\end{attention}

\courssub{Méthode de mémorisation : apprendre par familles de caractères}

\begin{methode}[Regrouper les mots par famille de caractères]
Le moyen le plus efficace de retenir ce vocabulaire est de le classer autour d'un caractère commun (clé sémantique).
\begin{enumerate}
\item Repérez le caractère qui revient : ici \zh{工} (gōng, travail, ouvrage).
\item Listez tous les mots de la famille du thème : \zh{工作} (gōngzuò, travailler), \zh{打工} (dǎgōng, travailler comme migrant), \zh{工资} (gōngzī, salaire), \zh{工厂} (gōngchǎng, usine), \zh{工人} (gōngrén, ouvrier).
\item Construisez une phrase personnelle avec chaque mot pour fixer son emploi.
\item Faites de même avec d'autres familles du thème : \zh{人} donne \zh{人口} (population), \zh{人们} (les gens), \zh{农民} (paysan) ; \zh{发} donne \zh{发展} (développement), \zh{发现} (découvrir).
\end{enumerate}
\end{methode}

\begin{popfigure}
\begin{center}\tcbox[colback=popblue,colframe=popblue,coltext=white,arc=9pt,boxsep=4pt,left=6pt,right=6pt]{\bfseries\small 工 gōng travail}\end{center}
\vspace{-2pt}
\begin{tcbraster}[raster columns=3,raster equal height=rows,raster column skip=6pt,raster row skip=6pt,enhanced,blankest]
\begin{tcolorbox}[enhanced,colback=poporangeL,colframe=poporangeL,coltext=white,boxrule=0.9pt,arc=3pt,left=4pt,right=4pt,top=3pt,bottom=3pt,boxsep=1pt,halign=left]\bfseries\scriptsize 工作 gōngzuò travailler\end{tcolorbox}
\begin{tcolorbox}[enhanced,colback=popgreenL,colframe=popgreenL,coltext=white,boxrule=0.9pt,arc=3pt,left=4pt,right=4pt,top=3pt,bottom=3pt,boxsep=1pt,halign=left]\bfseries\scriptsize 打工 dǎgōng travail migrant\end{tcolorbox}
\begin{tcolorbox}[enhanced,colback=poppinkL,colframe=poppinkL,coltext=white,boxrule=0.9pt,arc=3pt,left=4pt,right=4pt,top=3pt,bottom=3pt,boxsep=1pt,halign=left]\bfseries\scriptsize 工资 gōngzī salaire\end{tcolorbox}
\begin{tcolorbox}[enhanced,colback=popgoldL,colframe=popgoldL,coltext=white,boxrule=0.9pt,arc=3pt,left=4pt,right=4pt,top=3pt,bottom=3pt,boxsep=1pt,halign=left]\bfseries\scriptsize 工厂 gōngchǎng usine\end{tcolorbox}
\begin{tcolorbox}[enhanced,colback=poppurpleL,colframe=poppurpleL,coltext=white,boxrule=0.9pt,arc=3pt,left=4pt,right=4pt,top=3pt,bottom=3pt,boxsep=1pt,halign=left]\bfseries\scriptsize 工人 gōngrén ouvrier\end{tcolorbox}
\end{tcbraster}

\legende{Carte mentale de la famille du caractère 工 (gōng, travail) : un caractère, cinq mots.}
\end{popfigure}

\courssub{Entraînement à la prononciation : les tons}

Plusieurs mots nouveaux présentent des combinaisons de tons difficiles pour les francophones. Répétez chaque mot lentement, puis à vitesse normale. Respectez l'espace entre les mots en pinyin.

\begin{center}
\begin{tabularx}{\linewidth}{|X|X|X|}
\hline
\rowcolor{popblueL} Mot & Tons & Conseil de prononciation \\
\hline
农村 nóngcūn & 2--1 & Monter sur \emph{nóng} (ton montant), puis descendre d'un coup haut et plat sur \emph{cūn} (ton haut). \\
\hline
离开 líkāi & 2--1 & Même schéma : montée puis ton haut tendu. \\
\hline
工资 gōngzī & 1--1 & Deux tons hauts plats : garder la voix haute et égale, sans mélodie descendante. \\
\hline
发展 fāzhǎn & 1--3 & Ton haut plat, puis descendre bas et remonter (\emph{zhǎn}, ton 3 bas-remontant). \\
\hline
城市 chéngshì & 2--4 & Monter sur \emph{chéng}, puis tomber brusquement sur \emph{shì} (ton descendant). \\
\hline
剩下 shèngxià & 4--4 & Deux tons descendants, secs et décidés, comme un ordre ferme. \\
\hline
留守儿童 liúshǒu értóng & 2--3--2--2 & Attention au \emph{shǒu} (ton 3) et au \emph{ér} (ton 2, souvent neutre en pratique : \emph{ertong}). \\
\hline
\end{tabularx}
\end{center}

\begin{attention}[Piège fréquent : le ton 2 suivi du ton 1]
Les francophones ont tendance à prononcer \zh{农村} (nóngcūn) avec deux tons plats (1--1), ce qui le fait ressembler à d'autres mots. Entraînez-vous sur la paire minimale : \zh{农村} (nóngcūn, 2--1, la campagne) $\neq$ \zh{浓春} (nóngchūn, 2--1, printemps dense - rare). De même, \zh{工资} (gōngzī, 1--1) exige deux tons hauts bien plats : ne laissez pas la voix retomber sur la deuxième syllabe (évitez 1--4 ou 1--3).
\end{attention}

\courssub{Radicaux et ordre des traits}

Connaître le radical (clé) et l'ordre des traits permet de chercher un caractère dans le dictionnaire et de l'écrire correctement (équilibre, vitesse). Voici les caractères clés du thème.

\begin{center}
\begin{tabularx}{\linewidth}{|X|X|X|X|X|}
\hline
\rowcolor{popblueL} Caractère & Pinyin & Radical (Clé) & Nb traits & Ordre des traits (règles principales) \\
\hline
农 & nóng & \zh{曲} (qū, courbe) & 7 & Horizontal, vertical, horizontal, vertical, horizontal, horizontal, crochet final. \\
\hline
村 & cūn & \zh{木} (mù, bois/arbre) & 7 & Écrire le radical \zh{木} (haut→bas, gauche→droite), puis \zh{寸} à droite. \\
\hline
城 & chéng & \zh{土} (tǔ, terre) & 10 & Radical \zh{土} à gauche (horizontal, vertical, horizontal), puis \zh{成} à droite. \\
\hline
市 & shì & \zh{巾} (jīn, étoffe) & 5 & Vertical, horizontal (x3), vertical traversant. (Caractère simple, pas de clé séparée visible). \\
\hline
工 & gōng & \zh{工} (gōng, travail) & 3 & Horizontal (court), vertical, horizontal (long). Base de \zh{工厂}, \zh{工资}... \\
\hline
厂 & chǎng & \zh{厂} (hǎn, falaise/usine) & 2 & Horizontal, puis pli vertical-horizontal (comme un L retourné). \\
\hline
留 & liú & \zh{田} (tián, champ) & 10 & Haut \zh{刘} (simplifié) : \zh{田} + \zh{卯} (simplifié en bas). Écrire \zh{田} puis le bas. \\
\hline
守 & shǒu & \zh{宀} (mián, toit) & 6 & Point, pli, point (toit), puis \zh{寸} en dessous. \\
\hline
\end{tabularx}
\end{center}

\begin{aretenir}[Règles d'or de l'ordre des traits]
\begin{itemize}
\item De haut en haut (\zh{三}), de gauche à droite (\zh{川}), horizontal avant vertical (\zh{十}), extérieur avant intérieur (\zh{同}), fermer en dernier (\zh{口}, \zh{田}).
\item Pour \zh{村} : clé \zh{木} (4 traits) + phonétique \zh{寸} (3 traits) = 7 traits.
\item Pour \zh{城} : clé \zh{土} (3 traits) + phonétique \zh{成} (7 traits) = 10 traits.
\end{itemize}
\end{aretenir}

\courssub{Compréhension du texte support}

Répondez aux questions suivantes en français (ou en chinois si vous le pouvez) en vous basant \emph{uniquement} sur le texte \zh{农村的年轻人}.

\begin{exercice}[Questions de compréhension]
\begin{enumerate}
\item Pourquoi les jeunes aiment-ils leur maison natale mais la quittent-ils quand même ? Citez trois raisons données par le texte.
\item Où vont les jeunes ? Quel type de travail font-ils là-bas (vocabulaire précis) ?
\item Quelle est la conséquence démographique pour le village ? Quel terme chinois désigne les enfants restés au village ?
\item Quel est l'avis du gouvernement sur cette situation ? Que compte-t-il faire ?
\item Traduisez la phrase : \zh{农村人口外流很快，村子里只剩下老人和留守儿童。}
\end{enumerate}
\end{exercice}

\begin{corrige}[Exercice « Compréhension du texte »]
\begin{enumerate}
\item Raisons : peu d'usines (\zh{工厂不多}), salaires bas (\zh{工资低}), peu d'opportunités (\zh{机会少}).
\item Ils vont dans les grandes villes (\zh{大城市}) pour \zh{打工} (travailler comme migrants/ouvriers).
\item La population diminue vite (\zh{外流很快}) ; il ne reste que les vieux et les \zh{留守儿童} (enfants laissés au village).
\item Le gouvernement considère cela comme un grand problème social (\zh{大的社会问题}) et cherche des moyens (\zh{想很多办法}) pour développer la campagne (\zh{发展农村}) afin que les jeunes puissent travailler sur place.
\item « L'exode de la population rurale est très rapide, il ne reste dans le village que les personnes âgées et les enfants laissés au village. »
\end{enumerate}
\end{corrige}

\courssub{Collocations et exemples (Structures lexicales)}

Maîtriser un mot, c'est savoir avec quels autres mots il s'associe naturellement (collocations). Voici les associations fréquentes du thème.

\begin{propriete}[Collocations verbales et nominales du thème]
\begin{center}
\begin{tabularx}{\linewidth}{|X|X|X|}
\hline
\rowcolor{popblueL} Structure & Exemple (Caractères + Pinyin) & Traduction \\
\hline
离开 + Lieu & \zh{离开家乡} (líkāi jiāxiāng) & Quitter son pays natal \\
\hline
去 + Lieu + 打工 & \zh{去城市打工} (qù chéngshì dǎgōng) & Aller en ville travailler comme migrant \\
\hline
Adj + 机会 & \zh{很多机会} (hěn duō jīhuì) / \zh{机会少} (jīhuì shǎo) & Beaucoup / peu d'opportunités \\
\hline
工资 + Adj & \zh{工资高} (gōngzī gāo) / \zh{工资低} (gōngzī dī) & Salaire élevé / bas \\
\hline
人口 + 外流 & \zh{农村人口外流} (nóngcūn rénkǒu wàiliú) & Exode rural \\
\hline
社会 + 问题 & \zh{社会问题} (shèhuì wèntí) & Problème de société \\
\hline
想 + 办法 & \zh{想办法} (xiǎng bànfǎ) & Chercher une solution / un moyen \\
\hline
发展 + Lieu/Objet & \zh{发展农村} (fāzhǎn nóngcūn) / \zh{发展经济} (fāzhǎn jīngjì) & Développer la campagne / l'économie \\
\hline
只 + 剩下 + N & \zh{只剩下老人} (zhǐ shèngxià lǎorén) & Il ne reste que les vieillards \\
\hline
\end{tabularx}
\end{center}
\end{propriete}

\begin{exemplebox}[Phrases modèles pour l'expression écrite/orale]
\zh{为了赚钱，他不得不离开家乡。}\\
\emph{Wèile zhuànqián, tā bùdébù líkāi jiāxiāng.}\\
« Pour gagner de l'argent, il a dû quitter son pays natal. »\\[4pt]
\zh{解决农村人口外流的问题，必须发展农村经济。}\\
\emph{Jiějué nóngcūn rénkǒu wàiliú de wèntí, bìxū fāzhǎn nóngcūn jīngjì.}\\
« Pour résoudre le problème de l'exode rural, il faut développer l'économie rurale. »\\[4pt]
\zh{留守儿童很想念父母，一年只能见一次。}\\
\emph{Liúshǒu értóng hěn xiǎngniàn fùmǔ, yì nián zhǐ néng jiàn yí cì.}\\
« Les enfants laissés au village manquent beaucoup à leurs parents, ils ne peuvent les voir qu'une fois par an. »
\end{exemplebox}

\courssub{Production guidée}

\courssub{Expression écrite : Message / Petit paragraphe}

\begin{methode}[Rédiger un paragraphe sur l'exode rural (80--100 caractères)]
\begin{enumerate}
\item \textbf{Situation initiale} : Présentez le village (nom, situation, population). Utilisez \zh{是}, \zh{有}.
\item \textbf{Problème} : Les jeunes partent. Utilisez \zh{离开}, \zh{去}, \zh{打工}, \zh{因为}... \zh{所以}...
\item \textbf{Conséquence} : Qui reste ? Utilisez \zh{只剩下}, \zh{留守儿童}.
\item \textbf{Solution / Avenir} : Ce que fait le gouvernement / espoir. Utilisez \zh{想办法}, \zh{发展}, \zh{希望}.
\item \textbf{Relisez-vous} : Vérifiez l'ordre des mots (Sujet - Temps - Lieu - Verbe - Complément), les tons dans votre tête, la ponctuation chinoise 。
\end{enumerate}
\end{methode}

\begin{exemplebox}[Modèle de rédaction (niveau HSK 3)]
\zh{我的家乡是一个小村子，叫马鲁阿。以前村子里人很多，现在年轻人都离开家乡，去杜阿拉打工，因为城市里工资高，机会多。村子里只剩下老人和留守儿童，这是一个大问题。政府正在想办法发展农村，修路、建厂。我希望将来年轻人能在家乡找到好工作，不用离开父母。}\\
\emph{Wǒ de jiāxiāng shì yí ge xiǎo cūnzi, jiào Mǎlǔā. Yǐqián cūnzi lǐ rén hěn duō, xiànzài niánqīngrén dōu líkāi jiāxiāng, qù Dùālā dǎgōng, yīnwèi chéngshì lǐ gōngzī gāo, jīhuì duō. Cūnzi lǐ zhǐ shèngxià lǎorén hé liúshǒu értóng, zhè shì yí ge dà wèntí. Zhèngfǔ zhèngzài xiǎng bànfǎ fāzhǎn nóngcūn, xiū lù, jiàn chǎng. Wǒ xīwàng jiānglái niánqīngrén néng zài jiāxiāng zhǎodào hǎo gōngzuò, bùyòng líkāi fùmǔ.}
\end{exemplebox}

\courssub{Expression orale : Jeu de rôle et débat}

\begin{experience}[Jeu de rôle : « Le départ »]
\textbf{Contexte :} Un jeune (A) annonce à son grand-père (B) qu'il part pour Douala/Yaoundé travailler dans une usine.
\textbf{Rôles :}
\begin{itemize}
\item \textbf{A (Le jeune) :} Explique \emph{pourquoi} (salaire bas au village, pas d'usine, envie d'aider la famille). Utilise : \zh{工资低}, \zh{没机会}, \zh{想赚钱}, \zh{打工}.
\item \textbf{B (Le grand-père) :} Exprime son inquiétude (ville dangereuse, vie dure, manque de main-d'œuvre au village, petits-enfants \zh{留守儿童}). Utilise : \zh{城市生活苦}, \zh{担心}, \zh{留守儿童}, \zh{希望你回来}.
\end{itemize}
\textbf{Consigne :} Préparez 2 minutes, jouez 3 minutes. Changez de rôle. Le professeur note le vocabulaire du thème utilisé correctement.
\end{experience}

\begin{experience}[Débat guidé : « Faut-il partir ou rester ? »]
\textbf{Sujet :} \zh{年轻人应该留在农村发展，还是去城市打工好？} (Les jeunes doivent-ils rester développer la campagne ou aller travailler en ville ?)
\textbf{Groupes :} Groupe 1 (Partir) / Groupe 2 (Rester).
\textbf{Arguments suggérés :}
\begin{itemize}
\item \textbf{Partir :} \zh{工资高}, \zh{见世面} (voir le monde/élargir ses horizons), \zh{寄钱回家} (envoyer de l'argent à la maison), \zh{学技术} (apprendre un métier).
\item \textbf{Rester :} \zh{陪伴父母} (accompagner ses parents), \zh{建设家乡} (construire son pays natal), \zh{生活成本低} (coût de la vie bas), \zh{政策扶持} (soutien politique/gouvernemental).
\end{itemize}
\textbf{Sortie langagière attendue :} Utiliser \zh{我认为...因为...而且...但是...所以...} (Je pense que... parce que... de plus... mais... donc...).
\end{experience}

\courssub{Élément socioculturel : Cameroun -- Chine}

\begin{savaistu}[L'exode rural : regards croisés Cameroun -- Chine]
\begin{itemize}
\item \textbf{Chine :} Phénomène massif depuis les années 1980 (\zh{农民工} -- travailleurs migrants paysans). \zh{户口} (hùkǒu, registre de résidence) crée une distinction administrative : les migrants ruraux n'ont pas les mêmes droits sociaux (école, santé, logement) en ville que les urbains. \zh{留守儿童} (enfants laissés au village) : plus de 60 millions (chiffres récents). Politique actuelle : \zh{城镇化} (urbanisation) + \zh{乡村振兴} (revitalisation rurale) pour créer des emplois sur place.
\item \textbf{Cameroun :} Exode vers Douala, Yaoundé, Bafoussam, Ngaoundéré. Causes : chômage rural, manque d'infrastructures (routes, électricité, eau), attraction du secteur informel urbain. Conséquences : vieillissement des villages, pression sur l'urbain (bidonvilles, chômage urbain). Politiques : \emph{Programme d'urgence triennal}, \emph{Plan d'urgence pour la jeunesse}, décentralisation.
\item \textbf{Point commun :} Dans les deux pays, le \cle{développement rural} (amélioration des revenus agricoles, agro-industrie, numérique, tourisme) est la clé durable pour freiner l'exode.
\end{itemize}
\end{savaistu}

\courssub{Entraînement et évaluation}

\begin{exoresolu}[Complétez avec le mot du glossaire qui convient]
Complétez : \zh{很多年轻人离开农村，去城市\_\_\_\_，因为那里的\_\_\_\_比较高。（a. 工资 b. 打工）}
\tcblower
\textbf{Solution détaillée.} Premier blanc : après \zh{去城市}, il faut un verbe exprimant le but du déplacement ; seul \zh{打工} (dǎgōng, travailler comme migrant) est un verbe. Deuxième blanc : avant \zh{比较高} (relativement élevé), il faut un nom qui peut être « élevé » ; \zh{工资} (gōngzī, salaire) convient, un salaire peut être \zh{高} (élevé).\\
Phrase complète : \zh{很多年轻人离开农村，去城市打工，因为那里的工资比较高。}\\
\emph{Hěn duō niánqīngrén líkāi nóngcūn, qù chéngshì dǎgōng, yīnwèi nàlǐ de gōngzī bǐjiào gāo.}\\
« Beaucoup de jeunes quittent la campagne pour aller travailler en ville, parce que les salaires y sont relativement élevés. »
\end{exoresolu}

\begin{exercice}[Entraînement 1 : Traduction vers le chinois (Thème)]
Traduisez en chinois en utilisant \textbf{impérativement} le vocabulaire du glossaire.
\begin{enumerate}
\item Mon frère aîné a quitté la campagne pour aller travailler à Douala.
\item L'exode rural est un problème social sérieux.
\item Il ne reste que les personnes âgées et les enfants dans le village.
\item Le gouvernement cherche des moyens pour développer l'économie rurale.
\item Il y a peu d'usines dans mon pays natal, les salaires sont bas.
\end{enumerate}
\end{exercice}

\begin{corrige}[Exercice « Entraînement 1 »]
\begin{enumerate}
\item \zh{我哥哥离开农村，去杜阿拉打工。} \emph{Wǒ gēge líkāi nóngcūn, qù Dùālā dǎgōng.}
\item \zh{农村人口外流是一个严重的社会问题。} \emph{Nóngcūn rénkǒu wàiliú shì yí ge yánzhòng de shèhuì wèntí.} (Ajout de \zh{严重} yánzhòng, sérieux, niveau HSK 3).
\item \zh{村子里只剩下老人和孩子。} \emph{Cūnzi lǐ zhǐ shèngxià lǎorén hé háizi.} (Ou \zh{留守儿童}).
\item \zh{政府想办法发展农村经济。} \emph{Zhèngfǔ xiǎng bànfǎ fāzhǎn nóngcūn jīngjì.}
\item \zh{我家乡的工厂不多，工资低。} \emph{Wǒ jiāxiāng de gōngchǎng bù duō, gōngzī dī.}
\end{enumerate}
\end{corrige}

\begin{exercice}[Entraînement 2 : Version (Chinois $\rightarrow$ Français)]
Traduisez en français :
\begin{enumerate}
\item \zh{现在的年轻人不想在农村种地，都想去城市打工。}
\item \zh{解决留守儿童的问题，要从发展农村入手。}
\item \zh{虽然城市工资高，但生活压力也大，很累。}
\end{enumerate}
\end{exercice}

\begin{corrige}[Exercice « Entraînement 2 »]
\begin{enumerate}
\item « Les jeunes d'aujourd'hui ne veulent pas cultiver la terre à la campagne, ils veulent tous aller travailler en ville. »
\item « Pour résoudre le problème des enfants laissés au village, il faut commencer par développer la campagne. »
\item « Bien que les salaires en ville soient élevés, la pression de la vie y est aussi grande, c'est très fatigant. »
\end{enumerate}
\end{corrige}

\begin{exercice}[Entraînement 3 : Production libre (Devoir / Évaluation)]
\textbf{Sujet :} « Dans votre région, les jeunes partent-ils pour la ville ? Écrivez un court texte (6--8 phrases) pour décrire la situation, les causes, les conséquences et vos espoirs pour l'avenir. Utilisez au moins 8 mots du glossaire. »
\textbf{Mots obligatoires à intégrer (cochez-les) :} \zh{农村}, \zh{城市}, \zh{家乡}, \zh{离开}, \zh{打工}, \zh{工资}, \zh{机会}, \zh{剩下}, \zh{留守儿童}, \zh{发展}, \zh{办法}, \zh{问题}.
\end{exercice}

\begin{aretenir}[L'essentiel de la Leçon 23]
\begin{itemize}
\item \textbf{Thème :} Exode rural (\zh{农村人口外流}). Trois pôles lexicaux : Origine (\zh{农村}, \zh{家乡}) -- Départ (\zh{离开}, \zh{外流}, \zh{打工}) -- Destination (\zh{城市}, \zh{工厂}, \zh{工资}, \zh{机会}).
\item \textbf{Grammaire clé :} \zh{离开} (verbe transitif direct) ; \zh{去 + Lieu + V} (but) ; \zh{只 + 剩下} (restriction) ; \zh{想办法 + V} (chercher à).
\item \textbf{Distinction cruciale :} \zh{打工} (travail migrant précaire) $\neq$ \zh{工作} (travailler en général / emploi).
\item \textbf{Écriture :} Maîtriser l'ordre des traits des radicaux \zh{木} (村), \zh{土} (城), \zh{工}, \zh{厂}, \zh{宀} (守).
\item \textbf{Prononciation :} Paires tonales 2--1 (\zh{农村}, \zh{离开}), 1--1 (\zh{工资}), 4--4 (\zh{剩下}), 2--4 (\zh{城市}).
\item \textbf{Culture :} Phénomène commun Chine / Cameroun. Enjeu : \zh{留守儿童} / enfants laissés au village. Solution commune : \zh{发展农村} / développement rural.
\end{itemize}
\end{aretenir}


\coursec{Radicaux et ordre des traits}

Dans le glossaire de la section précédente, nous avons rencontré des mots comme \zh{农村} (campagne), \zh{城市} (ville), \zh{流动} (circuler), \zh{离开} (quitter) et \zh{工作} (travailler). Ces mots sont composés de caractères qui ne sont pas des dessins au hasard : chacun est construit à partir d'une \cle{clé} (ou \cle{radical}, \zh{部首} \emph{bùshǒu}) qui donne une indication de sens, et s'écrit selon un \cle{ordre des traits} (\zh{笔顺} \emph{bǐshùn}) précis et immuable. Dans cette section, nous allons démonter cinq caractères du thème de l'exode rural pour comprendre leur logique interne et apprendre à les écrire correctement, trait par trait.

\courssub{Qu'est-ce qu'un radical ? Rappel et observation}

\begin{definition}[Le radical (部首)]
Le \cle{radical} (ou clé) est la partie d'un caractère chinois qui sert à le classer dans les dictionnaires et qui indique le plus souvent sa \textbf{famille de sens}. L'autre partie du caractère donne généralement une indication de \textbf{prononciation} (on l'appelle la \emph{phonétique}). Un caractère composé est donc souvent construit ainsi : \textbf{radical (sens) + phonétique (son)}.
\end{definition}

Observez ces trois caractères du glossaire : \zh{流} (couler), \zh{河} (rivière), \zh{海} (mer). Qu'ont-ils en commun ? À gauche, trois petits traits : \zh{氵}. Et tous trois ont un rapport avec l'eau. Ce n'est pas un hasard.

\begin{propriete}[Le radical 氵, « trois gouttes d'eau »]
Le radical \zh{氵} (forme abrégée de \zh{水} \emph{shuǐ}, eau) apparaît à gauche du caractère et indique souvent un lien avec \textbf{l'eau, le liquide ou le mouvement} :
\begin{itemize}
\item \zh{流} \emph{liú} : couler, flux (d'où \zh{流动人口} \emph{liúdòng rénkǒu}, population flottante) ;
\item \zh{河} \emph{hé} : rivière ;
\item \zh{海} \emph{hǎi} : mer ;
\item \zh{洗} \emph{xǐ} : laver ;
\item \zh{汗} \emph{hàn} : sueur.
\end{itemize}
\end{propriete}

\begin{methode}[Deviner le sens d'un caractère inconnu]
Face à un caractère que vous ne connaissez pas, ne paniquez pas : cherchez d'abord son radical.
\begin{enumerate}
\item \textbf{Repérez le radical} : il se trouve le plus souvent à gauche (\zh{氵}, \zh{木}, \zh{土}) ou en haut (\zh{宀}, \zh{艹}) du caractère.
\item \textbf{Déduisez la famille de sens} : \zh{木} → bois, végétal ; \zh{土} → terre, lieu ; \zh{氵} → eau, mouvement.
\item \textbf{Regardez l'autre partie} : elle suggère souvent la prononciation. Exemple : dans \zh{城} \emph{chéng}, la partie droite \zh{成} se prononce aussi \emph{chéng}.
\item \textbf{Vérifiez dans le contexte} : la phrase confirme ou corrige votre hypothèse.
\end{enumerate}
Cette stratégie est précieuse à l'examen : elle permet de comprendre l'essentiel d'un texte même avec des caractères inconnus.
\end{methode}

\courssub{Étude détaillée des cinq caractères du thème}

\textbf{1. 村 \emph{cūn} : village — 7 traits}\par
Décomposition : \zh{木} (bois, arbre ; radical de sens) + \zh{寸} \emph{cùn} (pouce, unité de mesure ; phonétique proche de \emph{cūn}). Le bois évoque le monde rural, les arbres qui entourent les villages. Ordre des traits :
\begin{enumerate}
\item \zh{木} : trait horizontal 一 (de gauche à droite) ;
\item trait vertical 丨 (de haut en bas, il traverse l'horizontal) ;
\item trait tombant à gauche 丿 ;
\item trait tombant à droite 丶 (pression) ;
\item \zh{寸} : trait horizontal 一 ;
\item crochet vertical 亅 ;
\item point 丶.
\end{enumerate}
Règle appliquée : on écrit d'abord la partie gauche (\zh{木}), puis la partie droite (\zh{寸}) ; à l'intérieur de \zh{木}, l'horizontal avant le vertical, puis les deux tombants.

\textbf{2. 城 \emph{chéng} : ville — 9 traits}\par
Décomposition : \zh{土} (terre ; radical) + \zh{成} (accomplir ; donne le son \emph{chéng}). Autrefois, les villes étaient entourées de murs de terre : la terre + l'accomplissement donnent la « ville fortifiée ». Ordre des traits :
\begin{enumerate}
\item \zh{土} : petit horizontal 一, vertical 丨, trait remontant ㇀ (en position de radical à gauche, le grand horizontal du bas de \zh{土} devient un trait remontant) ;
\item \zh{成} : trait horizontal 一 ;
\item trait tombant 丿 ;
\item crochet horizontal-vertical (trait brisé) ;
\item trait oblique avec crochet ㇂ (le trait le plus long, de haut en bas à droite) ;
\item trait tombant interne 丿 ;
\item point final 丶.
\end{enumerate}

\textbf{3. 流 \emph{liú} : couler, flux — 10 traits}\par
Décomposition : \zh{氵} (eau ; radical) + \zh{㐬} (partie phonétique). Lien de sens : l'eau qui coule → le flux, le courant → par extension le \textbf{déplacement de population}, l'exode. C'est exactement le sens de \zh{外流} \emph{wàiliú} (émigrer, partir) dans notre texte. Ordre des traits : les trois gouttes 氵 d'abord (point 丶, point 丶, trait remontant ㇀), puis la partie droite de haut en bas : point 丶, horizontal 一, plié (trait brisé), point 丶, puis les trois traits du bas : tombant 丿, vertical 丨, vertical-courbe 乚.

\textbf{4. 离 \emph{lí} : quitter, s'éloigner de — 10 traits}\par
Radical : \zh{亠} (tête, couvercle). Sens fondamental : « s'éloigner de », d'où le verbe \zh{离开} \emph{líkāi} (quitter) et la structure \zh{离 + lieu + 远/近} (loin de / près de), par exemple : \zh{我家离学校很远。} \emph{Wǒ jiā lí xuéxiào hěn yuǎn.} « Ma maison est loin de l'école. » Ordre des traits : point 丶, horizontal 一 (le radical \zh{亠}), puis la partie centrale (tombant 丿, point 丶, plié vertical, plié horizontal-vertical), et enfin la partie basse (vertical, plié horizontal-crochet, plié, point). Règle appliquée : de haut en bas, l'extérieur avant l'intérieur pour la partie basse.

\textbf{5. 工 \emph{gōng} : travail — 3 traits}\par
Ce caractère est son propre radical (le radical \zh{工} existe dans les dictionnaires). Il représente un outil de charpentier. Ordre des traits : horizontal du haut 一, vertical 丨, horizontal du bas 一. Famille de mots : \zh{工作} \emph{gōngzuò} (travailler, travail), \zh{工厂} \emph{gōngchǎng} (usine), \zh{工资} \emph{gōngzī} (salaire), \zh{工人} \emph{gōngrén} (ouvrier). Tous ces mots sont au cœur du thème de l'exode rural : les jeunes quittent le village pour trouver du \zh{工作} dans une \zh{工厂} et gagner un meilleur \zh{工资}.

\begin{attention}[工 et 土 : deux caractères presque jumeaux]
\zh{工} (travail) et \zh{土} (terre) s'écrivent tous deux en 3 traits : horizontal, vertical, horizontal. La différence est dans les \textbf{longueurs} :
\begin{itemize}
\item \zh{工} : le trait horizontal du \textbf{haut} est court, celui du \textbf{bas} est long ; le vertical ne dépasse pas ;
\item \zh{土} : le trait horizontal du haut est court, celui du bas est \textbf{plus long}, et le vertical \textbf{dépasse} au-dessus du premier horizontal.
\end{itemize}
Astuce : dans \zh{土}, la terre « pousse » vers le haut (comme une plante) ; dans \zh{工}, l'outil reste bien rangé entre ses deux barres. Erreur fréquente des francophones : écrire \zh{城} avec un \zh{工} à gauche — c'est fautif, le radical de \zh{城} est bien \zh{土}.
\end{attention}

\begin{popfigure}
\begin{tikzpicture}[node distance=0.35cm and 0.9cm, every node/.style={font=\large}]
\node[draw=popblue, fill=popblueL, rounded corners, minimum width=1.1cm, minimum height=1.1cm] (cun) {村};
\node[right=of cun] (eq1) {=};
\node[draw=popgreen, fill=popgreenL, rounded corners, right=of eq1, minimum width=1.1cm, minimum height=1.1cm] (mu) {木};
\node[right=0.25cm of mu] (plus1) {+};
\node[draw=poporange, fill=poporangeL, rounded corners, right=0.25cm of plus1, minimum width=1.1cm, minimum height=1.1cm] (cun2) {寸};
\node[below=0.15cm of mu, font=\small, text=popdark] {bois (sens)};
\node[below=0.15cm of cun2, font=\small, text=popdark] {cùn (son)};

\node[draw=popblue, fill=popblueL, rounded corners, minimum width=1.1cm, minimum height=1.1cm, below=1.1cm of cun] (cheng) {城};
\node[right=of cheng] (eq2) {=};
\node[draw=popgreen, fill=popgreenL, rounded corners, right=of eq2, minimum width=1.1cm, minimum height=1.1cm] (tu) {土};
\node[right=0.25cm of tu] (plus2) {+};
\node[draw=poporange, fill=poporangeL, rounded corners, right=0.25cm of plus2, minimum width=1.1cm, minimum height=1.1cm] (cheng2) {成};
\node[below=0.15cm of tu, font=\small, text=popdark] {terre (sens)};
\node[below=0.15cm of cheng2, font=\small, text=popdark] {chéng (son)};

\node[draw=popblue, fill=popblueL, rounded corners, minimum width=1.1cm, minimum height=1.1cm, below=1.1cm of cheng] (liu) {流};
\node[right=of liu] (eq3) {=};
\node[draw=popgreen, fill=popgreenL, rounded corners, right=of eq3, minimum width=1.1cm, minimum height=1.1cm] (shui) {氵};
\node[right=0.25cm of shui] (plus3) {+};
\node[draw=poporange, fill=poporangeL, rounded corners, right=0.25cm of plus3, minimum width=1.1cm, minimum height=1.1cm] (liu2) {㐬};
\node[below=0.15cm of shui, font=\small, text=popdark] {eau (sens)};
\node[below=0.15cm of liu2, font=\small, text=popdark] {phonétique};
\end{tikzpicture}
\legende{Décomposition des caractères 村, 城 et 流 : en vert, le radical (clé de sens) ; en orange, la partie phonétique.}
\end{popfigure}

\courssub{Tableau récapitulatif}

\begin{center}
\begin{tabularx}{\linewidth}{|X|X|X|X|}
\hline
\rowcolor{popblueL} Caractère & Radical & Traits & Sens et famille de mots \\
\hline
村 cūn & 木 (bois) & 7 & village ; 农村 campagne, 村民 villageois \\
\hline
城 chéng & 土 (terre) & 9 & ville ; 城市 ville, 进城 entrer en ville \\
\hline
流 liú & 氵 (eau) & 10 & couler, flux ; 流动 circuler, 外流 émigrer \\
\hline
离 lí & 亠 (tête) & 10 & quitter ; 离开 partir, 离家 quitter la maison \\
\hline
工 gōng & 工 & 3 & travail ; 工作, 工厂, 工资, 工人 \\
\hline
\end{tabularx}
\end{center}

\courssub{Caractères proches : ne pas confondre}

\begin{attention}[村 / 树 et 城 / 诚]
Deux paires de caractères se ressemblent fortement ; seul le radical ou la phonétique change, et avec lui tout le sens :
\begin{itemize}
\item \zh{村} \emph{cūn} (village) : radical \zh{木} à gauche, phonétique \zh{寸} à droite. \zh{树} \emph{shù} (arbre) : même \zh{木}, mais la partie droite est \zh{对} et non \zh{寸}. Comparez : \zh{村里有很多树。} \emph{Cūn li yǒu hěn duō shù.} « Il y a beaucoup d'arbres dans le village. »
\item \zh{城} \emph{chéng} (ville) : radical \zh{土} (terre). \zh{诚} \emph{chéng} (sincère) : radical \zh{讠} (parole) — la sincérité se montre dans les paroles. Même prononciation, sens totalement différents : ce sont des \textbf{homophones}. Comparez : \zh{城市很大。} \emph{Chéngshì hěn dà.} « La ville est grande. » / \zh{他很诚实。} \emph{Tā hěn chéngshí.} « Il est très honnête. »
\end{itemize}
Stratégie : quand deux caractères se ressemblent, identifiez d'abord le radical à gauche — c'est lui qui tranche le sens.
\end{attention}

\begin{exoresolu}[Identifier le radical et deviner le sens]
Énoncé. Voici trois caractères inconnus : \zh{湖}, \zh{林}, \zh{坏}. Pour chacun : (a) identifiez le radical ; (b) proposez une hypothèse de sens ; (c) vérifiez avec la phrase : \zh{城市旁边的湖很美，湖边有树林。}
\tcblower
Solution détaillée.
\begin{enumerate}
\item \zh{湖} : radical \zh{氵} (eau) à gauche → hypothèse : un mot lié à l'eau. Le contexte « 湖很美 » près de la ville confirme : \zh{湖} \emph{hú} = lac.
\item \zh{林} : deux \zh{木} (bois) → hypothèse : beaucoup d'arbres. \zh{树林} \emph{shùlín} = bois, forêt. Confirmé par le contexte « 湖边有树林 » (« il y a un bois au bord du lac »).
\item \zh{坏} : radical \zh{土} (terre) → hypothèse : lié à la terre, au sol. En réalité \zh{坏} \emph{huài} signifie « mauvais, cassé » : le radical donne une piste historique (ce qui retourne à la terre se dégrade), mais le sens moderne s'est élargi. Leçon : le radical aide, mais le contexte et le dictionnaire décident.
\end{enumerate}
Traduction complète de la phrase : \zh{城市旁边的湖很美，湖边有树林。} \emph{Chéngshì pángbiān de hú hěn měi, hú biān yǒu shùlín.} « Le lac à côté de la ville est très beau, et il y a un bois au bord du lac. »
\end{exoresolu}

\begin{savaistu}[Pourquoi 家 (maison) contient-il un cochon ?]
Le caractère \zh{家} \emph{jiā} (maison, famille) se compose de \zh{宀} (toit) et de \zh{豕} (cochon). Dans la Chine ancienne, élever un cochon sous son toit était le signe d'un foyer établi et prospère : la maison, c'était le toit + le cochon ! Au Cameroun aussi, la possession d'animaux (chèvres, porcs, poules) marque traditionnellement la richesse d'une famille rurale. Quand les jeunes quittent le village — \zh{离开家} \emph{líkāi jiā} — c'est donc littéralement ce « toit » familial qu'ils abandonnent, ce qui donne tout son poids au mot \zh{离家} dans notre texte.
\end{savaistu}

\courssub{Exercice de tracé}

\begin{exercice}[Écriture guidée : chaque caractère cinq fois]
Sur votre cahier d'écriture (cases carrées), écrivez chaque caractère \textbf{cinq fois}, en prononçant le pinyin à voix haute à chaque tracé et en respectant l'ordre des traits :
\begin{enumerate}
\item \zh{村} : 木 (一 丨 丿 丶) puis 寸 (一 亅 丶) — 7 traits ;
\item \zh{城} : 土 (一 丨 ㇀) puis 成 — 9 traits ; attention au grand trait oblique ㇂ ;
\item \zh{流} : 氵 (丶 丶 ㇀) puis la partie droite de haut en bas — 10 traits ;
\item \zh{离} : de haut en bas, l'extérieur avant l'intérieur — 10 traits ;
\item \zh{工} : 一 丨 一 — 3 traits, le horizontal du bas plus long que celui du haut.
\end{enumerate}
Contrôle : demandez à votre voisin de vérifier que votre \zh{工} n'est pas devenu un \zh{土}, et que le radical de \zh{城} est bien \zh{土}.
\end{exercice}

\begin{corrige}[Exercice de tracé]
Barème de vérification (auto-correction en binôme) :
\begin{itemize}
\item \zh{村} : 7 traits comptés, partie gauche écrite avant la partie droite, les deux tombants de \zh{木} symétriques ;
\item \zh{城} : 9 traits, le \zh{土} de gauche est étroit (son dernier horizontal devient un trait remontant ㇀), le trait ㇂ de \zh{成} est le plus long du caractère ;
\item \zh{流} : 10 traits, les trois gouttes \zh{氵} alignées verticalement à gauche, la partie droite écrite de haut en bas ;
\item \zh{离} : 10 traits, le \zh{亠} du haut écrit en premier ;
\item \zh{工} : 3 traits, vertical sans dépassement, horizontal inférieur long.
\end{itemize}
Toute inversion d'ordre de traits doit être corrigée immédiatement : une mauvaise habitude d'écriture est très difficile à perdre ensuite.
\end{corrige}

\begin{aretenir}[L'essentiel de la section]
\begin{itemize}
\item Le \cle{radical} donne la famille de sens : \zh{木} bois, \zh{土} terre, \zh{氵} eau/mouvement, \zh{亠} tête ; la partie restante donne souvent le son.
\item Ordre général des traits : de haut en bas, de gauche à droite, horizontal avant vertical, extérieur avant intérieur.
\item \zh{村} 7 traits, \zh{城} 9 traits, \zh{流} 10 traits, \zh{离} 10 traits, \zh{工} 3 traits.
\item Pièges : \zh{工} ≠ \zh{土} (longueur des horizontaux) ; \zh{村} ≠ \zh{树} ; \zh{城} ≠ \zh{诚} (homophones, radicaux différents).
\item Face à un caractère inconnu : repérer le radical, deviner la famille de sens, confirmer par le contexte.
\end{itemize}
\end{aretenir}

\coursec{Compréhension du texte}

Après avoir appris à écrire les caractères clés du thème (\zh{农}, \zh{村}, \zh{城}, \zh{离}, \zh{剩}…), il est temps de revenir au texte \zh{《农村的年轻人》} pour vérifier que nous le comprenons vraiment. Comprendre un texte chinois, ce n'est pas seulement reconnaître les caractères : c'est pouvoir répondre à des questions précises, en phrases complètes, en réutilisant les mots du texte. C'est exactement la compétence demandée à l'examen.

\courssub{Rappel du texte support}

Relisons le texte étudié dans les sections précédentes. Lisez-le une première fois en silence, puis une deuxième fois à voix haute en suivant le pinyin.

\begin{textebox}[农村的年轻人 — Les jeunes de la campagne]
现在很多农村的年轻人离开家乡，去大城市工作。\\
\emph{Xiànzài hěn duō nóngcūn de niánqīngrén líkāi jiāxiāng, qù dà chéngshì gōngzuò.}\\
« Aujourd'hui, beaucoup de jeunes de la campagne quittent leur village natal pour aller travailler dans les grandes villes. »\\[4pt]
他们为什么离开农村？因为农村的工作机会比较少，工资也不高。\\
\emph{Tāmen wèishénme líkāi nóngcūn? Yīnwèi nóngcūn de gōngzuò jīhuì bǐjiào shǎo, gōngzī yě bù gāo.}\\
« Pourquoi quittent-ils la campagne ? Parce qu'à la campagne, les opportunités de travail sont assez rares et les salaires ne sont pas élevés non plus. »\\[4pt]
在城市里，他们可以找到工作，赚更多的钱，还可以学习新的东西。\\
\emph{Zài chéngshì li, tāmen kěyǐ zhǎodào gōngzuò, zhuàn gèng duō de qián, hái kěyǐ xuéxí xīn de dōngxi.}\\
« En ville, ils peuvent trouver du travail, gagner plus d'argent, et aussi apprendre des choses nouvelles. »\\[4pt]
他们离开以后，农村里只剩下老人和孩子。\\
\emph{Tāmen líkāi yǐhòu, nóngcūn li zhǐ shèngxià lǎorén hé háizi.}\\
« Après leur départ, il ne reste au village que les vieillards et les enfants. »\\[4pt]
现在，中国政府正在帮助农村发展：修新路，建新学校，让年轻人可以回家乡工作。\\
\emph{Xiànzài, Zhōngguó zhèngfǔ zhèngzài bāngzhù nóngcūn fāzhǎn : xiū xīn lù, jiàn xīn xuéxiào, ràng niánqīngrén kěyǐ huí jiāxiāng gōngzuò.}\\
« Aujourd'hui, le gouvernement chinois est en train d'aider la campagne à se développer : construire de nouvelles routes, bâtir de nouvelles écoles, pour permettre aux jeunes de retourner travailler au village. »
\end{textebox}

\begin{center}
\begin{tabularx}{\linewidth}{|X|X|X|X|}
\hline
\rowcolor{popblueL} Caractères & Pinyin & Nature & Traduction \\
\hline
离开 & líkāi & verbe & quitter, partir de \\
\hline
家乡 & jiāxiāng & nom & village natal, pays natal \\
\hline
机会 & jīhuì & nom & opportunité, occasion \\
\hline
工资 & gōngzī & nom & salaire \\
\hline
赚 & zhuàn & verbe & gagner (de l'argent) \\
\hline
剩下 & shèngxià & verbe & rester (il ne reste que) \\
\hline
政府 & zhèngfǔ & nom & gouvernement \\
\hline
发展 & fāzhǎn & verbe / nom & (se) développer ; développement \\
\hline
修 & xiū & verbe & construire, réparer (route, ouvrage) \\
\hline
让 & ràng & verbe & permettre, faire que, laisser \\
\hline
\end{tabularx}
\end{center}

\begin{attention}[Le verbe 修 : construire ou réparer ?]
\zh{修} (xiū) a deux sens proches : « construire » (une route, un pont) et « réparer » (un objet abîmé). Dans \zh{修新路}, il s'agit bien de \textbf{construire} de nouvelles routes, et non de « réparer des routes neuves ». C'est le contexte qui décide : \zh{修路} « construire une route », \zh{修自行车} « réparer un vélo ».
\end{attention}

\courssub{Méthode : comment répondre à une question de compréhension}

En chinois comme en français, une bonne réponse à une question de compréhension est une \cle{phrase complète}, pas un mot isolé. La technique la plus sûre consiste à \textbf{reprendre les mots de la question} comme squelette de la réponse, puis à \textbf{compléter avec l'information du texte}.

\begin{methode}[Répondre à une question de compréhension en 4 étapes]
\begin{enumerate}
\item \textbf{Repérer le mot interrogatif} : \zh{为什么} (pourquoi), \zh{什么} (quoi), \zh{谁} (qui), \zh{哪儿} (où), \zh{怎么} (comment).
\item \textbf{Repérer dans le texte la phrase qui contient la réponse} : cherchez les mots-clés de la question (\zh{离开}, \zh{城市}, \zh{政府}…).
\item \textbf{Recopier le squelette de la question} en remplaçant le mot interrogatif par l'information trouvée. Exemple : \zh{年轻人为什么离开农村？} → \zh{年轻人离开农村，因为…}
\item \textbf{Vérifier la phrase} : sujet + verbe + complément, ordre des mots correct, ponctuation chinoise 。
\end{enumerate}
\end{methode}

\begin{exemplebox}[La méthode en action]
Question : \zh{年轻人为什么离开农村？}\\
\emph{Niánqīngrén wèishénme líkāi nóngcūn?}\\
Étape 1 : mot interrogatif = \zh{为什么} (pourquoi).\\
Étape 2 : phrase du texte = \zh{因为农村的工作机会比较少，工资也不高。}\\
Étape 3 : réponse complète = \zh{年轻人离开农村，因为农村的工作机会比较少，工资也不高。}\\
« Les jeunes quittent la campagne parce que les opportunités de travail y sont assez rares et les salaires peu élevés. »
\end{exemplebox}

\courssub{Deux points de grammaire utiles pour répondre}

Avant de répondre aux questions, observons deux petits mots du texte qui posent souvent problème aux francophones : \zh{以后} et \zh{只}.

\begin{propriete}[La place de 以后 (yǐhòu) : « après »]
Contrairement au français où « après » se place \emph{avant} (« après leur départ »), le chinois \zh{以后} se place \textbf{après le verbe ou la proposition} qu'il suit :\\[4pt]
\centering
Proposition / Verbe + \zh{以后} , phrase principale\\[4pt]
\raggedright
Exemples :
\begin{itemize}
\item \zh{他们离开以后，农村里只剩下老人和孩子。} \emph{Tāmen líkāi yǐhòu, nóngcūn li zhǐ shèngxià lǎorén hé háizi.} « Après leur départ, il ne reste au village que les vieillards et les enfants. »
\item \zh{下课以后，我们去踢足球。} \emph{Xiàkè yǐhòu, wǒmen qù tī zúqiú.} « Après les cours, nous allons jouer au football. »
\item \zh{吃饭以后，他学习汉语。} \emph{Chīfàn yǐhòu, tā xuéxí Hànyǔ.} « Après le repas, il étudie le chinois. »
\end{itemize}
Remarque : \zh{以后} peut aussi signifier « plus tard, à l'avenir » (\zh{以后我想去中国。} « Plus tard, je voudrais aller en Chine. »). Son contraire est \zh{以前} (yǐqián, avant), qui se place au même endroit : \zh{吃饭以前} « avant le repas ».
\end{propriete}

\begin{attention}[La place de 只 (zhǐ) : « seulement »]
\zh{只} est un \textbf{adverbe} : il se place \textbf{avant le verbe}, jamais entre le verbe et son complément.\\[4pt]
Correct : \zh{农村里只剩下老人和孩子。} \emph{Nóngcūn li zhǐ shèngxià lǎorén hé háizi.} « Au village, il ne reste que les vieillards et les enfants. »\\[4pt]
Incorrect : \zh{剩下只老人和孩子} (le français « il ne reste \emph{que}… » pousse à mettre \zh{只} après le verbe — piège classique !)\\[4pt]
Autres exemples :
\begin{itemize}
\item \zh{我只有一块钱。} \emph{Wǒ zhǐ yǒu yí kuài qián.} « Je n'ai qu'un yuan. »
\item \zh{他只会说一点儿汉语。} \emph{Tā zhǐ huì shuō yìdiǎnr Hànyǔ.} « Il ne sait parler qu'un peu chinois. »
\end{itemize}
\end{attention}

\courssub{Questions de compréhension (avec réponses attendues)}

Voici les quatre questions principales du cours. Pour chacune, nous donnons la réponse attendue en phrase complète, avec pinyin et traduction.

\textbf{Q1.} \zh{年轻人为什么离开农村？} \emph{Niánqīngrén wèishénme líkāi nóngcūn?} — Pourquoi les jeunes quittent-ils la campagne ?

Réponse attendue : \zh{年轻人离开农村，因为农村的工作机会比较少，工资也不高。}\\
\emph{Niánqīngrén líkāi nóngcūn, yīnwèi nóngcūn de gōngzuò jīhuì bǐjiào shǎo, gōngzī yě bù gāo.}\\
« Les jeunes quittent la campagne parce que les opportunités de travail y sont assez rares et que les salaires n'y sont pas élevés non plus. »

\textbf{Q2.} \zh{他们在城市里可以找到什么？} \emph{Tāmen zài chéngshì li kěyǐ zhǎodào shénme?} — Que peuvent-ils trouver en ville ?

Réponse attendue : \zh{在城市里，他们可以找到工作，赚更多的钱，还可以学习新的东西。}\\
\emph{Zài chéngshì li, tāmen kěyǐ zhǎodào gōngzuò, zhuàn gèng duō de qián, hái kěyǐ xuéxí xīn de dōngxi.}\\
« En ville, ils peuvent trouver du travail, gagner plus d'argent, et aussi apprendre des choses nouvelles. »

\textbf{Q3.} \zh{他们离开以后，农村里剩下谁？} \emph{Tāmen líkāi yǐhòu, nóngcūn li shèngxià shéi?} — Qui reste au village après leur départ ?

Réponse attendue : \zh{他们离开以后，农村里只剩下老人和孩子。}\\
\emph{Tāmen líkāi yǐhòu, nóngcūn li zhǐ shèngxià lǎorén hé háizi.}\\
« Après leur départ, il ne reste au village que les vieillards et les enfants. »

\textbf{Q4.} \zh{中国政府正在做什么？} \emph{Zhōngguó zhèngfǔ zhèngzài zuò shénme?} — Que fait le gouvernement chinois ?

Réponse attendue : \zh{中国政府正在帮助农村发展：修新路，建新学校，让年轻人可以回家乡工作。}\\
\emph{Zhōngguó zhèngfǔ zhèngzài bāngzhù nóngcūn fāzhǎn : xiū xīn lù, jiàn xīn xuéxiào, ràng niánqīngrén kěyǐ huí jiāxiāng gōngzuò.}\\
« Le gouvernement chinois est en train d'aider la campagne à se développer : il construit de nouvelles routes, bâtit de nouvelles écoles, et permet aux jeunes de retourner travailler au village. »

\begin{aretenir}[L'action en cours : 正在 + verbe]
Dans la Q4, \zh{正在} (zhèngzài) marque une \textbf{action en train de se faire}, comme « être en train de » en français : \zh{政府正在帮助农村发展} « le gouvernement est en train d'aider la campagne ». Structure : Sujet + \zh{正在} + Verbe (+ \zh{呢} en fin de phrase, optionnel). On trouve aussi la forme courte \zh{在} : \zh{他在学习。} « Il est en train d'étudier. »
\end{aretenir}

\courssub{Schéma de compréhension : la logique du texte}

Pour bien mémoriser le texte, visualisons sa structure logique : les \textbf{causes} poussent les jeunes au \textbf{départ}, ce qui entraîne des \textbf{conséquences}, auxquelles le gouvernement répond par des \textbf{solutions}.

\begin{popfigure}
\begin{tikzpicture}[
  node distance=0.55cm and 1.1cm,
  box/.style={draw=popblue, fill=popblueL, rounded corners=3pt, align=center, inner sep=5pt, font=\small},
  cause/.style={draw=poporange, fill=poporangeL, rounded corners=3pt, align=center, inner sep=5pt, font=\small},
  cons/.style={draw=poppink, fill=poppinkL, rounded corners=3pt, align=center, inner sep=5pt, font=\small},
  sol/.style={draw=popgreen, fill=popgreenL, rounded corners=3pt, align=center, inner sep=5pt, font=\small},
  fleche/.style={-{Stealth[length=3mm]}, thick, popdark}
]
\node[cause, align=center] (c1){工作机会少\\ \emph{gōngzuò jīhuì shǎo}\\ peu d'emplois};
\node[cause, below=of c1, align=center] (c2){工资不高\\ \emph{gōngzī bù gāo}\\ salaires bas};
\node[box, right=of $(c1)!0.5!(c2)$, xshift=1.6cm, align=center] (dep){离开农村\\ \emph{líkāi nóngcūn}\\ départ du village};
\node[cons, right=of dep, align=center] (cons){只剩下老人和孩子\\ \emph{zhǐ shèngxià lǎorén hé háizi}\\ il ne reste que vieillards\\ et enfants};
\node[sol, below=of cons, align=center] (sol){政府帮助农村发展\\ \emph{zhèngfǔ bāngzhù nóngcūn fāzhǎn}\\ le gouvernement aide\\ la campagne};
\draw[fleche] (c1.east) -- (dep.west |- c1.east);
\draw[fleche] (c2.east) -- (dep.west |- c2.east);
\draw[fleche] (dep) -- (cons) node[midway, above, font=\scriptsize, popdark]{conséquence};
\draw[fleche] (cons) -- (sol) node[midway, right, font=\scriptsize, popdark]{solution};
\draw[fleche, dashed, popgreen] (sol.west) .. controls +(180:1.2) and +(270:1.2) .. (dep.south) node[midway, below, font=\scriptsize, popgreen]{retour possible};
\end{tikzpicture}
\legende{Organigramme de compréhension : causes → départ → conséquences → solutions.}
\end{popfigure}

\courssub{Vrai ou faux : vérifier la compréhension globale}

\begin{exoresolu}[Vrai (对) ou faux (错) ?]
Lisez chaque phrase et dites si elle est vraie (\zh{对} duì) ou fausse (\zh{错} cuò) d'après le texte. Corrigez les phrases fausses.
\begin{enumerate}
\item \zh{农村的年轻人离开家乡，去大城市工作。}
\item \zh{农村的工作机会很多，工资很高。}
\item \zh{在城市里，年轻人可以赚更多的钱。}
\item \zh{年轻人离开以后，农村里只剩下老人。}
\item \zh{中国政府正在帮助城市发展。}
\end{enumerate}
\tcblower
\textbf{Solutions détaillées :}
\begin{enumerate}
\item \textbf{对 (Vrai).} C'est la première phrase du texte : \zh{很多农村的年轻人离开家乡，去大城市工作。}
\item \textbf{错 (Faux).} Le texte dit le contraire : \zh{农村的工作机会比较少，工资也不高。} « Les opportunités sont assez rares et les salaires peu élevés. »
\item \textbf{对 (Vrai).} Le texte dit : \zh{他们可以找到工作，赚更多的钱。}
\item \textbf{错 (Faux).} Phrase incomplète : le texte dit \zh{只剩下老人和孩子} « il ne reste que les vieillards \emph{et les enfants} ». Il manque \zh{和孩子}.
\item \textbf{错 (Faux).} Attention à la substitution : le gouvernement aide la \textbf{campagne} (\zh{农村}), pas la ville (\zh{城市}) : \zh{中国政府正在帮助农村发展。}
\end{enumerate}
\end{exoresolu}

\begin{attention}[Pièges des questions vrai/faux]
Les phrases fausses jouent souvent sur trois pièges : (1) l'\textbf{inversion} (\zh{机会多} au lieu de \zh{机会少}) ; (2) l'\textbf{omission} (oublier \zh{和孩子}) ; (3) la \textbf{substitution} d'un mot proche (\zh{城市} à la place de \zh{农村}). Relisez toujours la phrase du texte mot à mot avant de répondre.
\end{attention}

\courssub{Question ouverte : et au Cameroun ?}

\begin{experience}[Discussion guidée : l'exode rural au Cameroun]
Question : \zh{喀麦隆也有农村人口外流吗？} \emph{Kāmàilóng yě yǒu nóngcūn rénkǒu wàiliú ma?} — L'exode rural existe-t-il aussi au Cameroun ?\\[4pt]
Travail par groupes de trois ou quatre. Préparez deux ou trois phrases en chinois avec le vocabulaire de la leçon, puis présentez-les à la classe. Modèles de phrases :
\begin{itemize}
\item \zh{喀麦隆的年轻人也离开农村，去雅温得和杜阿拉工作。} \emph{Kāmàilóng de niánqīngrén yě líkāi nóngcūn, qù Yǎwēndé hé Dùālā gōngzuò.} « Les jeunes du Cameroun quittent aussi la campagne pour aller travailler à Yaoundé et à Douala. »
\item \zh{他们离开以后，农村里只剩下老人和孩子。} \emph{Tāmen líkāi yǐhòu, nóngcūn li zhǐ shèngxià lǎorén hé háizi.} « Après leur départ, il ne reste au village que les vieillards et les enfants. »
\item \zh{政府应该帮助农村发展。} \emph{Zhèngfǔ yīnggāi bāngzhù nóngcūn fāzhǎn.} « Le gouvernement devrait aider la campagne à se développer. »
\end{itemize}
\end{experience}

\begin{savaistu}[Cameroun et Chine : un même phénomène]
L'exode rural est un phénomène mondial. En Chine, des centaines de millions de jeunes ont quitté les campagnes pour les villes depuis les années 1980, ce qui a transformé le pays. Au Cameroun, les jeunes des villages partent surtout vers Yaoundé et Douala à la recherche d'emplois et d'études. Dans les deux pays, les gouvernements cherchent à rendre la vie rurale plus attractive : routes, écoles, électricité, internet. Comparer les deux situations est un excellent sujet d'exposé pour l'oral !
\end{savaistu}

\courssub{Tableau récapitulatif : les mots interrogatifs et leurs réponses}

\begin{center}
\begin{tabularx}{\linewidth}{|X|X|X|}
\hline
\rowcolor{popblueL} Mot interrogatif & Question type & Élément de réponse \\
\hline
为什么 wèishénme (pourquoi) & 年轻人为什么离开农村？ & 因为 + cause : 因为工作机会少 \\
\hline
什么 shénme (quoi) & 他们可以找到什么？ & 找到工作，赚更多的钱 \\
\hline
谁 shéi (qui) & 农村里剩下谁？ & 只剩下老人和孩子 \\
\hline
哪儿 nǎr (où) & 年轻人去哪儿工作？ & 去大城市工作 \\
\hline
怎么 zěnme (comment) & 政府怎么帮助农村？ & 修新路，建新学校 \\
\hline
\end{tabularx}
\end{center}

\begin{exercice}[À vous de répondre]
Répondez en chinois, en phrases complètes, d'après le texte.
\begin{enumerate}
\item \zh{年轻人离开农村以后，去哪儿工作？}
\item \zh{政府怎么帮助农村发展？}
\item \zh{你的家乡有农村人口外流吗？为什么？}
\end{enumerate}
\end{exercice}

\begin{corrige}[Exercice « À vous de répondre »]
\begin{enumerate}
\item \zh{年轻人离开农村以后，去大城市工作。} \emph{Niánqīngrén líkāi nóngcūn yǐhòu, qù dà chéngshì gōngzuò.} « Après avoir quitté la campagne, les jeunes vont travailler dans les grandes villes. » (Attention : \zh{以后} se place après \zh{离开农村}.)
\item \zh{政府修新路，建新学校，让年轻人可以回家乡工作。} \emph{Zhèngfǔ xiū xīn lù, jiàn xīn xuéxiào, ràng niánqīngrén kěyǐ huí jiāxiāng gōngzuò.} « Le gouvernement construit de nouvelles routes et de nouvelles écoles pour permettre aux jeunes de retourner travailler au village. »
\item Réponse libre. Modèle : \zh{有。因为农村的工作机会少，工资不高，所以年轻人去雅温得工作。} \emph{Yǒu. Yīnwèi nóngcūn de gōngzuò jīhuì shǎo, gōngzī bù gāo, suǒyǐ niánqīngrén qù Yǎwēndé gōngzuò.} « Oui. Parce qu'il y a peu d'emplois et des salaires bas à la campagne, les jeunes partent travailler à Yaoundé. »
\end{enumerate}
\end{corrige}

\begin{aretenir}[L'essentiel de la section]
\begin{itemize}
\item Pour répondre à une question de compréhension : reprendre les mots de la question + compléter avec l'information du texte, en \textbf{phrase complète}.
\item \zh{以后} (après) se place \textbf{après} le verbe ou la proposition : \zh{他们离开以后} « après leur départ » ; son contraire \zh{以前} (avant) occupe la même place.
\item \zh{只} (seulement) se place \textbf{avant} le verbe : \zh{只剩下老人和孩子} « il ne reste que les vieillards et les enfants ».
\item \zh{正在} + verbe marque l'action en cours : \zh{政府正在帮助农村发展}.
\item La logique du texte : causes (\zh{工作机会少，工资不高}) → départ → conséquence (\zh{只剩下老人和孩子}) → solution (\zh{政府帮助农村发展}).
\end{itemize}
\end{aretenir}

\coursec{Collocations et exemples}

Dans la section précédente, nous avons vérifié notre compréhension globale du texte \zh{农村的年轻人} grâce aux questions. Mais comprendre un texte ne suffit pas : pour parler et écrire en chinois, il faut savoir \emph{combiner} les mots entre eux. En chinois comme en français, certains mots « vont naturellement ensemble » : on dit en français « chercher du travail » et non « regarder du travail ». Ces associations habituelles s'appellent des \cle{collocations}. C'est exactement ce que les correcteurs du Bac attendent : des combinaisons naturelles, et non des traductions mot à mot. Reprenons donc les six collocations-clés du texte sur l'exode rural, observons-les, puis entraînons-nous à les employer.

\courssub{Observer d'abord : les collocations dans leur contexte}

\begin{textebox}[Mini-texte de réinvestissement : \zh{进城打工的年轻人}]
李明是农村人。因为农村的工作太少，工资不高，所以他离开了家乡。\\
Lǐ Míng shì nóngcūn rén. Yīnwèi nóngcūn de gōngzuò tài shǎo, gōngzī bù gāo, suǒyǐ tā líkāile jiāxiāng.\\
Li Ming est un homme de la campagne. Parce qu'il y avait trop peu de travail à la campagne et que les salaires n'étaient pas élevés, il a quitté son pays natal.\\[4pt]
他和很多年轻人一样进城打工，在大城市里找工作。\\
Tā hé hěn duō niánqīngrén yīyàng jìn chéng dǎgōng, zài dà chéngshì lǐ zhǎo gōngzuò.\\
Comme beaucoup de jeunes, il est entré en ville pour travailler et cherche du travail dans une grande ville.\\[4pt]
农村人口外流是一个社会问题。政府正在想办法，帮助农村发展，让年轻人可以回家乡工作。\\
Nóngcūn rénkǒu wàiliú shì yī ge shèhuì wèntí. Zhèngfǔ zhèngzài xiǎng bànfǎ, bāngzhù nóngcūn fāzhǎn, ràng niánqīngrén kěyǐ huí jiāxiāng gōngzuò.\\
L'exode rural est un problème de société. Le gouvernement est en train de chercher des solutions pour aider la campagne à se développer, afin que les jeunes puissent retourner travailler au pays.
\end{textebox}

Ce court texte réutilise les six collocations de la leçon dans un contexte cohérent. Lisons-les une à une.

\courssub{Les six collocations du thème}

\begin{definition}[Qu'est-ce qu'une collocation ?]
Une \cle{collocation} est une association habituelle et stable de deux mots (ou plus) : verbe + nom (\zh{找工作}), verbe + complément (\zh{进城打工}), nom + nom (\zh{农村人口外流}). En chinois, ces combinaisons sont très codifiées : changer le verbe rend la phrase bizarre ou fausse, même si la traduction française semble logique.
\end{definition}

\begin{center}
\begin{tabularx}{\linewidth}{|X|X|X|X|}
\hline
\rowcolor{popblueL} Collocation & Pinyin & Structure & Traduction \\
\hline
离开家乡 & líkāi jiāxiāng & verbe + nom & quitter son pays natal \\
\hline
进城打工 & jìn chéng dǎgōng & verbe + nom + verbe & entrer en ville pour travailler \\
\hline
找工作 & zhǎo gōngzuò & verbe + nom & chercher du travail \\
\hline
农村人口外流 & nóngcūn rénkǒu wàiliú & nom + nom + verbe & exode rural \\
\hline
想办法 & xiǎng bànfǎ & verbe + nom & chercher une solution \\
\hline
帮助农村发展 & bāngzhù nóngcūn fāzhǎn & verbe + nom + verbe & aider la campagne à se développer \\
\hline
\end{tabularx}
\end{center}

\textbf{Collocation 1 : \zh{离开家乡} — quitter son pays natal}\par
Le verbe \zh{离开} (\emph{líkāi}, « quitter ») se construit \emph{directement} avec le lieu quitté, sans préposition : \zh{离开} + lieu. C'est une différence majeure avec le français, qui exige une préposition (« quitter \emph{son pays} », « partir \emph{de} son village »).

\begin{exemplebox}[\zh{离开} + lieu]
\zh{很多年轻人离开家乡。} \emph{Hěn duō niánqīngrén líkāi jiāxiāng.} « Beaucoup de jeunes quittent leur pays natal. »\\
\zh{他三年前离开了农村。} \emph{Tā sān nián qián líkāile nóngcūn.} « Il a quitté la campagne il y a trois ans. »\\
\zh{她不想离开雅温得。} \emph{Tā bù xiǎng líkāi Yǎwēndé.} « Elle ne veut pas quitter Yaoundé. »
\end{exemplebox}

\textbf{Collocation 2 : \zh{进城打工} — entrer en ville pour travailler}\par
Deux verbes se suivent : \zh{进城} (\emph{jìn chéng}, « entrer en ville ») indique le déplacement, \zh{打工} (\emph{dǎgōng}, « travailler, souvent comme ouvrier ou employé migrant ») indique le but. Le deuxième verbe exprime l'objectif du premier : « entrer en ville \emph{pour} travailler ».

\begin{exemplebox}[\zh{进城打工}]
\zh{他们进城打工。} \emph{Tāmen jìn chéng dǎgōng.} « Ils entrent en ville pour travailler. »\\
\zh{我哥哥去年进城打工了。} \emph{Wǒ gēge qùnián jìn chéng dǎgōng le.} « Mon grand frère est parti travailler en ville l'année dernière. »
\end{exemplebox}

\textbf{Collocation 3 : \zh{找工作} — chercher du travail}\par
C'est la collocation la plus fréquente du thème : \zh{找} (\emph{zhǎo}, « chercher ») + \zh{工作} (\emph{gōngzuò}, « travail »).

\begin{exemplebox}[\zh{找工作}]
\zh{年轻人在城市里找工作。} \emph{Niánqīngrén zài chéngshì lǐ zhǎo gōngzuò.} « Les jeunes cherchent du travail en ville. »\\
\zh{找工作不容易。} \emph{Zhǎo gōngzuò bù róngyì.} « Chercher du travail n'est pas facile. »
\end{exemplebox}

\begin{attention}[Le bon verbe : \zh{找}, pas \zh{看} ni \zh{问}]
En français, « chercher » peut se traduire par plusieurs verbes selon l'idée, mais en chinois, la collocation est fixe : on dit \zh{找工作} avec \zh{找} (chercher). Ne dites jamais \zh{看工作} (regarder le travail) ni \zh{问工作} (demander le travail) : ces phrases sont incorrectes. De même : \zh{找房子} (chercher un logement), \zh{找人} (chercher quelqu'un). Le verbe \zh{找} est le verbe de la recherche concrète.
\end{attention}

\textbf{Collocation 4 : \zh{农村人口外流} — l'exode rural}\par
C'est le terme sociologique de la leçon : \zh{农村} (campagne) + \zh{人口} (population) + \zh{外流} (s'écouler vers l'extérieur : \zh{外} « extérieur » + \zh{流} « couler »). Littéralement : « l'écoulement vers l'extérieur de la population rurale ». Belle image : la population \emph{coule} hors des villages, comme l'eau d'une rivière.

\begin{exemplebox}[\zh{农村人口外流}]
\zh{农村人口外流是一个社会问题。} \emph{Nóngcūn rénkǒu wàiliú shì yī ge shèhuì wèntí.} « L'exode rural est un problème de société. »\\
\zh{农村人口外流越来越严重。} \emph{Nóngcūn rénkǒu wàiliú yuè lái yuè yánzhòng.} « L'exode rural devient de plus en plus grave. »
\end{exemplebox}

\textbf{Collocation 5 : \zh{想办法} — chercher une solution}\par
\zh{想} (\emph{xiǎng}, « penser, réfléchir ») + \zh{办法} (\emph{bànfǎ}, « moyen, solution ») : littéralement « réfléchir à un moyen ». Très utile à l'écrit comme à l'oral.

\begin{exemplebox}[\zh{想办法}]
\zh{政府正在想办法。} \emph{Zhèngfǔ zhèngzài xiǎng bànfǎ.} « Le gouvernement est en train de chercher une solution. »\\
\zh{我们应该想办法帮助农村。} \emph{Wǒmen yīnggāi xiǎng bànfǎ bāngzhù nóngcūn.} « Nous devrions chercher des moyens d'aider la campagne. »
\end{exemplebox}

\textbf{Collocation 6 : \zh{帮助农村发展} — aider la campagne à se développer}\par
\zh{帮助} (\emph{bāngzhù}, « aider ») + nom + verbe : la personne aidée est suivie directement de l'action, sans préposition : \zh{帮助农村发展} = « aider la campagne (à) se développer ». Structure : \textbf{帮助 + quelqu'un/quelque chose + verbe}.

\begin{exemplebox}[\zh{帮助} + nom + verbe]
\zh{政府帮助农村发展。} \emph{Zhèngfǔ bāngzhù nóngcūn fāzhǎn.} « Le gouvernement aide la campagne à se développer. »\\
\zh{老师帮助学生学习汉字。} \emph{Lǎoshī bāngzhù xuéshēng xuéxí Hànzì.} « Le professeur aide les élèves à apprendre les caractères. »
\end{exemplebox}

\courssub{La corrélation \zh{因为...所以...}}

Dans le mini-texte, nous avons rencontré une structure essentielle pour exprimer la cause : \zh{因为...所以...}. Observons l'exemple de référence.

\begin{exemplebox}[Exemple traduit]
\zh{因为工资不高，所以他离开了农村。}\\
\emph{Yīnwèi gōngzī bù gāo, suǒyǐ tā líkāile nóngcūn.}\\
« Parce que le salaire n'était pas élevé, il a quitté la campagne. »
\end{exemplebox}

\begin{propriete}[La corrélation \zh{因为...所以...}]
Structure : \textbf{因为 + cause，所以 + conséquence}.\\
\zh{因为} (\emph{yīnwèi}) introduit la cause, \zh{所以} (\emph{suǒyǐ}) introduit la conséquence. Contrairement au français, où « parce que » et « donc » ne s'emploient pas ensemble dans la même phrase, le chinois utilise volontiers les deux à la fois : la corrélation \emph{encadre} la cause et la conséquence.\\
\textbf{Exception / souplesse :} on peut n'utiliser qu'un seul des deux mots — \zh{因为工资不高，他离开了农村。} ou \zh{工资不高，所以他离开了农村。} — et la phrase reste parfaitement correcte. Au Bac, les deux formes sont acceptées ; en revanche, ne traduisez jamais « parce que... donc... » par deux mots français dans votre version : en français, il faut choisir.
\end{propriete}

\begin{center}
\begin{tabularx}{\linewidth}{|X|X|X|}
\hline
\rowcolor{popblueL} Structure & Exemple (caractères + pinyin) & Traduction \\
\hline
因为 + cause，所以 + conséquence & \zh{因为工资不高，所以他离开了家乡。} \emph{Yīnwèi gōngzī bù gāo, suǒyǐ tā líkāile jiāxiāng.} & Parce que le salaire n'était pas élevé, il a quitté son pays. \\
\hline
因为 seul & \zh{因为农村工作太少，年轻人进城打工。} \emph{Yīnwèi nóngcūn gōngzuò tài shǎo, niánqīngrén jìn chéng dǎgōng.} & Comme il y a trop peu de travail à la campagne, les jeunes partent travailler en ville. \\
\hline
所以 seul & \zh{城市的机会多，所以人口外流。} \emph{Chéngshì de jīhuì duō, suǒyǐ rénkǒu wàiliú.} & Les villes offrent plus d'opportunités, donc la population s'exode. \\
\hline
\end{tabularx}
\end{center}

\courssub{Carte mentale des collocations}

\begin{popfigure}
\begin{center}\tcbox[colback=popblue,colframe=popblue,coltext=white,arc=9pt,boxsep=4pt,left=6pt,right=6pt]{\bfseries\small 农村 nóngcūn la campagne}\end{center}
\vspace{-2pt}
\begin{tcbraster}[raster columns=2,raster equal height=rows,raster column skip=6pt,raster row skip=6pt,enhanced,blankest]
\begin{tcolorbox}[enhanced,colback=white,colframe=popblue,boxrule=0.9pt,arc=3pt,title={离开家乡 líkāi jiāxiāng},fonttitle=\bfseries\scriptsize,coltitle=white,colbacktitle=popblue,left=4pt,right=4pt,top=2pt,bottom=2pt,boxsep=1pt,toptitle=1.5pt,bottomtitle=1.5pt,halign=left]
\begin{itemize}[nosep,leftmargin=*,itemsep=1pt,font=\scriptsize]
\item 离开农村 quitter la campagne
\item 离开城市 quitter la ville
\end{itemize}
\end{tcolorbox}
\begin{tcolorbox}[enhanced,colback=white,colframe=poporange,boxrule=0.9pt,arc=3pt,title={进城打工 jìn chéng dǎgōng},fonttitle=\bfseries\scriptsize,coltitle=white,colbacktitle=poporange,left=4pt,right=4pt,top=2pt,bottom=2pt,boxsep=1pt,toptitle=1.5pt,bottomtitle=1.5pt,halign=left]
\begin{itemize}[nosep,leftmargin=*,itemsep=1pt,font=\scriptsize]
\item 找工作 zhǎo gōngzuò
\end{itemize}
\end{tcolorbox}
\begin{tcolorbox}[enhanced,colback=white,colframe=popgreen,boxrule=0.9pt,arc=3pt,title={人口外流 rénkǒu wàiliú},fonttitle=\bfseries\scriptsize,coltitle=white,colbacktitle=popgreen,left=4pt,right=4pt,top=2pt,bottom=2pt,boxsep=1pt,toptitle=1.5pt,bottomtitle=1.5pt,halign=left]
\begin{itemize}[nosep,leftmargin=*,itemsep=1pt,font=\scriptsize]
\item 社会问题 shèhuì wèntí
\end{itemize}
\end{tcolorbox}
\begin{tcolorbox}[enhanced,colback=white,colframe=poppurple,boxrule=0.9pt,arc=3pt,title={帮助发展 bāngzhù fāzhǎn},fonttitle=\bfseries\scriptsize,coltitle=white,colbacktitle=poppurple,left=4pt,right=4pt,top=2pt,bottom=2pt,boxsep=1pt,toptitle=1.5pt,bottomtitle=1.5pt,halign=left]
\begin{itemize}[nosep,leftmargin=*,itemsep=1pt,font=\scriptsize]
\item 想办法 xiǎng bànfǎ
\end{itemize}
\end{tcolorbox}
\end{tcbraster}

\legende{Carte mentale des collocations du thème « exode rural » autour du nom \zh{农村}.}
\end{popfigure}

\courssub{Exercice résolu}

\begin{exoresolu}[Phrases à trous avec les collocations]
Énoncé. Complétez avec la collocation qui convient : \zh{离开家乡} / \zh{进城打工} / \zh{找工作} / \zh{想办法}.\\
1. 很多农村的年轻人 \_\_\_\_\_\_\_\_，因为家乡没有机会。\\
2. 他大学毕业以后，在城市里 \_\_\_\_\_\_\_\_。\\
3. 这个问题很难，我们一起 \_\_\_\_\_\_\_\_ 吧！\\
4. 我的叔叔去年 \_\_\_\_\_\_\_\_，现在在杜阿拉工作。
\tcblower
Solution détaillée.\\
1. \zh{离开家乡} : la cause « pas d'opportunités au pays » appelle l'idée de quitter son pays natal. Phrase complète : \zh{很多农村的年轻人离开家乡，因为家乡没有机会。} « Beaucoup de jeunes ruraux quittent leur pays, car il n'y a pas d'opportunités au pays. »\\
2. \zh{找工作} : après l'université, en ville, on cherche du travail. \zh{他大学毕业以后，在城市里找工作。} « Après l'université, il cherche du travail en ville. »\\
3. \zh{想办法} : devant un problème difficile, on cherche une solution ensemble. \zh{这个问题很难，我们一起想办法吧！} « Ce problème est difficile, cherchons une solution ensemble ! »\\
4. \zh{进城打工} : « maintenant à Douala » confirme le départ en ville pour travailler. \zh{我的叔叔去年进城打工，现在在杜阿拉工作。} « Mon oncle est parti travailler en ville l'année dernière ; il travaille maintenant à Douala. »
\end{exoresolu}

\courssub{Exercices d'entraînement}

\begin{exercice}[À vous : phrases à trous]
Complétez avec : \zh{农村人口外流} / \zh{帮助农村发展} / \zh{离开家乡} / \zh{找工作}.\\
1. \_\_\_\_\_\_\_\_ 是一个严重的社会问题。\\
2. 政府修了很多路，为了 \_\_\_\_\_\_\_\_。\\
3. 她不想 \_\_\_\_\_\_\_\_，她想在家乡工作。\\
4. 在喀麦隆，很多年轻人去雅温得 \_\_\_\_\_\_\_\_。
\end{exercice}

\begin{corrige}[Exercice « À vous : phrases à trous »]
1. \zh{农村人口外流是一个严重的社会问题。} « L'exode rural est un grave problème de société. »\\
2. \zh{政府修了很多路，为了帮助农村发展。} « Le gouvernement a construit beaucoup de routes pour aider la campagne à se développer. »\\
3. \zh{她不想离开家乡，她想在家乡工作。} « Elle ne veut pas quitter son pays, elle veut travailler au pays. »\\
4. \zh{在喀麦隆，很多年轻人去雅温得找工作。} « Au Cameroun, beaucoup de jeunes vont à Yaoundé chercher du travail. »
\end{corrige}

\begin{exercice}[Traduction : du français vers le chinois]
Traduisez en chinois en utilisant les collocations de la leçon.\\
1. Parce que la vie à la campagne est difficile, beaucoup de jeunes quittent leur pays natal.\\
2. Mon frère aîné est entré en ville pour travailler ; il cherche du travail à Douala.\\
3. L'exode rural est un problème de société : le gouvernement cherche une solution.\\
4. Nous devons aider la campagne à se développer pour que les jeunes restent au village.
\end{exercice}

\begin{corrige}[Exercice « Traduction »]
1. \zh{因为农村的生活很难，所以很多年轻人离开家乡。} \emph{Yīnwèi nóngcūn de shēnghuó hěn nán, suǒyǐ hěn duō niánqīngrén líkāi jiāxiāng.}\\
2. \zh{我哥哥进城打工，他在杜阿拉找工作。} \emph{Wǒ gēge jìn chéng dǎgōng, tā zài Dù'ālā zhǎo gōngzuò.}\\
3. \zh{农村人口外流是一个社会问题，政府正在想办法。} \emph{Nóngcūn rénkǒu wàiliú shì yī ge shèhuì wèntí, zhèngfǔ zhèngzài xiǎng bànfǎ.}\\
4. \zh{我们应该帮助农村发展，让年轻人留在村子里。} \emph{Wǒmen yīnggāi bāngzhù nóngcūn fāzhǎn, ràng niánqīngrén liú zài cūnzi lǐ.}
\end{corrige}

\begin{attention}[Pièges fréquents des francophones]
\begin{itemize}
\item Ne mettez pas de préposition après \zh{离开} : on dit \zh{离开家乡}, jamais une construction avec \zh{从} (« de ») entre le verbe et le lieu.
\item « Chercher du travail » = \zh{找工作} ; le verbe est figé, pas de \zh{看} ni de \zh{问}.
\item Attention au ton de \zh{找} (\emph{zhǎo}, 3e ton) : ne le confondez pas avec \zh{早} (\emph{zǎo}, « tôt »).
\item Avec \zh{因为...所以...}, vous pouvez garder les deux mots en chinois, mais en français, traduisez par « parce que » \emph{ou} « donc », jamais les deux.
\end{itemize}
\end{attention}

\begin{savaistu}[\zh{空村} : les villages vidés]
Au Cameroun comme en Chine, l'exode rural vide progressivement les villages : les jeunes partent vers Douala, Yaoundé ou les grandes métropoles chinoises, et il ne reste souvent que les personnes âgées et les enfants. En chinois, on parle même de \zh{空村} (\emph{kōng cūn}), littéralement « villages vides », ou de \zh{空心村} (\emph{kōngxīn cūn}), « villages au cœur vide ». Ce phénomène a poussé le gouvernement chinois à investir massivement dans les campagnes : routes, écoles, internet — exactement le sens de notre collocation \zh{帮助农村发展} !
\end{savaistu}

\begin{aretenir}[L'essentiel de la section]
\begin{itemize}
\item Six collocations à connaître par cœur : \zh{离开家乡}, \zh{进城打工}, \zh{找工作}, \zh{农村人口外流}, \zh{想办法}, \zh{帮助农村发展}.
\item \zh{离开} + lieu, sans préposition ; \zh{帮助} + nom + verbe, sans préposition.
\item Corrélation \zh{因为...所以...} : les deux mots ensemble en chinois, un seul en français.
\item Une collocation est un bloc : mémorisez-la entière, comme un seul mot.
\end{itemize}
\end{aretenir}

\coursec{Production guidée}

Après avoir étudié les collocations et les structures clés du thème « exode rural » dans la section précédente, nous passons maintenant à l'\textbf{expression écrite guidée}. L'objectif est de mobiliser le vocabulaire, la grammaire et les idées du texte support pour produire un paragraphe cohérent, correct et personnel sur la réalité camerounaise. Cette étape prépare directement l'évaluation sommative du module.

\courssub{Objectifs et démarche}

\begin{methode}[Rédiger un paragraphe guidé en 4 étapes]
\begin{enumerate}
\item \textbf{Analyser le plan imposé} : identifier les quatre blocs logiques (phénomène, cause, conséquence, solution) et les mots-clés chinois associés.
\item \textbf{Choisir dans la banque de mots} : sélectionner les verbes, noms, connecteurs et structures vus en classe ; ne pas inventer de vocabulaire hors programme.
\item \textbf{Rédiger une phrase par bloc} : respecter l'ordre des mots chinois (Sujet + Temps + Lieu + Verbe + Complément) et placer les connecteurs logiques (\cle{因为}...\cle{所以}, \cle{但是}, \cle{为了}) en début de proposition.
\item \textbf{Relire et corriger en binôme} : vérifier les tons (pinyin), l'écriture des caractères (ordre des traits), l'accord des classificateurs et la ponctuation chinoise pleine chasse (， 。).
\end{enumerate}
\end{methode}

\courssub{Banque de mots et collocations de la leçon}

Le tableau ci-dessous regroupe le lexique essentiel à réinvestir. Il est autorisé — et même recommandé — de réutiliser les phrases du texte support \zh{农村的年轻人} en remplaçant un ou deux éléments (lieu, chiffre, adjectif) pour les adapter au Cameroun.

\begin{center}
\begin{tabularx}{\linewidth}{|X|X|X|X|}
\hline
\rowcolor{popblueL} \textbf{Caractères} & \textbf{Pinyin} & \textbf{Nature} & \textbf{Traduction} \\
\hline
农村 & nóngcūn & nom & campagne, zone rurale \\
\hline
城市 & chéngshì & nom & ville \\
\hline
村庄 & cūnzhuāng & nom & village \\
\hline
大城市 & dà chéngshì & nom & grande ville \\
\hline
年轻人 & niánqīngrén & nom & jeune(s) \\
\hline
老人 & lǎorén & nom & personne âgée \\
\hline
孩子 & háizi & nom & enfant(s) \\
\hline
小孩 & xiǎohái & nom & enfant(s) (familier) \\
\hline
家乡 & jiāxiāng & nom & village natal, pays natal \\
\hline
离开 & líkāi & verbe & quitter, partir de \\
\hline
去 & qù & verbe & aller (vers) \\
\hline
留 & liú & verbe & rester, rester (quelque part) \\
\hline
找工作 & zhǎo gōngzuò & verbe + objet & chercher du travail / un emploi \\
\hline
种地 & zhòng dì & verbe + objet & cultiver la terre, faire de l'agriculture \\
\hline
工作机会 & gōngzuò jīhuì & nom & opportunité d'emploi \\
\hline
工厂 & gōngchǎng & nom & usine, fabrique \\
\hline
医院 & yīyuàn & nom & hôpital \\
\hline
学校 & xuéxiào & nom & école \\
\hline
路 & lù & nom & route \\
\hline
粮食 & liángshí & nom & vivres, céréales, production alimentaire \\
\hline
生产 & shēngchǎn & verbe / nom & produire / production \\
\hline
少 & shǎo & adj. & peu, insuffisant \\
\hline
多 & duō & adj. & beaucoup, nombreux \\
\hline
只剩下 & zhǐ shèngxià & structure & ne rester que / il ne reste plus que \\
\hline
不利 & búlì & adj. & défavorable, nuisible \\
\hline
更多 & gèng duō & adv. + adj. & plus, davantage \\
\hline
发展 & fāzhǎn & verbe / nom & développer / développement \\
\hline
政府 & zhèngfǔ & nom & gouvernement \\
\hline
应该 & yīnggāi & verbe modal & devoir, devrait \\
\hline
帮助 & bāngzhù & verbe & aider \\
\hline
改善 & gǎishàn & verbe & améliorer \\
\hline
创造 & chuàngzào & verbe & créer \\
\hline
修 & xiū & verbe & construire, réparer (route, maison) \\
\hline
建 & jiàn & verbe & construire, édifier (bâtiment, usine) \\
\hline
修好 & xiū hǎo & verbe + compl. & réparer/construire correctement, achever la construction de \\
\hline
生活条件 & shēnghuó tiáojiàn & nom & conditions de vie \\
\hline
因为……所以…… & yīnwèi……suǒyǐ…… & connecteurs & parce que… donc… \\
\hline
但是 & dànshì & conjonction & mais, cependant \\
\hline
为了 & wèile & préposition & pour, afin de \\
\hline
在 & zài & préposition & à, dans (lieu) \\
\hline
从……到…… & cóng……dào…… & structure & de… à… \\
\hline
喀麦隆 & Kāmàilóng & nom propre & Cameroun \\
\hline
杜阿拉 & Dù'ālā & nom propre & Douala \\
\hline
雅温得 & Yǎwēndé & nom propre & Yaoundé \\
\hline
巴门达 & Bāméndá & nom propre & Bamenda \\
\hline
布埃亚 & Bù'āiyà & nom propre & Buea \\
\hline
西部 & xībù & nom propre & Région de l'Ouest \\
\hline
附近 & fùjìn & nom / loc. & environs, à proximité de \\
\hline
\end{tabularx}
\end{center}

\courssub{Plan de production : 4 étapes fléchées}

Le schéma ci-dessous visualise la progression logique attendue. Chaque flèche correspond à une phrase (ou deux) de votre paragraphe.

\begin{popfigure}
\begin{tikzpicture}[
  node distance=1.8cm and 2.5cm,
  box/.style={rectangle, rounded corners=6pt, draw=popblue, thick, fill=popblueL, text width=3.2cm, align=center, font=\small},
  arrow/.style={->, >=Stealth, thick, popdark, shorten >=2pt, shorten <=2pt}
]
\node[box, align=center] (etape1){\textbf{1. Phénomène}\\\zh{现象}\\很多年轻人离开农村};
\node[box, right=of etape1, align=center] (etape2){\textbf{2. Cause}\\\zh{原因}\\因为工作机会少};
\node[box, right=of etape2, align=center] (etape3){\textbf{3. Conséquence}\\\zh{结果}\\农村里只剩下老人};
\node[box, right=of etape3, align=center] (etape4){\textbf{4. Solution}\\\zh{办法}\\政府应该帮助农村发展};

\draw[arrow] (etape1.east) -- (etape2.west) node[midway, above, font=\tiny, popmuted] {entraîne};
\draw[arrow] (etape2.east) -- (etape3.west) node[midway, above, font=\tiny, popmuted] {entraîne};
\draw[arrow] (etape3.east) -- (etape4.west) node[midway, above, font=\tiny, popmuted] {nécessite};
\end{tikzpicture}
\legende{Schéma du plan de production en 4 étapes : 现象 $\rightarrow$ 原因 $\rightarrow$ 结果 $\rightarrow$ 办法}
\end{popfigure}

\courssub{Modèle de paragraphe rédigé}

Voici un modèle complet de 5 phrases respectant le plan, le vocabulaire de la leçon et les structures grammaticales du niveau HSK 2-3. Étudiez-le attentivement avant de produire le vôtre.

\begin{textebox}[Modèle : L'exode rural au Cameroun]
在喀麦隆，很多年轻人离开农村，去杜阿拉和雅温得找工作。因为农村的工作机会少，所以年轻人不想留在家乡。现在农村里只剩下老人和孩子。这对农村的发展不利。政府应该帮助农村发展，改善生活条件，创造更多工作机会。
\end{textebox}

\textbf{Pinyin :} Zài Kāmàilóng, hěn duō niánqīngrén líkāi nóngcūn, qù Dù'ālā hé Yǎwēndé zhǎo gōngzuò. Yīnwèi nóngcūn de gōngzuò jīhuì shǎo, suǒyǐ niánqīngrén bù xiǎng liú zài jiāxiāng. Xiànzài nóngcūn lǐ zhǐ shèngxià lǎorén hé xiǎohái. Zhè duì nóngcūn de fāzhǎn búlì. Zhèngfǔ yīnggāi bāngzhù nóngcūn fāzhǎn, gǎishàn shēnghuó tiáojiàn, chuàngzào gèng duō gōngzuò jīhuì.

\textbf{Traduction :} Au Cameroun, beaucoup de jeunes quittent la campagne pour aller à Douala et Yaoundé chercher du travail. Parce que les opportunités d'emploi à la campagne sont rares, les jeunes ne veulent pas rester dans leur village natal. Maintenant, il ne reste plus que des personnes âgées et des enfants à la campagne. Cela est défavorable au développement de la campagne. Le gouvernement devrait aider la campagne à se développer, améliorer les conditions de vie et créer plus d'opportunités d'emploi.

\courssub{Analyse linguistique du modèle}

\begin{propriete}[Règle d'ordre : Complément de lieu (\cle{在} + lieu) avant le verbe]
En chinois, le complément de lieu (où se déroule l'action) se place \textbf{avant le verbe}, introduit par la préposition \cle{在} (ou directement par le nom de lieu pour les verbes de déplacement comme \cle{去}, \cle{来}, \cle{回}).
\begin{center}
\begin{tabularx}{\linewidth}{|X|X|X|}
\hline
\rowcolor{poporangeL} \textbf{Structure} & \textbf{Exemple (caractères + pinyin)} & \textbf{Traduction} \\
\hline
Sujet + \cle{在} + Lieu + Verbe + Complément & \zh{我在喀麦隆学习中文。} (Wǒ zài Kāmàilóng xuéxí Zhōngwén.) & J'étudie le chinois \textbf{au Cameroun}. \\
\hline
Sujet + Verbe de déplacement + Lieu + Action & \zh{年轻人去城市找工作。} (Niánqīngrén qù chéngshì zhǎo gōngzuò.) & Les jeunes vont \textbf{en ville} chercher du travail. \\
\hline
Sujet + \cle{从} + Lieu源 + \cle{到} + Lieu但 + Verbe & \zh{他们从农村到城市工作。} (Tāmen cóng nóngcūn dào chéngshì gōngzuò.) & Ils vont \textbf{de la campagne à la ville} travailler. \\
\hline
\end{tabularx}
\end{center}
\emph{Attention :} En français, le lieu se place souvent à la fin (« Je travaille au Cameroun »). En chinois, \cle{在喀麦隆} précède \cle{工作}. Ne dites jamais \zh{*我工作在喀麦隆} (erreur de calque francophone).
\end{propriete}

\begin{attention}[Pièges fréquents des francophones]
\begin{enumerate}
\item \textbf{Ordre des mots (Lieu) :} \zh{*年轻人找工作在城市} $\rightarrow$ \textbf{Correct :} \zh{年轻人在城市找工作} ou \zh{年轻人去城市找工作}.
\item \textbf{Oubli de \cle{所以} dans la corrélation \cle{因为}...\cle{所以}} : \zh{因为工作少，年轻人走} $\rightarrow$ \textbf{Mieux :} \zh{因为工作少，\textbf{所以}年轻人走}. À l'écrit formel, \cle{所以} est quasi obligatoire.
\item \textbf{Classificateur manquant} : \zh{*一个年轻人} (correct) mais \zh{*三年轻人} $\rightarrow$ \zh{三\textbf{个}年轻人} ; \zh{*这问题} $\rightarrow$ \zh{这\textbf{个}问题}.
\item \textbf{Confusion \cle{只} / \cle{只有}} : \cle{只} (adverbe « seulement ») se place \textbf{devant le verbe} : \zh{只剩下老人}. \cle{只有} (verbe « seulement avoir ») introduit un nom : \zh{只有老人}.
\item \textbf{Ponctuation} : Utilisez la virgule chinoise \zh{，} et le point \zh{。}. Pas d'espace entre les caractères. Majuscule seulement pour les noms propres en pinyin.
\end{enumerate}
\end{attention}

\courssub{Exercice résolu : transformation guidée}

\begin{exoresolu}[Transformer une phrase du texte support pour l'adapter au Cameroun]
\textbf{Énoncé :} Partez de la phrase du texte support : \zh{很多年轻人离开农村，去城市打工。} (Beaucoup de jeunes quittent la campagne pour aller travailler en ville.) Adaptez-la au contexte camerounais en précisant les villes (Douala, Yaoundé) et en utilisant \cle{找工作}.

\tcblower
\textbf{Solution détaillée :}
\begin{enumerate}
\item \textbf{Identifier les éléments fixes :} \zh{很多年轻人离开农村} (phénomène) + \zh{去 + Lieu + 找工作} (action but).
\item \textbf{Remplacer « 城市 » par les villes cibles :} \zh{杜阿拉和雅温得}.
\item \textbf{Remplacer « 打工 » (travailler/job temporaire) par « 找工作 » (chercher un emploi), plus formel.}
\item \textbf{Résultat :} \zh{在喀麦隆，很多年轻人离开农村，去杜阿拉和雅温得找工作。}
\item \textbf{Vérification :} \cle{在喀麦隆} (lieu global) en tête de phrase ; \cle{去} + villes + \cle{找工作} (structure verbe de déplacement + lieu + but) ; ponctuation chinoise.
\end{enumerate}
\end{exoresolu}

\courssub{Activité de production en binôme}

\begin{experience}[Production écrite collaborative : « Mon paragraphe sur l'exode rural »]
\textbf{Consignes :}
\begin{enumerate}
\item \textbf{Individuel (5 min) :} Rédigez votre paragraphe de 5--6 phrases sur votre cahier, en suivant strictement le plan en 4 étapes et en puisant dans la banque de mots. Écrivez : caractères + pinyin (au-dessus ou à côté) + traduction française.
\item \textbf{Binôme (8 min) :} Échangez vos cahiers. Votre partenaire vérifie : (a) les 4 étapes sont présentes ; (b) chaque phrase a un verbe ; (c) l'ordre Sujet--Temps--Lieu--Verbe est respecté ; (d) les tons du pinyin sont notés ; (e) les caractères sont lisibles et l'ordre des traits respecté ; (f) la ponctuation chinoise est utilisée. Il annote en vert (correct) ou rouge (à corriger).
\item \textbf{Réécriture (5 min) :} Intégrez les corrections et produisez une version finale propre.
\end{enumerate}
\textbf{Production attendue (exemple d'élève) :}
\zh{在喀麦隆，很多年轻人离开农村，去大城市找工作。因为农村没有工厂，所以工作机会很少。现在农村只剩下老人，土地没人种。这不利于国家发展。政府应该修路、建厂，帮助农村发展。}
\end{experience}

\courssub{Correction collective et critères de réussite}

Au tableau, l'enseignant corrige une production d'élève (anonyme) en projetant le texte. La classe analyse ensemble : respect du plan, justesse lexicale, exactitude grammaticale, écriture des caractères.

\begin{aretenir}[Critères de réussite pour l'évaluation (notés sur 10 pts)]
\begin{itemize}
\item \textbf{Vocabulaire du thème (3 pts) :} Utilisation correcte d'au moins 8 mots de la banque (ex. \zh{农村, 城市, 离开, 找工作, 工作机会, 只剩下, 老人, 发展, 政府, 应该, 帮助, 改善}).
\item \textbf{Exactitude des caractères (2 pts) :} Caractères lisibles, traits dans l'ordre, pas de confusion visuelle (ex. \zh{未} vs \zh{末}, \zh{人} vs \zh{入}, \zh{土} vs \zh{士}).
\item \textbf{Structures grammaticales (3 pts) :} Au moins 3 structures distinctes bien employées : \cle{因为}...\cle{所以}, \cle{在} + lieu + V, \cle{只} + V, \cle{应该} + V, \cle{去} + lieu + V (but).
\item \textbf{Cohérence et plan (2 pts) :} Les 4 étapes logiques sont présentes, ordonnées, reliées par des connecteurs ; paragraphe de 5--6 phrases complètes (sujet + verbe).
\end{itemize}
\end{aretenir}

\courssub{Devoirs et entraînement autonome}

\begin{exercice}[Entraînement à la maison]
\begin{enumerate}
\item Recopiez le modèle de paragraphe \textbf{3 fois} : une fois en caractères seulement, une fois en pinyin avec tons, une fois en français. Mémorisez les 4 phrases-clés.
\item Rédigez une \textbf{variante personnelle} (5 phrases) en changeant : (a) les villes (ex. \zh{布埃亚} Bù'āiyà / Buea, \zh{巴门达} Bāméndá / Bamenda) ; (b) la cause (ex. \zh{因为农村没有大学} « pas d'université à la campagne ») ; (c) la solution (ex. \zh{政府应该在农村建大学}).
\item Préparez une \textbf{présentation orale de 1 minute} (sans notes) pour la prochaine séance : « L'exode rural dans ma région ». Utilisez le même plan.
\end{enumerate}
\end{exercice}

\begin{corrige}[Exercice 1 -- Variantes attendues (exemples)]
\textbf{Variante A (région de l'Ouest) :} \zh{在喀麦隆西部，很多年轻人离开农村，去巴门达和杜阿拉找工作。因为农村没有工厂，所以工作机会少。现在农村只剩下老人种地。这对粮食生产不利。政府应该在农村建工厂，创造工作机会。}

\textbf{Variante B (région du Sud-Ouest) :} \zh{在布埃亚附近，年轻人离开村庄去城市。因为农村的学校少，医院也远。所以只剩下老人和孩子。政府应该修好路，建医院和学校，改善农村生活。}
\end{corrige}

\newpage

\coursec{Activité d'intégration}

\coursec{Leçon 23 — Vocabulaire et compréhension de texte : Exode rural}

\courssub{Texte support : 农村的年轻人}

\begin{textebox}[Texte original : 农村的年轻人]
在中国和喀麦隆，很多年轻人离开农村，去城市生活。\\
\emph{Zài Zhōngguó hé Kāmàilóng, hěn duō niánqīngrén líkāi nóngcūn, qù chéngshì shēnghuó.}\\
« En Chine et au Cameroun, beaucoup de jeunes quittent la campagne pour vivre en ville. »\\[0.5em]
因为城市的工作机会多，工资也高。\\
\emph{Yīnwèi chéngshì de gōngzuò jīhuì duō, gōngzī yě gāo.}\\
« Parce qu'en ville il y a plus d'offres d'emploi et les salaires sont aussi plus élevés. »\\[0.5em]
可是城市的空气不干净，交通不方便，生活很忙。\\
\emph{Kěshì chéngshì de kōngqì bù gānjìng, jiāotōng bù fāngbiàn, shēnghuó hěn máng.}\\
« Mais l'air en ville n'est pas pur, les transports ne sont pas pratiques, la vie est très occupée. »\\[0.5em]
农村的空气干净，环境安静，可是医院少，学校也不好。\\
\emph{Nóngcūn de kōngqì gānjìng, huánjìng ānjìng, kěshì yīyuàn shǎo, xuéxiào yě bù hǎo.}\\
« La campagne a l'air pur, l'environnement est calme, mais il y a peu d'hôpitaux et les écoles ne sont pas bonnes non plus. »\\[0.5em]
我觉得农村应该发展旅游业，让年轻人留在家乡工作。\\
\emph{Wǒ juéde nóngcūn yīnggāi fāzhǎn lǚyóuyè, ràng niánqīngrén liú zài jiāxiāng gōngzuò.}\\
« Je pense que la campagne devrait développer le tourisme, pour que les jeunes restent travailler au pays natal. »
\end{textebox}

\courssub{Glossaire du thème : L'exode rural}

\begin{definition}[Vocabulaire clé du texte et de la leçon]
Le tableau ci-dessous regroupe le lexique essentiel pour comprendre le texte et participer aux activités. La nature grammaticale (nom, verbe, adjectif, auxiliaire, locution) est indiquée pour guider l'emploi en phrase.
\end{definition}

\begin{center}
\begin{tabularx}{\linewidth}{|X|X|X|X|}
\hline
\rowcolor{popblueL} Caractères & Pinyin & Nature & Traduction \\
\hline
农村 & nóngcūn & nom & campagne, milieu rural \\
\hline
城市 & chéngshì & nom & ville \\
\hline
年轻人 & niánqīngrén & nom & jeune, les jeunes \\
\hline
离开 & líkāi & verbe & quitter, partir de \\
\hline
生活 & shēnghuó & nom / verbe & vie ; vivre, mener sa vie \\
\hline
工作机会 & gōngzuò jīhuì & loc. nom. & offre d'emploi, débouché \\
\hline
工资 & gōngzī & nom & salaire, paye \\
\hline
空气 & kōngqì & nom & air \\
\hline
干净 & gānjìng & adjectif & propre, pur \\
\hline
安静 & ānjìng & adjectif & calme, tranquille \\
\hline
方便 & fāngbiàn & adjectif & pratique, commode \\
\hline
交通 & jiāotōng & nom & transports, circulation \\
\hline
环境 & huánjìng & nom & environnement \\
\hline
医院 & yīyuàn & nom & hôpital \\
\hline
学校 & xuéxiào & nom & école \\
\hline
应该 & yīnggāi & auxiliaire & devoir (il faudrait) \\
\hline
发展 & fāzhǎn & verbe & développer \\
\hline
旅游业 & lǚyóuyè & nom & industrie touristique, tourisme \\
\hline
留在 & liú zài & loc. verb. & rester à, demeurer à \\
\hline
家乡 & jiāxiāng & nom & ville natale, pays natal \\
\hline
喀麦隆 & Kāmàilóng & nom propre & Cameroun \\
\hline
中国 & Zhōngguó & nom propre & Chine \\
\hline
\end{tabularx}
\end{center}

\courssub{Radicaux et ordre des traits}

\begin{propriete}[Analyse de caractères clés : radicaux et ordre des traits]
Maîtriser les radicaux (clés) aide à mémoriser le sens et à classer les caractères dans le dictionnaire. L'ordre des traits suit les règles : haut $\to$ bas, gauche $\to$ droite, horizontal avant vertical, extérieur avant intérieur.
\end{propriete}

\begin{exemplebox}[Caractères à analyser]
\begin{enumerate}
\item \textbf{农} (\emph{nóng}) — agriculture, paysan.\\
\cle{Radical} : \zh{曲} (qū, « courbe ») — 6 traits.\\
\cle{Ordre} : 丿 (piě) → 一 (héng) → 丨 (shù) → 丿 (piě) → 乀 (zhé) → 丶 (diǎn). (Le trait horizontal traverse le vertical).
\item \textbf{村} (\emph{cūn}) — village.\\
\cle{Radical} : \zh{木} (mù, « bois, arbre ») — 4 traits.\\
\cle{Structure} : Haut-bas (\zh{屯} au-dessus de \zh{木}). \cle{Ordre} : Écrire \zh{屯} (丿 一 丨 丿 乀) puis \zh{木} (一 丨 丿 丶).
\item \textbf{城} (\emph{chéng}) — ville, enceinte.\\
\cle{Radical} : \zh{土} (tǔ, « terre ») — 3 traits.\\
\cle{Structure} : Gauche-droite. \cle{Ordre} : \zh{土} (一 丨 一) puis partie droite \zh{成} (一 丨 丿 乀 一 一 丨).
\item \textbf{离} (\emph{lí}) — quitter, s'éloigner.\\
\cle{Radical} : \zh{隹} (zhuī, « oiseau court ») — 8 traits (forme simplifiée \zh{离} n'a plus de radical évident, historiquement \zh{隹}).\\
\cle{Ordre} : \zh{禸} (liǎo) en haut (丿 丨 ㇆ 丶) puis \zh{隹} simplifié en bas.
\item \textbf{工} (\emph{gōng}) — travail, ouvrier.\\
\cle{Radical} : \zh{工} (lui-même, radical n°48).\\
\cle{Ordre} : 一 (héng) → 丨 (shù) → 一 (héng). (Le trait du milieu ne touche pas les bords).
\item \textbf{空} (\emph{kōng}) — vide, air, ciel.\\
\cle{Radical} : \zh{穴} (xué, « caverne, trou ») — 5 traits.\\
\cle{Structure} : Haut-bas. \cle{Ordre} : \zh{穴} (丶 丶 ㇆ 丿 丶) puis \zh{工} + \zh{小} (partie basse).
\item \textbf{安} (\emph{ān}) — paix, calme.\\
\cle{Radical} : \zh{宀} (mián, « toit ») — 3 traits.\\
\cle{Structure} : Haut-bas. \cle{Ordre} : \zh{宀} (丶 丿 丶) puis \zh{女} (nǚ, « femme ») : 丿 ㇟ 丶.
\end{enumerate}
\end{exemplebox}

\courssub{Compréhension du texte}

\begin{exercice}[Questions de compréhension sur le texte « 农村的年轻人 »]
Répondez en français (ou en chinois si possible) en vous basant strictement sur le texte.
\begin{enumerate}
\item \textbf{Question globale :} Quel est le thème principal de ce court texte ?
\item \textbf{Détail (Cause) :} Selon le texte, quelles sont les deux raisons principales qui poussent les jeunes à partir vers la ville ?
\item \textbf{Détail (Contraste) :} Citez deux avantages de la campagne et deux inconvénients mentionnés dans le texte.
\item \textbf{Vocabulaire en contexte :} Dans la phrase « \zh{农村的空气干净，环境安静} », que signifie \zh{环境} ? Donnez un synonyme chinois possible vu en classe.
\item \textbf{Inférence / Opinion :} Quelle solution l'auteur propose-t-il à la fin pour retenir les jeunes au village ? Cette solution vous semble-t-elle adaptée au contexte camerounais ? Justifiez brièvement.
\end{enumerate}
\end{exercice}

\begin{corrige}[Exercice : Compréhension du texte]
\begin{enumerate}
\item \textbf{Thème :} L'exode rural (le départ des jeunes de la campagne vers la ville) et ses causes / conséquences, avec une proposition de solution (tourisme).
\item \textbf{Raisons (Ville) :} 1. Beaucoup d'offres d'emploi (\zh{工作机会多}). 2. Salaires plus élevés (\zh{工资也高}).
\item \textbf{Campagne — Avantages :} Air pur (\zh{空气干净}), environnement calme (\zh{环境安静}). \textbf{Inconvénients :} Peu d'hôpitaux (\zh{医院少}), écoles pas bonnes (\zh{学校也不好}).
\item \textbf{\zh{环境} (huánjìng) :} Environnement, cadre de vie. Synonyme possible : \zh{周围} (zhōuwéi, « alentours ») ou \zh{条件} (tiáojiàn, « conditions »).
\item \textbf{Solution :} Développer l'industrie touristique (\zh{发展旅游业}) pour créer de l'emploi local (\zh{让年轻人留在家乡工作}). \emph{Avis personnel attendu :} Oui, le Cameroun a un potentiel touristique (parcs, plages, culture) mais nécessite des infrastructures (routes, hôtels, formation).
\end{enumerate}
\end{corrige}

\courssub{Collocations et exemples}

\begin{propriete}[Collocations fréquentes (associations de mots fixes)]
Ces combinaisons sont des « briques » toutes faites à mémoriser pour parler naturellement de l'exode rural. Respectez l'ordre Verbe + Objet ou Adjectif + Nom.
\end{propriete}

\begin{center}
\begin{tabularx}{\linewidth}{|X|X|X|}
\hline
\rowcolor{popblueL} Collocation (Caractères + Pinyin) & Structure & Traduction / Exemple d'emploi \\
\hline
\zh{离开农村} (\emph{líkāi nóngcūn}) & V + O & Quitter la campagne. \zh{他不想离开农村。} (Il ne veut pas quitter la campagne.) \\
\hline
\zh{去城市} (\emph{qù chéngshì}) & V + O & Aller en ville. \zh{很多年轻人去城市。} (Beaucoup de jeunes vont en ville.) \\
\hline
\zh{找工作} (\emph{zhǎo gōngzuò}) & V + O & Chercher du travail. \zh{他去雅温得找工作。} (Il va à Yaoundé chercher du travail.) \\
\hline
\zh{工作机会多} (\emph{gōngzuò jīhuì duō}) & Suj + Adj & Les offres d'emploi sont nombreuses. \zh{城市的工作机会多。} \\
\hline
\zh{空气干净} (\emph{kōngqì gānjìng}) & Suj + Adj & L'air est pur. \zh{农村的空气干净。} \\
\hline
\zh{生活安静} (\emph{shēnghuó ānjìng}) & Suj + Adj & La vie est calme. \zh{我喜欢生活安静的地方。} \\
\hline
\zh{交通方便} (\emph{jiāotōng fāngbiàn}) & Suj + Adj & Les transports sont pratiques. \zh{城市里交通方便。} \\
\hline
\zh{留在家乡} (\emph{liú zài jiāxiāng}) & V + Prep + O & Rester au pays natal. \zh{我们要留在家乡建设家乡。} (Nous devons rester au pays pour le construire.) \\
\hline
\zh{发展经济} (\emph{fāzhǎn jīngjì}) & V + O & Développer l'économie. \zh{农村应该发展经济。} \\
\hline
\end{tabularx}
\end{center}

\begin{attention}[Pièges de collocation pour francophones]
\begin{itemize}
\item Ne dites pas \zh{找一个工作} (chercher *un* travail) : \zh{工作} est souvent indénombrable ou générique ici, \zh{找工作} est l'expression figée.
\item \zh{机会} (chance/opportunité) prend le classifieur \zh{个} : \zh{一个机会}, mais \zh{工作机会} est un nom composé « offre d'emploi ».
\item \zh{空气} « air » (physique) vs \zh{气} « air/qi/souffle » ou « météo » (\zh{天气}). Ne confondez pas.
\item \zh{方便} est un adjectif (étatif), pas un verbe. « C'est pratique » = \zh{很方便} / \zh{比较方便}. Ne dites pas \zh{这方便我} pour « ça m'arrange » (dire \zh{这对我很方便}).
\end{itemize}
\end{attention}

\courssub{Production guidée : Débat et paragraphe d'opinion}

\textbf{Objectifs de l'activité}

À l'issue de cette activité d'intégration, l'élève sera capable de :
\begin{itemize}
\item prendre la parole en chinois pour défendre une position simple dans un débat organisé (expression orale) ;
\item réinvestir le vocabulaire de la leçon sur l'exode rural (\zh{农村}, \zh{城市}, \zh{年轻人}, \zh{工作机会}, \zh{空气}, \zh{离开}, \zh{生活}…) dans des phrases argumentées ;
\item utiliser correctement la structure d'argumentation \zh{因为……所以……} et la comparaison avec \zh{比} ;
\item écouter et comprendre les arguments de ses camarades (compréhension orale) ;
\item rédiger individuellement un court paragraphe de conclusion (4 à 5 phrases) donnant son avis sur l'exode rural au Cameroun, en suivant le plan de production guidée.
\end{itemize}

\textbf{Situation}

Dans le cadre de la semaine des langues du lycée de Bafoussam, le club de chinois organise un grand débat en langue chinoise. Le thème choisi cette année touche directement la vie des élèves : beaucoup de jeunes quittent leur village pour aller étudier ou travailler à Douala ou à Yaoundé, et le même phénomène existe en Chine. Deux équipes s'affrontent : l'équipe \zh{农村} (la campagne) et l'équipe \zh{城市} (la ville). Chaque équipe doit convaincre le jury, composé du professeur et de trois élèves.

Voici l'affiche du débat, affichée au tableau du club :

\begin{textebox}[Affiche du club de chinois]
辩论会：农村好还是城市好？\\
Biànlùnhuì : Nóngcūn hǎo háishi chéngshì hǎo ?\\
« Débat : vaut-il mieux la campagne ou la ville ? »\\[0.5em]
很多年轻人离开农村，去大城市找工作。\\
Hěn duō niánqīngrén líkāi nóngcūn, qù dà chéngshì zhǎo gōngzuò.\\
« Beaucoup de jeunes quittent la campagne pour aller chercher du travail dans les grandes villes. »\\[0.5em]
你觉得呢？欢迎你参加！\\
Nǐ juéde ne ? Huānyíng nǐ cānjiā !\\
« Et toi, qu'en penses-tu ? Tu es le bienvenu ! »\\[0.5em]
时间：星期五下午四点。地点：中文教室。\\
Shíjiān : Xīngqīwǔ xiàwǔ sì diǎn. Dìdiǎn : Zhōngwén jiàoshì.\\
« Date : vendredi à 16 heures. Lieu : salle de chinois. »
\end{textebox}

\textbf{Consignes}

\begin{enumerate}
\item \textbf{Avant le débat — préparation du lexique (10 min).} Relis le glossaire de la leçon. Avec ton équipe, choisis dans le tableau cinq mots que vous voulez absolument employer pendant le débat (par exemple \zh{空气}, \zh{工作机会}, \zh{交通}, \zh{医院}, \zh{安静}). Écris chaque mot avec son pinyin et sa traduction sur une fiche.
\item \textbf{Préparation des arguments (10 min).} Chaque élève rédige par écrit \textbf{deux arguments} en chinois, en suivant les modèles : \zh{我觉得农村好，因为……} / \zh{城市的工作机会比较多。} Utilise au moins une fois la structure \zh{因为……所以……} et une fois la comparaison avec \zh{比}. Le professeur vérifie les phrases avant le débat.
\item \textbf{Le débat (20 min).} À tour de rôle, un élève de chaque équipe lit ou dit un argument. L'équipe adverse peut poser une question simple (\zh{为什么？} « pourquoi ? »). Le rapporteur de chaque équipe note les arguments au tableau en caractères chinois.
\item \textbf{Écoute et compréhension.} Pendant le débat, note sur ta feuille deux arguments de l'équipe adverse que tu as compris, avec leur traduction française.
\item \textbf{Production écrite individuelle (15 min).} Après le débat, rédige un court paragraphe de conclusion (4 à 5 phrases) donnant ton avis personnel sur l'exode rural au Cameroun, en suivant ce plan :
\begin{enumerate}
\item Phrase 1 : présente le phénomène (\zh{很多年轻人离开农村……}).
\item Phrase 2 : donne une cause avec \zh{因为}.
\item Phrase 3 : compare la ville et la campagne avec \zh{比}.
\item Phrase 4 : donne ton avis avec \zh{我觉得……}.
\item Phrase 5 (facultative) : propose une idée pour le village (\zh{农村应该有……}).
\end{enumerate}
\end{enumerate}

\begin{experience}[Déroulement du débat en classe]
Le professeur tire au sort l'équipe qui commence. Chaque prise de parole dure au maximum 30 secondes. Le jury attribue un point par argument correct en chinois (grammaire et tons) et un point bonus si l'argument utilise une collocation de la leçon (\zh{找工作}, \zh{离开农村}, \zh{工作机会多}). L'équipe qui a le plus de points gagne le débat, mais tous les élèves sont évalués sur leur paragraphe écrit final.
\end{experience}

\textbf{Ressources à mobiliser}

\begin{propriete}[Structures utiles pour argumenter]
\textbf{1. La cause :} \cle{因为} + cause, \cle{所以} + résultat.\\
\zh{因为城市的工作机会多，所以很多年轻人离开农村。}\\
\emph{Yīnwèi chéngshì de gōngzuò jīhuì duō, suǒyǐ hěn duō niánqīngrén líkāi nóngcūn.}\\
« Parce qu'il y a beaucoup d'offres d'emploi en ville, beaucoup de jeunes quittent la campagne. »\\[0.4em]
\textbf{2. La comparaison :} A + \cle{比} + B + adjectif.\\
\zh{农村的空气比城市干净。}\\
\emph{Nóngcūn de kōngqì bǐ chéngshì gānjìng.}\\
« L'air de la campagne est plus pur que celui de la ville. »\\[0.4em]
\textbf{3. L'avis personnel :} \cle{我觉得} + phrase.\\
\zh{我觉得农村的生活很安静。} \emph{Wǒ juéde nóngcūn de shēnghuó hěn ānjìng.} « Je trouve que la vie à la campagne est très calme. »\\[0.4em]
\textbf{4. La concession / opposition :} \cle{可是} / \cle{但是} + phrase.\\
\zh{城市工作机会多，可是生活很忙。} \emph{Chéngshì gōngzuò jīhuì duō, kěshì shēnghuó hěn máng.} « La ville a du travail, mais la vie est très occupée. »
\end{propriete}

\begin{attention}[Erreurs fréquentes des francophones]
\begin{itemize}
\item Ne traduis pas « plus de » par \zh{更} seul : avec \zh{比}, l'adjectif suffit : \zh{城市比农村方便} (et non \zh{城市比农村很方便}, le \zh{很} est de trop après \zh{比}).
\item \zh{因为} et \zh{所以} peuvent se combiner dans la même phrase en chinois, contrairement au français (« parce que… donc… » est incorrect en français mais normal en chinois).
\item Attention aux tons de \zh{líkāi} (离开, 2\textsuperscript{e} ton + 1\textsuperscript{er} ton) et de \zh{jīhuì} (机会, 1\textsuperscript{e} ton + 4\textsuperscript{e} ton).
\item On dit \zh{在农村} « à la campagne » et \zh{在城市里} « en ville » : n'oublie pas la préposition \zh{在} + lieu devant le verbe.
\item \zh{安静} s'écrit \emph{ānjìng} (un seul mot, pas d'espace) ; \zh{干净} s'écrit \emph{gānjìng}.
\end{itemize}
\end{attention}

\textbf{Proposition de production et corrigé commenté}

\begin{exoresolu}[Paragraphe de conclusion : modèle rédigé]
Énoncé : rédige un paragraphe de 4 à 5 phrases donnant ton avis sur l'exode rural au Cameroun, en suivant le plan de production guidée.
\tcblower
\textbf{Production modèle :}\\[0.4em]
\zh{在喀麦隆，很多年轻人离开农村，去雅温得和杜阿拉。}\\
\emph{Zài Kāmàilóng, hěn duō niánqīngrén líkāi nóngcūn, qù Yǎwēndé hé Dùālā.}\\
« Au Cameroun, beaucoup de jeunes quittent la campagne pour Yaoundé et Douala. »\\[0.3em]
\zh{因为城市的工作机会比较多，学校也比较好。}\\
\emph{Yīnwèi chéngshì de gōngzuò jīhuì bǐjiào duō, xuéxiào yě bǐjiào hǎo.}\\
« Parce qu'en ville il y a plus d'offres d'emploi et les écoles sont aussi meilleures. »\\[0.3em]
\zh{可是农村的空气比城市干净，生活也比城市安静。}\\
\emph{Kěshì nóngcūn de kōngqì bǐ chéngshì gānjìng, shēnghuó yě bǐ chéngshì ānjìng.}\\
« Mais l'air de la campagne est plus pur que celui de la ville, et la vie y est plus calme. »\\[0.3em]
\zh{我觉得农村和城市都很好。}\\
\emph{Wǒ juéde nóngcūn hé chéngshì dōu hěn hǎo.}\\
« Je pense que la campagne et la ville sont toutes les deux bien. »\\[0.3em]
\zh{农村应该有更多的医院和工作机会。}\\
\emph{Nóngcūn yīnggāi yǒu gèng duō de yīyuàn hé gōngzuò jīhuì.}\\
« La campagne devrait avoir plus d'hôpitaux et d'offres d'emploi. »\\[0.5em]
\textbf{Commentaire des choix de langue :}
\begin{itemize}
\item Phrase 1 : le lieu est placé juste après le sujet (\zh{在喀麦隆} en tête), conformément à l'ordre chinois Sujet + Lieu + Verbe ; les noms propres \zh{雅温得} (Yaoundé) et \zh{杜阿拉} (Douala) prennent la majuscule en pinyin.
\item Phrase 2 : la cause est introduite par \zh{因为} ; \zh{比较多} « relativement plus » évite la faute \zh{比……很} (on n'utilise pas \zh{比} ici, mais \zh{比较} adverbe de degré).
\item Phrase 3 : la concession est marquée par \zh{可是} « mais » ; la comparaison A + \zh{比} + B + adjectif est appliquée deux fois, sans \zh{很} après \zh{比}.
\item Phrase 4 : l'avis personnel s'exprime avec \zh{我觉得} ; \zh{都} « tous les deux » résume les deux positions du débat.
\item Phrase 5 : la proposition utilise l'auxiliaire \zh{应该} « il faudrait » + verbe \zh{有}, structure simple et correcte au niveau HSK 2-3. \zh{更多} « plus » s'emploie sans \zh{比} ici.
\end{itemize}
\end{exoresolu}

\textbf{Bilan de l'activité}

Cette activité d'intégration a permis de mobiliser l'ensemble des savoir-faire de la leçon dans une tâche finale de communication réelle :
\begin{itemize}
\item \textbf{compréhension orale} : écouter et noter les arguments de l'équipe adverse ;
\item \textbf{lexique} : réemployer le vocabulaire de l'exode rural dans des phrases personnelles ;
\item \textbf{grammaire} : manier la cause (\zh{因为……所以……}), la comparaison (\zh{比}) et l'avis (\zh{我觉得}) ;
\item \textbf{expression orale} : prendre la parole en public, avec un temps limité, dans un cadre ludique ;
\item \textbf{expression écrite} : rédiger un paragraphe structuré en cinq phrases, étape par étape.
\end{itemize}

\begin{aretenir}[Ce que je dois retenir]
Pour donner mon avis en chinois, je peux suivre le plan : phénomène (\zh{很多年轻人离开农村}) → cause (\zh{因为……}) → comparaison (\zh{A 比 B + adjectif}) → avis (\zh{我觉得……}) → proposition (\zh{应该……}). Après \zh{比}, jamais de \zh{很} devant l'adjectif.
\end{aretenir}

\begin{savaistu}[L'exode rural : Cameroun et Chine]
Au Cameroun comme en Chine, de nombreux jeunes quittent les villages pour les grandes villes : Douala et Yaoundé d'un côté, Pékin, Shanghai et Guangzhou de l'autre. En chinois, les travailleurs venus de la campagne sont appelés \zh{农民工} (\emph{nóngmíngōng}, « ouvriers paysans »). Aujourd'hui, dans les deux pays, certains jeunes font le chemin inverse et retournent au village pour y créer des activités, par exemple dans l'agriculture ou le tourisme : on parle alors de « retour à la terre » (\zh{返乡创业} \emph{fǎnxiāng chuàngyè}).
\end{savaistu}

\newpage

\coursec{Méthodes et exercices résolus}

\courssub{Méthodes et exercices résolus}

\begin{methode}[Lire et comprendre un texte authentique sur l'exode rural]
\textbf{Objectif :} Extraire l'information principale et les détails d'un texte chinois (HSK 3-4) sur un thème de société.
\textbf{Démarche en 4 étapes :}
\begin{enumerate}
    \item \textbf{Repérage global :} Lire le titre (\zh{农村的年轻人}) et la première phrase de chaque paragraphe. Identifier le sujet (qui ? \zh{年轻人}), le lieu (où ? \zh{农村} vs \zh{城市}), le problème (quoi ? \zh{外出打工} / \zh{留守}).
    \item \textbf{Lecture détaillée avec annotations :} Surligner les \textbf{mots-clés connus} (vocabulaire HSK 1-3). Entourer les \textbf{connecteurs logiques} (\zh{因为、所以、但是、越来越}) qui structurent l'argumentation. Noter le pinyin des caractères nouveaux au-dessus du texte.
    \item \textbf{Segmentation en unités de sens :} Découper le texte en 3 parties : \emph{Constat (départ)}, \emph{Causes (raisons)}, \emph{Conséquences et avenir (retour / développement)}.
    \item \textbf{Vérification croisée :} Relire la traduction française fournie pour valider sa compréhension. Noter les écarts de structure (ex. : \zh{越来越 + adjectif} = « de plus en plus... »).
\end{enumerate}
\textbf{Réflexe :} Ne pas traduire mot à mot. Chercher l'idée directrice de chaque phrase.
\end{methode}

\begin{exoresolu}[Compréhension globale du texte « 农村的年轻人 »]
\textbf{Énoncé :} Lisez le texte ci-dessous. Répondez en français aux questions 1 à 3.

\begin{textebox}[Texte support original : 农村的年轻人]
农村的年轻人越来越多外出打工。\\
\emph{Nóngcūn de niánqīng rén yuè lái yuè duō wàichū dǎgōng.}\\
« Les jeunes de la campagne partent de plus en plus nombreux travailler ailleurs. »\\
因为城市的工作机会多，生活条件好。\\
\emph{Yīnwèi chéngshì de gōngzuò jīhuì duō, shēnghuó tiáojiàn hǎo.}\\
« Parce qu'en ville, les opportunités d'emploi sont nombreuses et les conditions de vie sont bonnes. »\\
但是，农村留下了很多老人和孩子。\\
\emph{Dànshì, nóngcūn liúxià le hěn duō lǎorén hé háizi.}\\
« Mais la campagne a gardé (derrière elle) beaucoup de personnes âgées et d'enfants. »\\
现在，国家发展农村经济，年轻人回来建设家乡。\\
\emph{Xiànzài, guójiā fāzhǎn nóngcūn jīngjì, niánqīng rén huílái jiànshè jiāxiāng.}\\
« Maintenant, l'État développe l'économie rurale, et les jeunes reviennent construire leur village natal. »
\end{textebox}

\begin{enumerate}
    \item Quel est le sujet principal du texte ?
    \item Citez deux raisons données pour expliquer le départ des jeunes.
    \item Quelle est la situation actuelle décrite dans la dernière phrase ?
\end{enumerate}

\tcblower
\textbf{Solution détaillée :}

\textbf{Étape 1 : Analyse du vocabulaire clé.}
\begin{itemize}
    \item \zh{农村} (nóngcūn) : campagne / zone rurale.
    \item \zh{年轻人} (niánqīng rén) : jeunes / jeunes gens.
    \item \zh{越来越多} (yuè lái yuè duō) : de plus en plus (nombreux).
    \item \zh{外出打工} (wàichū dǎgōng) : partir travailler ailleurs (migrer pour le travail).
    \item \zh{因为...所以...} (yīnwèi... suǒyǐ...) : parce que... donc... (structure cause-conséquence).
    \item \zh{工作机会} (gōngzuò jīhuì) : opportunités d'emploi.
    \item \zh{留下} (liúxià) : laisser derrière soi.
    \item \zh{留守} (liúshǒu) : « restés sur place » (terme spécifique pour les enfants et les personnes âgées restés au village).
    \item \zh{发展} (fāzhǎn) : développer / développement.
    \item \zh{建设家乡} (jiànshè jiāxiāng) : construire / développer sa région natale.
\end{itemize}

\textbf{Étape 2 : Réponses justifiées.}
\begin{enumerate}
    \item \textbf{Sujet principal :} L'exode rural des jeunes (\zh{农村的年轻人外出打工}) et le renversement de tendance actuel (\zh{年轻人回来建设家乡}).
    \item \textbf{Deux raisons (phrase 2) :}
        \begin{itemize}
            \item \zh{城市的工作机会多} : beaucoup d'opportunités d'emploi en ville.
            \item \zh{生活条件好} : meilleures conditions de vie en ville.
        \end{itemize}
    \item \textbf{Situation actuelle (phrase 4) :} L'État développe l'économie rurale (\zh{国家发展农村经济}), ce qui incite les jeunes à revenir construire leur village natal (\zh{年轻人回来建设家乡}).
\end{enumerate}

\textbf{Conclusion :} Le texte suit une structure classique \textbf{Constat $\rightarrow$ Causes $\rightarrow$ Problème (\zh{留守}) $\rightarrow$ Solution / Perspective}. Maîtriser les connecteurs (\zh{因为/所以、但是、现在}) est essentiel pour suivre la logique.
\end{exoresolu}

\begin{methode}[Identifier, prononcer et classer le vocabulaire thématique (radicaux et tons)]
\textbf{Objectif :} Maîtriser l'orthographe (caractères), la prononciation (pinyin + tons) et la nature grammaticale du lexique « Exode rural / Société ».
\textbf{Démarche en 3 étapes :}
\begin{enumerate}
    \item \textbf{Analyse graphique (clés / radicaux) :} Identifier la \textbf{clé de dictionnaire} (部首 bùshǒu) pour deviner le sens ou classer le caractère. Ex. : \zh{农} (clé \zh{冖} « couverture »), \zh{城} (clé \zh{土} « terre » $\rightarrow$ rempart de terre, ville), \zh{留} (clé \zh{田} « champ » $\rightarrow$ ce qui reste au champ).
    \item \textbf{Ordre des traits (笔顺 bǐshùn) :} Appliquer les règles universelles : \emph{haut $\rightarrow$ bas, gauche $\rightarrow$ droite, horizontal avant vertical, extérieur avant intérieur, fermer en dernier}. Ex. : \zh{打} (5 traits : la main \zh{扌} en 3 traits, puis \zh{丁} en 2 traits). \zh{工} : horizontale, verticale, horizontale (3 traits).
    \item \textbf{Mémorisation par collocations (搭配 dāpèi) :} Ne pas apprendre le mot isolé. Associer : \zh{外出打工} (verbe + verbe), \zh{留守儿童} (nom composé), \zh{建设家乡} (verbe + objet), \zh{生活条件} (nom + nom).
\end{enumerate}
\textbf{Structure à retenir :} \cle{Caractère = clé (sens) + partie phonétique (son) + ordre des traits (mémoire motrice)}.
\end{methode}

\begin{exoresolu}[Vocabulaire : radicaux, ordre des traits et emploi]
\textbf{Énoncé :} Complétez le tableau ci-dessous pour les 5 caractères cibles. Ensuite, écrivez une phrase correcte pour chacun en utilisant la collocation indiquée.

\begin{center}
\begin{tabularx}{\linewidth}{|X|X|X|X|X|}
\hline
\rowcolor{popblueL} \textbf{Caractère} & \textbf{Clé (radical)} & \textbf{Nb traits} & \textbf{Collocation imposée} & \textbf{Phrase chinoise (caractères + pinyin)} \\
\hline
\zh{农} & & & \zh{农村} & \\
\hline
\zh{城} & & & \zh{城市} & \\
\hline
\zh{打} & & & \zh{打工} & \\
\hline
\zh{留} & & & \zh{留守} & \\
\hline
\zh{建} & & & \zh{建设} & \\
\hline
\end{tabularx}
\end{center}

\tcblower
\textbf{Solution détaillée :}

\begin{center}
\begin{tabularx}{\linewidth}{|X|X|X|X|X|}
\hline
\rowcolor{popblueL} \textbf{Caractère} & \textbf{Clé (radical)} & \textbf{Nb traits} & \textbf{Collocation} & \textbf{Phrase chinoise} \\
\hline
\zh{农} & \zh{冖} (mì, couverture) & 6 & \zh{农村} & \zh{我的家乡在农村。} (Wǒ de jiāxiāng zài nóngcūn.) \\
\hline
\zh{城} & \zh{土} (tǔ, terre) & 9 & \zh{城市} & \zh{他去城市工作了。} (Tā qù chéngshì gōngzuò le.) \\
\hline
\zh{打} & \zh{扌} (tíshǒu, main) & 5 & \zh{打工} & \zh{哥哥在广州打工。} (Gēge zài Guǎngzhōu dǎgōng.) \\
\hline
\zh{留} & \zh{田} (tián, champ) & 10 & \zh{留守} & \zh{农村有很多留守儿童。} (Nóngcūn yǒu hěn duō liúshǒu értóng.) \\
\hline
\zh{建} & \zh{廴} (yǐn, longue marche) & 8 & \zh{建设} & \zh{我们要建设美丽的家乡。} (Wǒmen yào jiànshè měilì de jiāxiāng.) \\
\hline
\end{tabularx}
\end{center}

\textbf{Justifications graphiques :}
\begin{itemize}
    \item \zh{农} (6 traits) : forme simplifiée de \zh{農} (traditionnel). Clé \zh{冖}. Ordre : point, horizontale-crochet (\zh{冖}), puis la partie inférieure : tombant, vertical-relevé, tombant, appuyée (de haut en bas).
    \item \zh{城} (9 traits) : clé \zh{土} (terre, construction) à gauche ; partie phonétique \zh{成} (chéng) à droite. Ordre : \zh{土} (horizontale, verticale, horizontale) puis \zh{成} (6 traits, de haut en bas).
    \item \zh{打} (5 traits) : clé \zh{扌} (action de la main) ; phonétique \zh{丁} (dīng). Ordre : \zh{扌} (horizontale, verticale-crochet, horizontale montante) puis \zh{丁} (horizontale, verticale-crochet).
    \item \zh{留} (10 traits) : clé \zh{田} (champ) en bas ; partie supérieure \zh{卯} (mǎo, phonétique). Ordre : partie supérieure d'abord (de gauche à droite), puis \zh{田} (fermer la boîte en dernier).
    \item \zh{建} (8 traits) : clé \zh{廴} ; intérieur \zh{聿} (yù, pinceau). Ordre : \zh{聿} d'abord (6 traits, horizontales puis verticale), puis \zh{廴} (2 traits) qui enveloppe par le bas à gauche.
\end{itemize}

\textbf{Point de grammaire :} \zh{留守} fonctionne comme un verbe « rester garder (les lieux) » ou un déterminant : \zh{留守儿童} = « enfants restés au village » (terme officiel). \zh{建设} est un verbe transitif « construire, édifier » $\rightarrow$ \zh{建设家乡、建设国家}.
\end{exoresolu}

\begin{methode}[Analyser les structures grammaticales clés du texte (comparatif, cause/effet, changement d'état)]
\textbf{Objectif :} Reconnaître et reproduire les 3 structures grammaticales HSK 3-4 majeures du texte pour l'expression écrite.
\textbf{Structures cibles :}
\begin{enumerate}
    \item \textbf{Comparatif d'inégalité :} \cle{A + 比 + B + adjectif}. \textit{Ex. : \zh{城市比农村发达。}} (Chéngshì bǐ nóngcūn fādá.) « La ville est plus développée que la campagne. » \textbf{Attention :} pas de \zh{很} après \zh{比}. Négation : \cle{A + 没有 + B + adjectif} = « A n'est pas aussi... que B ».
    \item \textbf{Structure cause-conséquence :} \cle{因为 ..., 所以 ...} « parce que..., donc... ». En chinois, les deux marqueurs sont souvent présents ensemble (corrélatifs), contrairement au français où « parce que » suffit. \textit{Ex. : \zh{因为下雨，所以我不去。}} (Yīnwèi xià yǔ, suǒyǐ wǒ bú qù.)
    \item \textbf{Changement d'état progressif :} \cle{越来越 + adjectif / verbe psychologique} « de plus en plus... ». \textit{Ex. : \zh{天气越来越冷。}} (Tiānqì yuè lái yuè lěng.) \zh{外出打工的年轻人越来越多。}
\end{enumerate}
\textbf{Méthode de transformation :} Phrase française $\rightarrow$ identifier la relation logique $\rightarrow$ choisir la structure chinoise $\rightarrow$ ordonner les éléments (sujet en premier !).
\end{methode}

\begin{exoresolu}[Transformation de phrases et emploi des structures]
\textbf{Énoncé :} Traduisez les phrases suivantes en chinois (caractères + pinyin) en utilisant la structure imposée entre parenthèses.

\begin{enumerate}
    \item Les conditions de vie en ville sont meilleures qu'à la campagne. (Structure \zh{比})
    \item Parce qu'il y a du travail, il part à la ville. (Structure \zh{因为...所以...})
    \item Les jeunes qui reviennent au village sont de plus en plus nombreux. (Structure \zh{越来越})
    \item Mon village natal n'est pas aussi grand que Yaoundé. (Structure \zh{没有} pour comparatif négatif)
\end{enumerate}

\tcblower
\textbf{Solution détaillée :}

\begin{enumerate}
    \item \textbf{Structure \zh{比} :} \zh{城市的生活条件比农村好。} \\
    (Chéngshì de shēnghuó tiáojiàn bǐ nóngcūn hǎo.) \\
    \textit{Analyse :} Sujet A (\zh{城市的生活条件}) + \zh{比} + Sujet B (\zh{农村}) + adjectif (\zh{好}). \textbf{Piège évité :} ne pas dire \zh{比农村很好} (pas de \zh{很} dans un comparatif).
    \item \textbf{Structure \zh{因为...所以...} :} \zh{因为有工作，所以他去城市。} \\
    (Yīnwèi yǒu gōngzuò, suǒyǐ tā qù chéngshì.) \\
    \textit{Analyse :} \zh{因为} introduit la cause (en tête de phrase), \zh{所以} introduit la conséquence (en début de seconde proposition). La cause est ici impersonnelle : \zh{有工作} « il y a du travail ».
    \item \textbf{Structure \zh{越来越} :} \zh{回来建设家乡的年轻人越来越多。} \\
    (Huílái jiànshè jiāxiāng de niánqīng rén yuè lái yuè duō.) \\
    \textit{Analyse :} sujet complexe (\zh{回来建设家乡的年轻人}, « les jeunes qui reviennent construire le village natal ») + \zh{越来越多}. \zh{多} est l'adjectif « nombreux » ; ne pas le remplacer par \zh{大} (grand) ou \zh{好} (bon).
    \item \textbf{Structure \zh{没有} (comparatif négatif) :} \zh{我的家乡没有雅温得大。} \\
    (Wǒ de jiāxiāng méiyǒu Yǎwēndé dà.) \\
    \textit{Analyse :} A + \zh{没有} + B + adjectif = « A n'est pas aussi... que B ». \zh{雅温得} (Yǎwēndé) = Yaoundé. \textbf{Erreur fréquente :} employer \zh{不比} au sens de « pas aussi... que » ; \zh{不比} signifie plutôt « n'est pas plus... que ».
\end{enumerate}

\textbf{Bilan :} La place du sujet est immuable (début de phrase). Les compléments de lieu (\zh{在城市}) se placent \textbf{avant} le verbe (\zh{工作}), jamais après.
\end{exoresolu}

\begin{methode}[Répondre aux questions de compréhension détaillée en chinois]
\textbf{Objectif :} Produire des réponses courtes, grammaticalement correctes, en réemployant le vocabulaire et les structures du texte.
\textbf{Démarche en 3 étapes :}
\begin{enumerate}
    \item \textbf{Analyser la question :} Identifier le mot interrogatif (\zh{谁} qui, \zh{什么} quoi, \zh{为什么} pourquoi, \zh{怎么} comment, \zh{多少} combien). Identifier l'aspect (\zh{了} accompli, \zh{过} expérience).
    \item \textbf{Construire la réponse :} Reprendre le \textbf{sujet de la question} + \textbf{verbe/structure du texte} + \textbf{information demandée}. Éviter le « oui/non » nu (\zh{是/不是}, \zh{有/没有}) quand la question porte sur un contenu.
    \item \textbf{Vérifier tons et caractères :} Vérifier les tons de sandhi (\zh{一}, \zh{不}) et l'écriture des caractères clés (\zh{留守、建设、打工}).
\end{enumerate}
\textbf{Règle d'or :} La réponse doit pouvoir remplacer la question dans le dialogue sans incohérence syntaxique.
\end{methode}

\begin{exoresolu}[Questions / réponses détaillées sur le texte]
\textbf{Énoncé :} Répondez en chinois (caractères + pinyin) aux questions suivantes d'après le texte de la leçon.

\begin{enumerate}
    \item \zh{为什么农村的年轻人外出打工？}
    \item \zh{农村留下了谁？}
    \item \zh{现在国家怎么做？}
    \item \zh{年轻人现在怎么办？}
\end{enumerate}

\tcblower
\textbf{Solution détaillée :}

\begin{enumerate}
    \item \textbf{Q :} \zh{为什么...？} (pourquoi) $\rightarrow$ réponse attendue : la cause (\zh{因为...}).\\
    \textbf{R :} \zh{因为城市的工作机会多，生活条件好。} \\
    (Yīnwèi chéngshì de gōngzuò jīhuì duō, shēnghuó tiáojiàn hǎo.) \\
    \textit{Note :} dans une réponse courte, on peut omettre \zh{所以}, mais garder \zh{因为} marque clairement la cause.
    \item \textbf{Q :} \zh{谁} (qui) est le complément de \zh{留下} (laisser derrière soi).\\
    \textbf{R :} \zh{农村留下了很多老人和孩子。} \\
    (Nóngcūn liúxià le hěn duō lǎorén hé háizi.) \\
    \textit{Grammaire :} \zh{了} (le) marque le résultat : ils « sont restés là ». \zh{老人} (personnes âgées) et \zh{孩子} (enfants) $\rightarrow$ d'où les termes \zh{留守老人、留守儿童}.
    \item \textbf{Q :} \zh{怎么做} (que fait-on / comment fait-on).\\
    \textbf{R :} \zh{国家发展农村经济。} \\
    (Guójiā fāzhǎn nóngcūn jīngjì.) \\
    \textit{Structure :} sujet (\zh{国家}) + verbe (\zh{发展}) + objet (\zh{农村经济}).
    \item \textbf{Q :} \zh{怎么办} (que faire).\\
    \textbf{R :} \zh{年轻人回来建设家乡。} \\
    (Niánqīng rén huílái jiànshè jiāxiāng.) \\
    \textit{Structure :} verbe de mouvement \zh{回来} (revenir) + verbe d'action \zh{建设} (construire) + objet \zh{家乡}.
\end{enumerate}

\textbf{Entraînement oral :} Lire les questions/réponses à haute voix en respectant les tons (ex. : \zh{huílái} 2-2, \zh{jiànshè} 4-4). Attention au ton neutre sur \zh{了} (le) et \zh{的} (de).
\end{exoresolu}

\begin{methode}[Rédiger un court message / paragraphe argumentatif (expression écrite guidée)]
\textbf{Objectif :} Produire un texte cohérent (60-80 caractères) sur le thème « Mon avis sur le retour des jeunes au village », en réinvestissant vocabulaire et grammaire.
\textbf{Structure type « message / micro-essai » (3 phrases) :}
\begin{enumerate}
    \item \textbf{Constat / accroche :} \zh{现在，回农村的年轻人越来越多。} (Xiànzài, huí nóngcūn de niánqīng rén yuè lái yuè duō.) Contexte actuel + \zh{越来越}.
    \item \textbf{Argumentation (cause / bénéfice) :} \zh{因为家乡发展快，工作机会多。} (\zh{因为} + raison économique). \zh{而且，照顾留守老人很重要。} (\zh{而且} « de plus » + raison sociale et familiale).
    \item \textbf{Engagement / conclusion :} \zh{我也想回去建设家乡。} (\zh{也} « aussi » + \zh{想} volonté + \zh{建设家乡} but).
\end{enumerate}
\textbf{Connecteurs logiques obligatoires :} \zh{现在} (temps), \zh{因为/所以} (cause), \zh{而且} (addition), \zh{但是} (concession).
\textbf{Contraintes examen :} écriture soignée (caractères carrés), pinyin exact si demandé, pas de français.
\end{methode}

\begin{exoresolu}[Production écrite : « Le retour au village »]
\textbf{Énoncé :} Vous êtes \zh{李明} (Lǐ Míng), élève camerounais à Yaoundé. Écrivez un court message à votre ami \zh{王伟} (Wáng Wěi) pour lui dire que vous rentrez au village (\zh{老家} lǎojiā) pour les vacances et aider vos parents. Utilisez au moins 5 mots du vocabulaire de la leçon et la structure \zh{因为...所以...}.

\textbf{Contraintes :} 60-80 caractères chinois. Format message : destinataire, corps, signature, date.

\tcblower
\textbf{Solution détaillée (modèle de rédaction) :}

\begin{textebox}[Production de l'élève (modèle corrigé)]
王伟：\\
现在我回老家了。因为农村发展快，有工作机会，而且要照顾留守的老人，所以我决定回来建设家乡，不去城市打工了。\\
\emph{Wáng Wěi: Xiànzài wǒ huí lǎojiā le. Yīnwèi nóngcūn fāzhǎn kuài, yǒu gōngzuò jīhuì, érqiě yào zhàogù liúshǒu de lǎorén, suǒyǐ wǒ juédìng huílái jiànshè jiāxiāng, bú qù chéngshì dǎgōng le.}\\
« Wang Wei : je suis maintenant rentré au village. Parce que la campagne se développe vite, qu'il y a des opportunités d'emploi et qu'il faut s'occuper des personnes âgées restées au village, j'ai décidé de revenir construire mon village natal et de ne plus partir travailler en ville. »\\
李明\\
十月五日
\end{textebox}

\textbf{Analyse linguistique du modèle :}
\begin{itemize}
    \item \textbf{Vocabulaire réinvesti (6 items) :} \zh{老家} (village natal), \zh{发展} (développer), \zh{工作机会} (opportunités d'emploi), \zh{留守} (restés sur place), \zh{建设家乡} (construire le village natal), \zh{打工} (travailler comme migrant).
    \item \textbf{Grammaire :}
        \begin{itemize}
            \item \zh{现在我回老家了} : \zh{了} (le) marque le changement d'état (« je suis rentré »).
            \item \zh{因为...，所以我决定...} : structure cause-effet complète. \zh{决定} (juédìng) « décider » + verbe.
            \item \zh{而且要照顾...} : \zh{而且} (érqiě) « de plus » coordonne deux raisons. \zh{照顾} (zhàogù) « s'occuper de ».
            \item \zh{不去城市打工了} : négation \zh{不} + verbe + \zh{了} (cessation d'action : « je ne vais plus... »).
        \end{itemize}
    \item \textbf{Format :} destinataire \zh{王伟：} (deux-points chinois), corps du message, signature \zh{李明}, date \zh{十月五日} (5 octobre).
\end{itemize}

\textbf{Variante plus formelle (lettre) :}\\
\zh{亲爱的王伟：}\\
\zh{你好！现在国家发展农村经济，我的家乡变化很大。因为那里有好的工作机会，而且我要照顾留守的爷爷奶奶，所以我决定毕业后回来建设家乡。祝你学习进步！}\\
\emph{Qīn'ài de Wáng Wěi: Nǐ hǎo! Xiànzài guójiā fāzhǎn nóngcūn jīngjì, wǒ de jiāxiāng biànhuà hěn dà. Yīnwèi nàli yǒu hǎo de gōngzuò jīhuì, érqiě wǒ yào zhàogù liúshǒu de yéye nǎinai, suǒyǐ wǒ juédìng bìyè hòu huílái jiànshè jiāxiāng. Zhù nǐ xuéxí jìnbù!}\\
« Cher Wang Wei : bonjour ! Maintenant, l'État développe l'économie rurale et mon village natal a beaucoup changé. Parce qu'il y a là-bas de bonnes opportunités d'emploi et que je dois m'occuper de mes grands-parents restés au village, j'ai décidé de revenir construire mon village après mes études. Je te souhaite de progresser dans tes études ! »\\
\zh{李明}\\
\zh{十月五日}\\
\textbf{Note :} \zh{亲爱} (qīn'ài) « cher », \zh{毕业后} (bìyè hòu) « après les études », \zh{祝...进步} formule de politesse.
\end{exoresolu}

\begin{attention}[Les 5 erreurs qui coûtent des points à l'examen (ZHA23)]
\begin{enumerate}
    \item \textbf{Tons et sandhi (les pièges \zh{一} / \zh{不} / \zh{两}) :} \zh{一} (yī) devient \zh{yí} (2\up{e} ton) devant un 4\up{e} ton (\zh{一个} yí ge), et \zh{yì} (4\up{e} ton) devant un 1\up{er}, 2\up{e} ou 3\up{e} ton (\zh{一起} yìqǐ). \zh{不} (bù) devient \zh{bú} (2\up{e} ton) devant un 4\up{e} ton (\zh{不是} bú shì). \zh{两} (liǎng) s'emploie pour « deux » (quantité, devant un classificateur) au lieu de \zh{二} (èr) (comptage). \textit{Ex. : \zh{两个年轻人} (liǎng gè niánqīng rén), jamais \zh{二个}.}
    \item \textbf{Caractères « jumeaux » confondus (形近字 xíngjìnzì) :}
    \begin{itemize}
        \item \zh{农} (nóng, campagne) vs \zh{衣} (yī, vêtement) : formes proches, sens très différents.
        \item \zh{留} (liú, rester) vs \zh{流} (liú, couler) : clé \zh{田} (champ) vs clé \zh{氵} (eau). \textit{Astuce : l'eau coule (\zh{流}), ce qui reste au champ demeure (\zh{留}).}
        \item \zh{建} (jiàn, construire) vs \zh{健} (jiàn, santé) : clé \zh{廴} vs clé \zh{亻} (homme).
    \end{itemize}
    \item \textbf{Ordre des mots : le complément de lieu se place AVANT le verbe.} \textbf{FAUX :} \zh{我工作在城市。} \textbf{JUSTE :} \zh{我在城市工作。} (Wǒ zài chéngshì gōngzuò.) ou \zh{我去城市工作。} (Wǒ qù chéngshì gōngzuò.) Règle : \cle{Sujet + temps + lieu + verbe + complément}.
    \item \textbf{Les mots-outils \zh{的、地、得} (les « trois frères » de) :}
    \begin{itemize}
        \item \zh{的} (de) : après un adjectif ou un nom, devant un nom $\rightarrow$ \zh{农村的年轻人} (les jeunes \textbf{de} la campagne).
        \item \zh{地} (de) : après un adverbe, devant un verbe $\rightarrow$ \zh{努力地建设} (construire \textbf{avec ardeur}).
        \item \zh{得} (de) : après un verbe, devant un complément de manière ou de degré $\rightarrow$ \zh{建设得很好} (construit \textbf{bien}).
    \end{itemize}
    \item \textbf{Classificateurs (量词 liàngcí) obligatoires :} jamais de nombre + nom sans classificateur : \zh{一个年轻人} (yí gè niánqīng rén), \zh{一条路} (yì tiáo lù, une route), \zh{一座城市} (yí zuò chéngshì, une ville). \textbf{Erreur fatale :} \zh{两个工作} $\rightarrow$ dire \zh{两份工作} (liǎng fèn gōngzuò, deux emplois ; \zh{份} fèn pour les emplois et les tâches).
\end{enumerate}
\textbf{Conseil final :} Relisez votre copie en \textbf{chuchotant le pinyin} : l'oreille capte les erreurs de tons et d'ordre des mots que l'œil ne voit pas.
\end{attention}

\newpage

\coursec{Exercices}

\courssub{Je vérifie mes connaissances}

\begin{exercice}[QCM : le vocabulaire de l'exode rural]
Pour chaque question, choisis la bonne réponse.
\begin{enumerate}
\item \zh{农村} (nóngcūn) signifie :
\begin{enumerate}
\item la ville \item la campagne \item la capitale
\end{enumerate}
\item \zh{外流} (wàiliú) signifie :
\begin{enumerate}
\item rentrer au village \item partir vers l'extérieur \item voyager en avion
\end{enumerate}
\item \zh{打工} (dǎgōng) signifie :
\begin{enumerate}
\item travailler comme ouvrier ou employé \item faire du sport \item étudier à l'université
\end{enumerate}
\item \zh{机会} (jīhuì) signifie :
\begin{enumerate}
\item le problème \item le salaire \item l'occasion, la chance
\end{enumerate}
\end{enumerate}
\end{exercice}

\begin{exercice}[Vrai ou faux]
Lis le texte de la leçon \zh{农村的年轻人} puis dis si les phrases suivantes sont vraies ou fausses. Justifie ta réponse par une phrase du texte.
\begin{enumerate}
\item \zh{很多年轻人离开农村去城市。} (Hěn duō niánqīngrén líkāi nóngcūn qù chéngshì.)
\item \zh{城市里工作很少，工资很低。} (Chéngshì lǐ gōngzuò hěn shǎo, gōngzī hěn dī.)
\item \zh{年轻人走了以后，农村里只剩下老人和孩子。} (Niánqīngrén zǒu le yǐhòu, nóngcūn lǐ zhǐ shèngxià lǎorén hé háizi.)
\item \zh{政府不想帮助农村发展。} (Zhèngfǔ bù xiǎng bāngzhù nóngcūn fāzhǎn.)
\end{enumerate}
\end{exercice}

\begin{exercice}[Trouve le mot du glossaire]
Complète chaque définition avec le mot chinois qui convient (caractères et pinyin).
\begin{enumerate}
\item L'argent qu'on reçoit pour son travail : \_\_\_\_\_\_\_\_
\item Les personnes âgées : \_\_\_\_\_\_\_\_
\item Quitter un endroit : \_\_\_\_\_\_\_\_
\item Une solution, un moyen : \_\_\_\_\_\_\_\_
\item Le nombre d'habitants d'un lieu : \_\_\_\_\_\_\_\_
\end{enumerate}
\end{exercice}

\begin{exercice}[Associe le pinyin aux caractères]
Associe chaque caractère ou mot à son pinyin, puis donne sa traduction française.
\begin{center}
\begin{tabularx}{\linewidth}{|X|X|}\hline
\rowcolor{popblueL} Caractères & Pinyin \\ \hline
1. 村 & a. fāzhǎn \\ \hline
2. 城市 & b. líkāi \\ \hline
3. 离开 & c. cūn \\ \hline
4. 发展 & d. chéngshì \\ \hline
5. 剩下 & e. shèngxià \\ \hline
\end{tabularx}
\end{center}
\end{exercice}

\courssub{Je m'entraîne}

\begin{exercice}[Traduction vers le chinois]
Traduis les phrases suivantes en chinois (caractères obligatoires, pinyin en plus si possible).
\begin{enumerate}
\item Beaucoup de jeunes quittent leur village parce qu'il y a peu de travail à la campagne.
\item En ville, il y a plus d'opportunités et les salaires sont plus élevés.
\item Mon grand-père ne veut pas quitter son village.
\item Le gouvernement cherche des solutions pour développer la campagne.
\end{enumerate}
\end{exercice}

\begin{exercice}[Traduction vers le français]
Traduis les phrases suivantes en français.
\begin{enumerate}
\item \zh{他们离开以后，农村里只剩下老人和孩子。} (Tāmen líkāi yǐhòu, nóngcūn lǐ zhǐ shèngxià lǎorén hé háizi.)
\item \zh{很多农村的年轻人去城市打工。} (Hěn duō nóngcūn de niánqīngrén qù chéngshì dǎgōng.)
\item \zh{人口外流是农村发展的大问题。} (Rénkǒu wàiliú shì nóngcūn fāzhǎn de dà wèntí.)
\end{enumerate}
\end{exercice}

\begin{exercice}[Texte à trous]
Complète le texte suivant avec les mots de la liste : \zh{离开} (líkāi), \zh{机会} (jīhuì), \zh{工资} (gōngzī), \zh{剩下} (shèngxià), \zh{办法} (bànfǎ).

\begin{textebox}[Texte à compléter]
在农村，工作很少，所以很多年轻人\_\_\_\_\_\_\_\_村子，去城市找工作。\\
Zài nóngcūn, gōngzuò hěn shǎo, suǒyǐ hěn duō niánqīngrén \_\_\_\_\_\_\_\_ cūnzi, qù chéngshì zhǎo gōngzuò.\\
城市里的\_\_\_\_\_\_\_\_比较多，\_\_\_\_\_\_\_\_也比较高。\\
Chéngshì lǐ de \_\_\_\_\_\_\_\_ bǐjiào duō, \_\_\_\_\_\_\_\_ yě bǐjiào gāo.\\
年轻人走了，村子里只\_\_\_\_\_\_\_\_老人和孩子。\\
Niánqīngrén zǒu le, cūnzi lǐ zhǐ \_\_\_\_\_\_\_\_ lǎorén hé háizi.\\
政府正在想\_\_\_\_\_\_\_\_帮助农村发展。\\
Zhèngfǔ zhèngzài xiǎng \_\_\_\_\_\_\_\_ bāngzhù nóngcūn fāzhǎn.
\end{textebox}
\end{exercice}

\begin{exercice}[Remets les mots dans l'ordre]
Remets les mots dans l'ordre pour faire des phrases correctes, puis traduis-les en français.
\begin{enumerate}
\item \zh{年轻人 / 很多 / 农村 / 离开 / 的}
\item \zh{打工 / 去 / 他们 / 城市 / 想}
\item \zh{里 / 村子 / 老人 / 只剩下 / 和 / 孩子}
\item \zh{发展 / 农村 / 帮助 / 政府 / 想}
\end{enumerate}
\end{exercice}

\courssub{Je me prépare au Bac}

\begin{exercice}[Compréhension de texte — type Bac]
Lis le texte suivant, puis réponds en chinois, par des phrases complètes, aux cinq questions de compréhension et aux deux questions de langue.

\begin{textebox}[农村的年轻人]
王明的家乡是一个小村子。那里风景很美，但是工作很少。\\
Wáng Míng de jiāxiāng shì yí ge xiǎo cūnzi. Nàlǐ fēngjǐng hěn měi, dànshì gōngzuò hěn shǎo.\\
去年，王明离开了村子，去大城市打工。\\
Qùnián, Wáng Míng líkāi le cūnzi, qù dà chéngshì dǎgōng.\\
在城市里，他找到了一份工作，工资比农村高很多。\\
Zài chéngshì lǐ, tā zhǎodào le yí fèn gōngzuò, gōngzī bǐ nóngcūn gāo hěn duō.\\
但是，他很想念家乡的父母。\\
Dànshì, tā hěn xiǎngniàn jiāxiāng de fùmǔ.\\
现在，村子里只剩下老人和孩子，人口外流越来越严重。\\
Xiànzài, cūnzi lǐ zhǐ shèngxià lǎorén hé háizi, rénkǒu wàiliú yuèláiyuè yánzhòng.\\
政府想了很多办法，希望年轻人能回家乡发展。\\
Zhèngfǔ xiǎng le hěn duō bànfǎ, xīwàng niánqīngrén néng huí jiāxiāng fāzhǎn.\\
\emph{Traduction : Le village natal de Wang Ming est un petit village. Le paysage y est très beau, mais il y a peu de travail. L'année dernière, Wang Ming a quitté le village pour aller travailler dans une grande ville. En ville, il a trouvé un emploi, et le salaire est beaucoup plus élevé qu'à la campagne. Mais ses parents du village lui manquent beaucoup. Maintenant, il ne reste dans le village que des personnes âgées et des enfants ; l'exode de la population devient de plus en plus grave. Le gouvernement a imaginé beaucoup de solutions et espère que les jeunes pourront retourner développer leur village natal.}
\end{textebox}

\textbf{Questions de compréhension} (réponds en chinois, en phrases complètes) :
\begin{enumerate}
\item \zh{王明的家乡在哪里？那里怎么样？}
\item \zh{王明为什么离开村子？}
\item \zh{他在城市里找到了什么？}
\item \zh{年轻人走了以后，村子里剩下什么人？}
\item \zh{政府希望年轻人做什么？}
\end{enumerate}

\textbf{Questions de langue} :
\begin{enumerate}
\item Trouve dans le texte une phrase avec la comparaison \zh{比} et traduis-la en français.
\item Relève dans le texte deux mots du glossaire de la leçon et emploie-les chacun dans une phrase personnelle en chinois.
\end{enumerate}
\end{exercice}

\begin{exercice}[Thème et version — type Bac]
\textbf{A. Thème.} Traduis en chinois (caractères obligatoires) :
\begin{enumerate}
\item Beaucoup de jeunes Camerounais quittent leur village pour aller travailler à Douala ou à Yaoundé.
\item Après leur départ, il ne reste à la campagne que des personnes âgées et des enfants.
\item Le gouvernement cherche des solutions pour aider les villages à se développer.
\end{enumerate}

\textbf{B. Version.} Traduis en français :
\begin{enumerate}
\item \zh{农村的机会比较少，所以人口外流越来越多。} (Nóngcūn de jīhuì bǐjiào shǎo, suǒyǐ rénkǒu wàiliú yuèláiyuè duō.)
\item \zh{在城市打工虽然工资高，但是很多年轻人想念家乡。} (Zài chéngshì dǎgōng suīrán gōngzī gāo, dànshì hěn duō niánqīngrén xiǎngniàn jiāxiāng.)
\end{enumerate}
\end{exercice}

\begin{exercice}[Expression écrite — type Bac]
\textbf{Sujet :} \zh{喀麦隆的农村人口外流} (Kāmàilóng de nóngcūn rénkǒu wàiliú) — L'exode rural au Cameroun.

Rédige en chinois un texte de \textbf{80 à 100 caractères} sur l'exode rural au Cameroun. Tu présenteras le problème, ses causes et une ou deux solutions possibles.

Consignes obligatoires :
\begin{itemize}
\item utilise au moins \textbf{trois mots du glossaire} de la leçon (par exemple \zh{离开}, \zh{机会}, \zh{工资}, \zh{剩下}, \zh{办法}, \zh{发展}) ;
\item utilise au moins \textbf{une structure de la leçon} : \zh{因为……所以……}, \zh{虽然……但是……} ou la comparaison avec \zh{比} ;
\item écris en caractères chinois (le pinyin n'est pas exigé) ;
\item soigne l'ordre des mots et la ponctuation chinoise.
\end{itemize}
\end{exercice}

\newpage

\coursec{Corrigés des exercices}

\begin{corrige}[Exercice 1 — QCM : le vocabulaire de l'exode rural]
\begin{enumerate}
\item \textbf{Bonne réponse : 2.} \zh{农村} (nóngcūn) = la campagne.
\item \textbf{Bonne réponse : 2.} \zh{外流} (wàiliú) = partir vers l'extérieur, exode (de population).
\item \textbf{Bonne réponse : 1.} \zh{打工} (dǎgōng) = travailler comme ouvrier, employé, faire un travail manuel / salarié (souvent loin de chez soi).
\item \textbf{Bonne réponse : 3.} \zh{机会} (jīhuì) = l'occasion, la chance, l'opportunité.
\end{enumerate}
\end{corrige}

\begin{corrige}[Exercice 2 — Vrai ou faux]
Les réponses se justifient par le texte \zh{农村的年轻人} (Wang Ming).
\begin{enumerate}
\item \textbf{Vrai.} Justification : \zh{去年，王明离开了村子，去大城市打工。} et \zh{人口外流越来越严重。}
\item \textbf{Faux.} Justification : Le texte dit \zh{在城市里，他找到了一份工作，工资比农村高很多。} (En ville, il a trouvé un travail, le salaire est bien plus élevé qu'à la campagne).
\item \textbf{Vrai.} Justification : \zh{现在，村子里只剩下老人和孩子。}
\item \textbf{Faux.} Justification : Le texte dit \zh{政府想了很多办法，希望年轻人能回家乡发展。} (Le gouvernement a cherché beaucoup de solutions, il espère que les jeunes pourront revenir développer leur village).
\end{enumerate}
\end{corrige}

\begin{corrige}[Exercice 3 — Trouve le mot du glossaire]
\begin{enumerate}
\item \zh{工资} (gōngzī) — salaire, paye.
\item \zh{老人} (lǎorén) — personne âgée, vieux.
\item \zh{离开} (líkāi) — quitter, partir, laisser.
\item \zh{办法} (bànfǎ) — moyen, solution,办法.
\item \zh{人口} (rénkǒu) — population, nombre d'habitants.
\end{enumerate}
\end{corrige}

\begin{corrige}[Exercice 4 — Associe le pinyin aux caractères]
\begin{center}
\begin{tabularx}{\linewidth}{|X|X|X|X|}\hline
\rowcolor{popblueL} Caractères & Pinyin & Traduction & Nature \\ \hline
1. 村 & c. cūn & village & nom \\ \hline
2. 城市 & d. chéngshì & ville & nom \\ \hline
3. 离开 & b. líkāi & quitter, partir & verbe \\ \hline
4. 发展 & a. fāzhǎn & développer, développement & verbe / nom \\ \hline
5. 剩下 & e. shèngxià & rester (il reste), rester en surplus & verbe \\ \hline
\end{tabularx}
\end{center}
\end{corrige}

\begin{corrige}[Exercice 5 — Traduction vers le chinois]
\begin{enumerate}
\item \zh{很多年轻人因为农村工作少，所以离开村子。} (Hěn duō niánqīngrén yīnwèi nóngcūn gōngzuò shǎo, suǒyǐ líkāi cūnzi.) \\
  \emph{Structure : 因为……所以……}
\item \zh{城市里机会比较多，工资也比较高。} (Chéngshì lǐ jīhuì bǐjiào duō, gōngzī yě bǐjiào gāo.) \\
  \emph{On peut aussi dire : 城市里的机会多，工资高。}
\item \zh{我爷爷不想离开村子。} (Wǒ yéye bù xiǎng líkāi cūnzi.) \\
  \emph{Grand-père paternel = 爷爷 ; grand-père maternel = 外公.}
\item \zh{政府想办法帮助农村发展。} (Zhèngfǔ xiǎng bànfǎ bāngzhù nóngcūn fāzhǎn.) \\
  \emph{想办法 = chercher des solutions, trouver un moyen.}
\end{enumerate}
\end{corrige}

\begin{corrige}[Exercice 6 — Traduction vers le français]
\begin{enumerate}
\item Après leur départ, il ne reste dans la campagne que des personnes âgées et des enfants.
\item Beaucoup de jeunes de la campagne vont en ville travailler (comme ouvriers/employés).
\item L'exode de la population est un grand problème pour le développement de la campagne.
\end{enumerate}
\end{corrige}

\begin{corrige}[Exercice 7 — Texte à trous]
\begin{textebox}[Texte complété]
在农村，工作很少，所以很多年轻人\zh{离开}村子，去城市找工作。\\
Zài nóngcūn, gōngzuò hěn shǎo, suǒyǐ hěn duō niánqīngrén \zh{líkāi} cūnzi, qù chéngshì zhǎo gōngzuò.\\
城市里的\zh{机会}比较多，\zh{工资}也比较高。\\
Chéngshì lǐ de \zh{jīhuì} bǐjiào duō, \zh{gōngzī} yě bǐjiào gāo.\\
年轻人走了，村子里只\zh{剩下}老人和孩子。\\
Niánqīngrén zǒu le, cūnzi lǐ zhǐ \zh{shèngxià} lǎorén hé háizi.\\
政府正在想\zh{办法}帮助农村发展。\\
Zhèngfǔ zhèngzài xiǎng \zh{bànfǎ} bāngzhù nóngcūn fāzhǎn.
\end{textebox}
\end{corrige}

\begin{corrige}[Exercice 8 — Remets les mots dans l'ordre]
\begin{enumerate}
\item \zh{很多农村的年轻人离开。} (Hěn duō nóngcūn de niánqīngrén líkāi.) — Beaucoup de jeunes de la campagne partent.
\item \zh{他们想去城市打工。} (Tāmen xiǎng qù chéngshì dǎgōng.) — Ils veulent aller travailler en ville.
\item \zh{村子里只剩下老人和孩子。} (Cūnzi lǐ zhǐ shèngxià lǎorén hé háizi.) — Il ne reste dans le village que des personnes âgées et des enfants.
\item \zh{政府想帮助农村发展。} (Zhèngfǔ xiǎng bāngzhù nóngcūn fāzhǎn.) — Le gouvernement veut aider la campagne à se développer.
\end{enumerate}
\end{corrige}

\begin{corrige}[Exercice 9 — Compréhension de texte — type Bac]
\textbf{Barème indicatif :} 5 questions de compréhension (2 pts chacune = 10 pts) + 2 questions de langue (5 pts chacune = 10 pts). Total 20 pts.

\textbf{Questions de compréhension (réponses en chinois, phrases complètes) :}
\begin{enumerate}
\item \zh{王明的家乡在一个小村子。那里风景很美，但是工作很少。}
\item \zh{因为村子里工作很少，所以王明离开村子去大城市打工。}
\item \zh{他在城市里找到了一份工作，工资比农村高很多。}
\item \zh{年轻人走了以后，村子里只剩下老人和孩子。}
\item \zh{政府希望年轻人能回家乡发展。}
\end{enumerate}

\textbf{Questions de langue :}
\begin{enumerate}
\item \textbf{Phrase avec \zh{比} :} \zh{工资比农村高很多。} (Gōngzī bǐ nóngcūn gāo hěn duō.) \\
  \textbf{Traduction :} Le salaire est beaucoup plus élevé qu'à la campagne.
\item \textbf{Deux mots du glossaire employés en phrases :} (exemples acceptables)
  \begin{itemize}
  \item \zh{离开} (líkāi) : \zh{我明天要离开喀麦隆去中国。} (Wǒ míngtiān yào líkāi Kāmàilóng qù Zhōngguó.) — Demain je dois quitter le Cameroun pour la Chine.
  \item \zh{机会} (jīhuì) : \zh{学习中文给我带来了很多机会。} (Xuéxí Zhōngwén gěi wǒ dàilái le hěn duō jīhuì.) — Apprendre le chinois m'a apporté beaucoup d'opportunités.
  \item \zh{发展} (fāzhǎn) : \zh{我们要帮助家乡发展。} (Wǒmen yào bāngzhù jiāxiāng fāzhǎn.) — Nous devons aider notre village natal à se développer.
  \end{itemize}
\end{enumerate}
\end{corrige}

\begin{corrige}[Exercice 10 — Thème et version — type Bac]
\textbf{A. Thème (Traduction vers le chinois)}
\begin{enumerate}
\item \zh{很多喀麦隆年轻人离开村子，去杜阿拉或雅温得打工。} (Hěn duō Kāmàilóng niánqīngrén líkāi cūnzi, qù Dù'ālā huò Yǎwēndé dǎgōng.) \\
  \emph{Dualla = 杜阿拉 ; Yaoundé = 雅温得.}
\item \zh{他们离开后，农村里只剩下老人和孩子。} (Tāmen líkāi hòu, nóngcūn lǐ zhǐ shèngxià lǎorén hé háizi.)
\item \zh{政府想办法帮助村子发展。} (Zhèngfǔ xiǎng bànfǎ bāngzhù cūnzi fāzhǎn.) \emph{ou : 政府想帮助农村发展。}
\end{enumerate}

\textbf{B. Version (Traduction vers le français)}
\begin{enumerate}
\item Les opportunités à la campagne sont relativement peu nombreuses, c'est pourquoi l'exode rural est de plus en plus important.
\item Bien que les salaires soient élevés pour travailler en ville, beaucoup de jeunes ont le mal du pays (pensent à leur village natal / regrettent leur foyer).
\end{enumerate}
\end{corrige}

\begin{corrige}[Exercice 11 — Expression écrite — type Bac]
\textbf{Barème indicatif (sur 20 pts) :}
\begin{itemize}
\item Respect du nombre de caractères (80–100) : 3 pts
\item Utilisation d'au moins 3 mots du glossaire : 3 pts
\item Utilisation d'au moins une structure requise (\zh{因为……所以……}, \zh{虽然……但是……}, \zh{比}) : 4 pts
\item Correction grammaticale et ordre des mots : 5 pts
\item Richesse du vocabulaire et cohérence du propos : 5 pts
\end{itemize}

\textbf{Proposition de rédaction modèle (environ 95 caractères) :}

\begin{textebox}[Rédaction modèle]
喀麦隆很多农村年轻人因为工作少，所以离开家乡去城市打工。\\
Kāmàilóng hěn duō nóngcūn niánqīngrén yīnwèi gōngzuò shǎo, suǒyǐ líkāi jiāxiāng qù chéngshì dǎgōng.\\
城市里机会多，工资比农村高。\\
Chéngshì lǐ jīhuì duō, gōngzī bǐ nóngcūn gāo.\\
虽然生活好一些，但他们想念父母。\\
Suīrán shēnghuó hǎo yìxiē, dànshì tāmen xiǎngniàn fùmǔ.\\
村子里只剩下老人和孩子。\\
Cūnzi lǐ zhǐ shèngxià lǎorén hé háizi.\\
政府想办法发展农村，希望年轻人回来。\\
Zhèngfǔ xiǎng bànfǎ fāzhǎn nóngcūn, xīwàng niánqīngrén huílái.
\end{textebox}

\textbf{Points de vigilance pour l'élève :}
\begin{itemize}
\item Vérifier l'ordre des mots : \zh{因为} en début de proposition, \zh{所以} en début de proposition principale.
\item \zh{比} structure : Sujet 1 + \zh{比} + Sujet 2 + Adjectif (sans \zh{很}).
\item \zh{虽然……但是……} : ne pas mettre \zh{但是} après \zh{虽然} dans la même proposition.
\item Ponctuation chinoise : 。 ， 、 ; pas d'espaces entre les caractères.
\item Mots du glossaire utilisés ici : \zh{离开}, \zh{机会}, \zh{工资}, \zh{剩下}, \zh{办法}, \zh{发展} (6 mots).
\item Structures utilisées : \zh{因为……所以……}, \zh{比}, \zh{虽然……但是……}.
\end{itemize}
\end{corrige}

\newpage

\coursec{Fiche bilan}

\begin{popfigure}
\begin{center}\tcbox[colback=popblue,colframe=popblue,coltext=white,arc=9pt,boxsep=4pt,left=6pt,right=6pt]{\bfseries\small Exode rural \zh{农村的年轻人} \emph{nóngcūn de niánqīngrén}}\end{center}
\vspace{-2pt}
\begin{tcbraster}[raster columns=3,raster equal height=rows,raster column skip=6pt,raster row skip=6pt,enhanced,blankest]
\begin{tcolorbox}[enhanced,colback=poporange,colframe=poporange,coltext=white,boxrule=0.9pt,arc=3pt,left=4pt,right=4pt,top=3pt,bottom=3pt,boxsep=1pt,halign=left]\bfseries\scriptsize Lieux \zh{农村 / 城市} village / ville\end{tcolorbox}
\begin{tcolorbox}[enhanced,colback=popgreen,colframe=popgreen,coltext=white,boxrule=0.9pt,arc=3pt,left=4pt,right=4pt,top=3pt,bottom=3pt,boxsep=1pt,halign=left]\bfseries\scriptsize Personnes \zh{年轻人 / 农民工} jeunes / ouvriers migrants\end{tcolorbox}
\begin{tcolorbox}[enhanced,colback=poppurple,colframe=poppurple,coltext=white,boxrule=0.9pt,arc=3pt,left=4pt,right=4pt,top=3pt,bottom=3pt,boxsep=1pt,halign=left]\bfseries\scriptsize Causes \zh{找工作 / 上大学} emploi, études\end{tcolorbox}
\begin{tcolorbox}[enhanced,colback=poppink,colframe=poppink,coltext=white,boxrule=0.9pt,arc=3pt,left=4pt,right=4pt,top=3pt,bottom=3pt,boxsep=1pt,halign=left]\bfseries\scriptsize Conséquences \zh{老人留在农村} vieillissement\end{tcolorbox}
\begin{tcolorbox}[enhanced,colback=popgold,colframe=popgold,coltext=white,boxrule=0.9pt,arc=3pt,left=4pt,right=4pt,top=3pt,bottom=3pt,boxsep=1pt,halign=left]\bfseries\scriptsize Solutions \zh{发展 / 机会} développement\end{tcolorbox}
\begin{tcolorbox}[enhanced,colback=popblueL!80!popblue,colframe=popblueL!80!popblue,coltext=white,boxrule=0.9pt,arc=3pt,left=4pt,right=4pt,top=3pt,bottom=3pt,boxsep=1pt,halign=left]\bfseries\scriptsize Grammaire \zh{越来越} + adj. \zh{为了} + but\end{tcolorbox}
\end{tcbraster}

\legende{Carte mentale de la leçon : l'exode rural (lexique, causes, conséquences, structures).}
\end{popfigure}

\begin{bilan}
\textbf{1. Lexique clé à connaître par cœur}
\begin{itemize}
\item \zh{农村} \emph{nóngcūn} : campagne, monde rural ; \zh{城市} \emph{chéngshì} : ville.
\item \zh{年轻人} \emph{niánqīngrén} : les jeunes ; \zh{老人} \emph{lǎorén} : les personnes âgées.
\item \zh{农民工} \emph{nóngmíngōng} : ouvrier migrant ; \zh{打工} \emph{dǎgōng} : travailler (comme ouvrier migrant).
\item \zh{离开} \emph{líkāi} : quitter ; \zh{搬到} \emph{bāndào} : déménager vers ; \zh{留在} \emph{liúzài} : rester à.
\item \zh{机会} \emph{jīhuì} : occasion, opportunité ; \zh{工作} \emph{gōngzuò} : travail ; \zh{生活} \emph{shēnghuó} : vie.
\item \zh{发展} \emph{fāzhǎn} : développement, se développer ; \zh{问题} \emph{wèntí} : problème.
\item \zh{越来越多} \emph{yuè lái yuè duō} : de plus en plus nombreux.
\end{itemize}

\textbf{2. Règles de grammaire à connaître par cœur}
\begin{center}
\begin{tabularx}{\linewidth}{|X|X|X|}
\hline
\rowcolor{popblueL} Structure & Exemple & Traduction \\
\hline
越来越 + adjectif & \zh{农村的年轻人越来越少。} \emph{Nóngcūn de niánqīngrén yuè lái yuè shǎo.} & Les jeunes à la campagne sont de moins en moins nombreux. \\
\hline
为了 + but，+ proposition & \zh{为了找工作，很多年轻人去了城市。} \emph{Wèile zhǎo gōngzuò, hěn duō niánqīngrén qù le chéngshì.} & Pour trouver du travail, beaucoup de jeunes sont partis en ville. \\
\hline
Sujet + 离开 + lieu & \zh{他们离开了农村。} \emph{Tāmen líkāi le nóngcūn.} & Ils ont quitté la campagne. \\
\hline
Sujet + 留在 / 搬到 + lieu & \zh{老人留在农村。} \emph{Lǎorén liú zài nóngcūn.} & Les personnes âgées restent à la campagne. \\
\hline
因为 + cause，所以 + conséquence & \zh{因为农村工作少，所以年轻人去城市。} \emph{Yīnwèi nóngcūn gōngzuò shǎo, suǒyǐ niánqīngrén qù chéngshì.} & Parce qu'il y a peu de travail à la campagne, les jeunes vont en ville. \\
\hline
\end{tabularx}
\end{center}

\textbf{3. Écriture des caractères (clés et ordre des traits)}
\begin{itemize}
\item \zh{农} : clé 冖 (couverture) ; 6 traits, du haut vers le bas.
\item \zh{村} : clé 木 (bois) à gauche ; d'abord 木 (gauche → droite), puis 寸.
\item \zh{城} : clé 土 (terre) à gauche ; 土 puis 成, de gauche à droite.
\item \zh{轻} : clé 车 (voiture) à gauche ; la partie gauche s'écrit avant la droite.
\item \zh{留} : clé 田 (champ) en bas ; écrire la partie haute 卯 puis 田, du haut vers le bas.
\end{itemize}

\textbf{4. Méthodes express}
\begin{itemize}
\item Comprendre un texte : lire d'abord le titre et la dernière phrase, repérer les mots connus (\zh{农村}, \zh{城市}, \zh{工作}), puis répondre en phrase complète avec les mots de la question.
\item Traduire une cause : repérer \zh{因为} (parce que) et \zh{所以} (donc) ; traduire dans le même ordre.
\item Parler de l'exode rural au Cameroun : modèle \zh{很多年轻人离开……去……找工作。} (Beaucoup de jeunes quittent... pour aller chercher du travail à...).
\end{itemize}
\end{bilan}

\begin{aretenir}[Checklist avant le Bac]
\begin{itemize}
\item[$\square$] Je sais écrire et prononcer \zh{农村}, \zh{城市}, \zh{年轻人}, \zh{农民工} avec les bons tons.
\item[$\square$] Je connais les verbes \zh{离开}, \zh{搬到}, \zh{留在} et leur construction avec un lieu.
\item[$\square$] Je sais employer la structure \zh{越来越} + adjectif sans oublier le deuxième \zh{越}.
\item[$\square$] Je sais exprimer un but avec \zh{为了} + verbe, placé en tête de phrase.
\item[$\square$] Je sais exprimer cause et conséquence avec \zh{因为……所以……}.
\item[$\square$] Je place correctement le complément de lieu après \zh{在} : \zh{留在农村}, jamais \zh{留农村在}.
\item[$\square$] Je reconnais les clés 木, 土, 田, 车 et je respecte l'ordre des traits (gauche → droite, haut → bas).
\item[$\square$] Je sais résumer le texte \zh{农村的年轻人} en trois phrases (cause, départ, conséquence).
\item[$\square$] Je sais comparer l'exode rural au Cameroun et en Chine avec des phrases simples.
\item[$\square$] Je réponds aux questions de compréhension par des phrases complètes en chinois.
\end{itemize}
\end{aretenir}

\findecours

\end{document}
